% Sampling and Estimation — Further Mathematics Upper Sixth — Cours signé OnBuch+
% Official MINESEC syllabus. Compiler avec : tectonic FMA22-sampling-and-estimation.tex
\newcommand{\DOCMATIERE}{Further Mathematics}
\newcommand{\DOCNIVEAU}{Upper Sixth}
\newcommand{\DOCMODULE}{Upper Sixth — Further mathematics}
\newcommand{\DOCLECON}{22}
\newcommand{\DOCTITRE}{Sampling and Estimation}
% OnBuch+ — préambule « cours signés » ENRICHI (compilé avec tectonic / XeLaTeX)
% Système visuel « Pop » d'OnBuch+ : fond crème (jamais blanc), bordures encre
% épaisses, ombres « dures » décalées, titres Archivo Black en capitales.
% Version MATHÉMATIQUES : graphes (pgfplots/tikz), boîtes pédagogiques
% (définition, propriété, méthode, attention, le savais-tu, exercice résolu,
% exercice, TP, bilan), page de garde et signature OnBuch+ ancrée en bas.
%
% Macros attendues dans main.tex AVANT \input{preamble} :
%   \DOCMATIERE  \DOCNIVEAU  \DOCMODULE  \DOCLECON (numéro)  \DOCTITRE
\documentclass[11pt]{article}

\usepackage{fontspec}
\usepackage[english]{babel}
\usepackage[a4paper,margin=2cm,top=3.4cm,bottom=2.8cm,headheight=1.4cm,footskip=1.3cm]{geometry}
\usepackage{amsmath,amssymb,amsfonts}
\usepackage{mathtools} % \xmapsto, \coloneqq... souvent utilisés par les modèles
\usepackage{cancel} % simplifications barrées (bilans de forces)
% \abs{z} (module, valeur absolue) et \norm{u} (norme), utilisés par les modèles
\providecommand{\abs}{}\let\abs\relax\DeclarePairedDelimiter{\abs}{\lvert}{\rvert}
\providecommand{\norm}{}\let\norm\relax\DeclarePairedDelimiter{\norm}{\lVert}{\rVert}
\usepackage[table]{xcolor} % [table] : fournit aussi \rowcolors (alternance de lignes)
\usepackage{pagecolor}
\usepackage{tikz}
\usetikzlibrary{arrows.meta,positioning,calc,shapes.geometric,shapes.misc,shapes.arrows,decorations.pathmorphing,decorations.pathreplacing,decorations.markings,patterns,shadows,mindmap,trees,fit,backgrounds,matrix,intersections,angles,quotes}
\usepackage{pgfplots}
\usepackage[version=4]{mhchem} % équations nucléaires \ce{^{235}_{92}U}
\usepackage[european,siunitx]{circuitikz} % schémas électriques
\pgfplotsset{compat=1.18}
\usepgfplotslibrary{fillbetween,groupplots} % aires entre courbes (calcul intégral)
\usepackage{enumitem}
\usepackage{fancyhdr}
% [most] charge toutes les bibliothèques tcolorbox (listings, minted, raster,
% documentation...) et gonfle inutilement la mémoire TeX sur un document long
% (~300 boîtes) : on ne charge que ce qui est réellement utilisé.
\usepackage[skins,breakable]{tcolorbox}
\usepackage{graphicx}
\usepackage{colortbl}
\usepackage{booktabs}
\usepackage{tabularx}
\usepackage{diagbox} % case d'en-tête diagonale des tableaux à double entrée
% Colonnes extensibles centrée / gauche / droite (C, L, R), souvent utilisées par les modèles
\newcolumntype{C}{>{\centering\arraybackslash}X}
\newcolumntype{L}{>{\raggedright\arraybackslash}X}
\newcolumntype{R}{>{\raggedleft\arraybackslash}X}
\usepackage{array}
\usepackage{multicol}
\usepackage{siunitx}
% parse-numbers=false : accepte aussi \SI{2,5 \times 10^{-3}}{...}
\sisetup{locale=UK,output-decimal-marker={.},group-minimum-digits=4,parse-numbers=false}
\DeclareSIUnit\ohm{\ensuremath{\symup{\Omega}}} % U+2126 absent de la police
\DeclareSIUnit\division{div} % réglages d'oscilloscope (ms/div, V/div)
\DeclareSIUnit\tr{tr} % tours (vitesse de rotation, ex. \SI{1500}{\tr\per\minute})
\usepackage{qrcode}
% \degree : commande fréquemment utilisée par les modèles pour un angle ou une
% température en dehors de \SI{}{} (non fournie par les paquets chargés).
\providecommand{\degree}{^{\circ}}
% \qed (fin de démonstration), utilisé par les modèles sans amsthm
\providecommand{\qed}{\hfill$\square$}
% \barre : double barre des valeurs interdites dans un tableau de variations (tkz-tab)
\providecommand{\barre}{\ensuremath{\|}}
% environnement proof (amsthm non chargé) utilisé par les modèles
\ifdefined\proof\else\newenvironment{proof}[1][Proof]{\par\noindent\textit{#1.}\ }{\qed\par}\fi
\usepackage{lastpage}
\usepackage{hyperref}
\hypersetup{hidelinks}

% ============================================================
% Palette « Pop » OnBuch+ — mêmes hex que la classe P (plus_ui.dart)
% ============================================================
\definecolor{poporange}{HTML}{F59321}
\definecolor{popdark}{HTML}{E07A0C}
\definecolor{popcream}{HTML}{FFF6E8}
\definecolor{popink}{HTML}{1C1714}
\definecolor{poppurple}{HTML}{7A5AE0}
\definecolor{popgreen}{HTML}{2E9E6A}
\definecolor{poppink}{HTML}{E0457B}
\definecolor{popblue}{HTML}{2F7DE1}
\definecolor{popgold}{HTML}{C98A12}
\definecolor{popmuted}{HTML}{8A7F76}
\definecolor{popline}{HTML}{EADFD2}
\definecolor{popgray}{HTML}{8A7F76}
\colorlet{popgrey}{popgray}
% Alias tolérants (le modèle nomme parfois les couleurs différemment)
\colorlet{popwhite}{white}
\colorlet{popblack}{popink}
\colorlet{popred}{poppink}
\colorlet{popyellow}{popgold}
% Teintes claires pour fonds de tableaux / zones
\colorlet{popblueL}{popblue!10!white}
\colorlet{poporangeL}{poporange!14!white}
\colorlet{popgreenL}{popgreen!12!white}
\colorlet{poppurpleL}{poppurple!12!white}
\colorlet{poppinkL}{poppink!12!white}
\colorlet{popgoldL}{popgold!14!white}

\pagecolor{popcream}

% ============================================================
% Typographie
% ============================================================
\defaultfontfeatures{Ligatures=TeX}
\setmainfont{PlusJakartaSans-Medium.ttf}[Path=../../../fonts/,BoldFont=PlusJakartaSans-Bold.ttf]
% Maths en sans-serif (Fira Math) avec chiffres et lettres droites de Plus
% Jakarta, pour que les formules se fondent dans le texte.
\usepackage[math-style=ISO,bold-style=ISO]{unicode-math}
\setmathfont{FiraMath-Regular.otf}[Path=../../../fonts/]
\setmathfont{PlusJakartaSans-Medium.ttf}[Path=../../../fonts/,range={up/num,up/{Latin,latin},\mathup}]
% (icomma retiré : décimales avec point en anglais)
% Caractères mathématiques Unicode absents des polices de texte (Plus Jakarta,
% Archivo Black) : rendus par leur équivalent mathématique, en texte comme en
% formule — sinon ils s'affichent en carré vide (ex. « Arithmétique dans ℤ »).
\usepackage{newunicodechar}
\newunicodechar{ℝ}{\ensuremath{\mathbb{R}}}
\newunicodechar{ℤ}{\ensuremath{\mathbb{Z}}}
\newunicodechar{ℂ}{\ensuremath{\mathbb{C}}}
\newunicodechar{ℕ}{\ensuremath{\mathbb{N}}}
\newunicodechar{ℚ}{\ensuremath{\mathbb{Q}}}
\newunicodechar{𝔻}{\ensuremath{\mathbb{D}}}
\newunicodechar{α}{\ensuremath{\alpha}}
\newunicodechar{β}{\ensuremath{\beta}}
\newunicodechar{γ}{\ensuremath{\gamma}}
\newunicodechar{δ}{\ensuremath{\delta}}
\newunicodechar{ε}{\ensuremath{\varepsilon}}
\newunicodechar{θ}{\ensuremath{\theta}}
\newunicodechar{λ}{\ensuremath{\lambda}}
\newunicodechar{μ}{\ensuremath{\mu}}
\newunicodechar{σ}{\ensuremath{\sigma}}
\newunicodechar{τ}{\ensuremath{\tau}}
\newunicodechar{φ}{\ensuremath{\varphi}}
\newunicodechar{ψ}{\ensuremath{\psi}}
\newunicodechar{ω}{\ensuremath{\omega}}
\newunicodechar{Σ}{\ensuremath{\Sigma}}
% majuscules produites par \MakeUppercase dans les titres (α-aminés -> Α)
\newunicodechar{Α}{\ensuremath{\alpha}}
\newunicodechar{Β}{\ensuremath{\beta}}
\newunicodechar{Γ}{\ensuremath{\Gamma}}
\newunicodechar{Δ}{\ensuremath{\Delta}}
\newunicodechar{Ε}{\ensuremath{\varepsilon}}
\newunicodechar{Θ}{\ensuremath{\Theta}}
\newunicodechar{Λ}{\ensuremath{\Lambda}}
\newunicodechar{Μ}{\ensuremath{\mu}}
\newunicodechar{Τ}{\ensuremath{\tau}}
\newunicodechar{Φ}{\ensuremath{\Phi}}
\newunicodechar{Ψ}{\ensuremath{\Psi}}
\newunicodechar{Ω}{\ensuremath{\Omega}}
\newunicodechar{↦}{\ensuremath{\mapsto}}
\newunicodechar{∈}{\ensuremath{\in}}
\newunicodechar{∉}{\ensuremath{\notin}}
\newunicodechar{∧}{\ensuremath{\wedge}}
\newunicodechar{⇔}{\ensuremath{\Leftrightarrow}}
\newunicodechar{⟺}{\ensuremath{\Longleftrightarrow}}
\newunicodechar{⇒}{\ensuremath{\Rightarrow}}
\newunicodechar{⟹}{\ensuremath{\Longrightarrow}}
\newunicodechar{∘}{\ensuremath{\circ}}
\newunicodechar{∪}{\ensuremath{\cup}}
\newunicodechar{∩}{\ensuremath{\cap}}
\newunicodechar{≡}{\ensuremath{\equiv}}
\newunicodechar{⊂}{\ensuremath{\subset}}
\newunicodechar{⊥}{\ensuremath{\perp}}
\newunicodechar{∃}{\ensuremath{\exists}}
\newunicodechar{∀}{\ensuremath{\forall}}
\newunicodechar{∛}{\ensuremath{\sqrt[3]{\,}}}
\newunicodechar{∜}{\ensuremath{\sqrt[4]{\,}}}
\newunicodechar{⃗}{\ensuremath{{}^{\to}}}
\newunicodechar{ⁿ}{\ifmmode^{n}\else\textsuperscript{\ensuremath{n}}\fi}
\newunicodechar{ˣ}{\ifmmode^{x}\else\textsuperscript{\ensuremath{x}}\fi}
\newunicodechar{ᵇ}{\ifmmode^{b}\else\textsuperscript{\ensuremath{b}}\fi}
\newunicodechar{ʳ}{\ifmmode^{r}\else\textsuperscript{\ensuremath{r}}\fi}
\newunicodechar{ᵘ}{\ifmmode^{u}\else\textsuperscript{\ensuremath{u}}\fi}
\newunicodechar{ʸ}{\ifmmode^{y}\else\textsuperscript{\ensuremath{y}}\fi}
\newunicodechar{ᵃ}{\ifmmode^{a}\else\textsuperscript{\ensuremath{a}}\fi}
\newunicodechar{ᵉ}{\ifmmode^{e}\else\textsuperscript{\ensuremath{e}}\fi}
\newunicodechar{ᵏ}{\ifmmode^{k}\else\textsuperscript{\ensuremath{k}}\fi}
\newunicodechar{ᵖ}{\ifmmode^{p}\else\textsuperscript{\ensuremath{p}}\fi}
\newunicodechar{ᴺ}{\ifmmode^{N}\else\textsuperscript{\ensuremath{N}}\fi}
\newunicodechar{ᶜ}{\ifmmode^{c}\else\textsuperscript{\ensuremath{c}}\fi}
\newunicodechar{ʰ}{\ifmmode^{h}\else\textsuperscript{\ensuremath{h}}\fi}
\newunicodechar{ᵠ}{\ifmmode^{q}\else\textsuperscript{\ensuremath{q}}\fi}
\newunicodechar{ₙ}{\ifmmode_{n}\else\textsubscript{\ensuremath{n}}\fi}
\newunicodechar{ₐ}{\ifmmode_{a}\else\textsubscript{\ensuremath{a}}\fi}
\newunicodechar{ᵢ}{\ifmmode_{i}\else\textsubscript{\ensuremath{i}}\fi}
\newunicodechar{ᵣ}{\ifmmode_{r}\else\textsubscript{\ensuremath{r}}\fi}
\newunicodechar{ₖ}{\ifmmode_{k}\else\textsubscript{\ensuremath{k}}\fi}
\newunicodechar{ᵦ}{\ifmmode_{\beta}\else\textsubscript{\ensuremath{\beta}}\fi}
\newunicodechar{ₓ}{\ifmmode_{x}\else\textsubscript{\ensuremath{x}}\fi}
\newfontfamily\popdisplay{ArchivoBlack-Regular.ttf}[Path=../../../fonts/]
\newfontfamily\poplogo{SpaceGrotesk-Bold.ttf}[Path=../../../fonts/]
\newfontfamily\popsemi{PlusJakartaSans-SemiBold.ttf}[Path=../../../fonts/]
\newfontfamily\popxbold{PlusJakartaSans-ExtraBold.ttf}[Path=../../../fonts/]
\newfontfamily\popxboldls{PlusJakartaSans-ExtraBold.ttf}[Path=../../../fonts/,LetterSpace=3.0]

\setlist[enumerate]{leftmargin=1.7em,itemsep=5pt,topsep=4pt}
\setlist[itemize]{leftmargin=1.7em,itemsep=4pt,topsep=4pt,label={\color{poporange}\textbullet}}
\setlength{\parindent}{0pt}
\setlength{\parskip}{5pt}
\renewcommand{\arraystretch}{1.25}

% pgfplots : style Pop par défaut
\pgfplotsset{
  every axis/.append style={axis line style={popink,line width=1pt},tick style={popink},
    grid style={popline,line width=0.5pt},label style={font=\small},tick label style={font=\footnotesize}},
  popcurve/.style={poporange,line width=1.8pt,no markers,smooth},
}
% Même style disponible dans un \draw TikZ simple (hors pgfplots)
\tikzset{popcurve/.style={poporange,line width=1.8pt}}
\tikzset{popgrid/.style={popline,very thin}}
\colorlet{popcurve}{poporange} % le nom du style est parfois utilisé comme couleur

% ============================================================
% Pill / squiggle
% ============================================================
\newcommand{\poppill}[3]{%
  \tikz[baseline=-0.55ex]\node[draw=popink,line width=1.2pt,rounded corners=2.6pt,%
    fill=#2,inner xsep=6.5pt,inner ysep=3pt]{%
    \popxboldls\fontsize{7}{7}\selectfont\color{#3}\MakeUppercase{#1}};%
}
\newcommand{\popsquiggle}[1]{%
  \tikz[baseline=-0.4ex]{
    \draw[#1,line width=2.6pt,line cap=round]
      plot [smooth] coordinates {(0,0.09) (0.34,0.01) (0.68,0.21) (1.02,0.01) (1.36,0.10)};
  }%
}
% Mot-clé surligné (marqueur orange sous le texte)
\newcommand{\cle}[1]{{\popxbold\color{popdark}#1}}

% ============================================================
% En-tête / pied
% ============================================================
\pagestyle{fancy}
\fancyhf{}
\renewcommand{\headrulewidth}{0pt}
\renewcommand{\footrulewidth}{0pt}
\lhead{\raisebox{-0.15ex}{\poplogo\fontsize{13}{13}\selectfont\color{popink}OnBuch\color{poporange}+}}
\rhead{\poppill{\DOCMATIERE\ \textbullet\ \DOCNIVEAU\ \textbullet\ Chapter \DOCLECON}{white}{popink}}
\lfoot{\popxbold\fontsize{7.6}{7.6}\selectfont\color{popmuted}\MakeUppercase{Signed course OnBuch+}\ \textbullet\ {\popsemi\DOCTITRE}}
\rfoot{\tikz[baseline=-0.6ex]\node[rounded corners=8.5pt,fill=popink,minimum height=17pt,minimum width=17pt,inner xsep=5pt,inner sep=0]{\popxbold\fontsize{8}{8}\selectfont\color{popcream}\thepage\,/\,\pageref*{LastPage}};}

% ============================================================
% Titres
% ============================================================
\newcommand{\chaptitre}[2]{%
  \begin{tcolorbox}[enhanced,colback=poporange,colframe=popink,boxrule=2.6pt,arc=11pt,
      drop shadow=popink,left=15pt,right=15pt,top=12pt,bottom=11pt,before skip=3pt,after skip=18pt]
    \poppill{#2}{popink}{popcream}\par\vspace{9pt}
    {\raggedright\hyphenpenalty=10000\exhyphenpenalty=10000\popdisplay\fontsize{22}{24}\selectfont\color{popcream}\MakeUppercase{#1}\par}
  \end{tcolorbox}
}
% Section numérotée : pastille orange avec le numéro + titre Archivo
\newcounter{popsec}
\newcommand{\coursec}[1]{%
  \stepcounter{popsec}\setcounter{popsub}{0}%
  \par\vspace{16pt}\noindent
  \begin{minipage}{\linewidth}
  \tikz[baseline=-0.7ex]\node[draw=popink,line width=1.6pt,fill=poporange,rounded corners=4pt,minimum size=22pt,inner sep=0]{\popdisplay\fontsize{12}{12}\selectfont\color{popcream}\thepopsec};%
  \hspace{8pt}{\popdisplay\fontsize{15}{17}\selectfont\color{popink}\MakeUppercase{#1}}\par
  \vspace{4pt}\noindent{\color{poporange}\rule{36pt}{3.6pt}}
  \end{minipage}\par\nopagebreak\vspace{8pt}
}
\newcounter{popsub}
\newcommand{\courssub}[1]{%
  \stepcounter{popsub}%
  \par\vspace{9pt}\noindent{\popxbold\fontsize{11.5}{13}\selectfont\color{popink}{\color{poporange}\thepopsec.\thepopsub}\hspace{6pt}#1}\par\nopagebreak\vspace{4pt}
}

% ============================================================
% Boîtes pédagogiques — même recette : fond blanc, bordure épaisse colorée,
% ombre dure encre, badge plein. Titre optionnel [#1].
% ============================================================
\tcbset{popbase/.style={breakable,enhanced,colback=white,boxrule=2.3pt,arc=9pt,
  shadow={3pt}{-3pt}{0pt}{popink},left=14pt,right=14pt,top=11pt,bottom=11pt,before skip=11pt,after skip=15pt,
  fontupper=\popsemi\fontsize{10.6}{14.5}\selectfont\color{popink},
  fontlower=\popsemi\fontsize{10.6}{14.5}\selectfont\color{popink}}}
% Coche dessinée (le glyphe ✓ n'existe pas dans les polices du document)
\newcommand{\popcheck}[1]{\tikz[baseline=-0.5ex]\draw[#1,line width=1.6pt,line cap=round,line join=round](-0.13,0.0)--(-0.04,-0.09)--(0.13,0.1);}
\providecommand{\coche}{\popcheck{popgreen}}
\providecommand{\resultat}[1]{\par\smallskip\noindent\fcolorbox{popgreen}{popgreenL}{\parbox{\dimexpr\linewidth-2\fboxsep-2\fboxrule\relax}{\textbf{Result.} #1}}\par\smallskip}
\newcommand{\popboxhead}[3]{\poppill{#1}{#2}{white}\ifx\relax#3\relax\else\hspace{6pt}{\popxbold\fontsize{10.6}{12}\selectfont\color{popink}#3}\fi\par\vspace{7pt}}

\newtcolorbox{aretenir}[1][]{popbase,colframe=popgreen,before upper={\popboxhead{Key points}{popgreen}{#1}}}
\newtcolorbox{exemplebox}[1][]{popbase,colframe=poporange,before upper={\popboxhead{Exemple}{poporange}{#1}}}
\newtcolorbox{definition}[1][]{popbase,colframe=popblue,before upper={\popboxhead{Definition}{popblue}{#1}}}
\newtcolorbox{propriete}[1][]{popbase,colframe=poppurple,before upper={\popboxhead{Property}{poppurple}{#1}}}
\newtcolorbox{methode}[1][]{popbase,colframe=popgold,before upper={\popboxhead{Method}{popgold}{#1}}}
\newtcolorbox{attention}[1][]{popbase,colframe=poppink,before upper={\popboxhead{Warning}{poppink}{#1}}}
\newtcolorbox{experience}[1][]{popbase,colframe=popblue,colback=popblueL,before upper={\popboxhead{Experiment}{popblue}{#1}}}
% « Le savais-tu ? » : carte violette pleine (contexte, culture, Cameroun)
\newtcolorbox{savaistu}[1][]{popbase,colback=poppurple,colframe=popink,
  fontupper=\popsemi\fontsize{10.4}{14.2}\selectfont\color{white},
  before upper={\poppill{Did you know?}{white}{poppurple}\ifx\relax#1\relax\else\hspace{6pt}{\popxbold\color{white}#1}\fi\par\vspace{7pt}}}
% Exercice résolu : énoncé en haut, \tcblower puis la solution sur fond crème
\newcounter{popexo}\newcounter{popexr}
\newtcolorbox{exoresolu}[1][]{popbase,colframe=poporange,colbacklower=popcream,
  segmentation style={popink,line width=1.2pt,dashed},
  before upper={\stepcounter{popexr}\popboxhead{Worked example \thepopexr}{poporange}{#1}},
  before lower={\poppill{Solution}{popink}{popcream}\par\vspace{6pt}}}
% Exercice (non résolu en place) — la correction est dans la partie Corrigés
\newtcolorbox{exercice}[1][]{popbase,colframe=popink,
  before upper={\stepcounter{popexo}\popboxhead{Exercise \thepopexo}{popink}{#1}}}
% Corrigé
\newtcolorbox{corrige}[1][]{popbase,colframe=popgreen,colback=popgreenL,
  before upper={\popboxhead{Solution}{popgreen}{#1}}}
% Objectifs / prérequis / bilan
\newtcolorbox{objectifs}{popbase,colframe=popink,colback=poporangeL,
  before upper={\popboxhead{Chapter objectives}{popink}{}}}
\newtcolorbox{prerequis}{popbase,colframe=popink,colback=popblueL,
  before upper={\popboxhead{Prerequisites}{popblue}{}}}
\newtcolorbox{bilan}{popbase,colframe=popink,colback=white,boxrule=2.8pt,
  before upper={\popboxhead{Fiche bilan}{poporange}{L'essentiel à connaître pour l'examen}}}
% Figure encadrée avec légende
\newtcolorbox{popfigure}[1][]{enhanced,breakable=false,colback=white,colframe=popline,boxrule=1.4pt,arc=8pt,
  left=10pt,right=10pt,top=10pt,bottom=8pt,before skip=10pt,after skip=12pt,halign=center,
  lower separated=false,halign lower=center,
  fontlower=\popsemi\fontsize{9}{11.5}\selectfont\color{popmuted},
  #1}
\newcommand{\legende}[1]{\par\vspace{6pt}{\popsemi\fontsize{9}{11.5}\selectfont\color{popmuted}#1}}

% ============================================================
% Signature OnBuch+
% ============================================================
\newcommand{\poplogoinline}[1]{{\poplogo\fontsize{#1}{#1}\selectfont\color{popink}OnBuch\color{poporange}+}}
\newcommand{\signatureonbuch}{%
  \begin{tcolorbox}[enhanced,colback=popink,colframe=popink,boxrule=2.2pt,arc=10pt,
      drop shadow=poporange,left=14pt,right=14pt,top=10pt,bottom=10pt,before skip=0pt,after skip=0pt]
    \tikz[baseline=-0.7ex]\node[circle,fill=popgreen,draw=popcream,line width=1.3pt,minimum size=22pt,inner sep=0]{\popcheck{white}};%
    \hspace{9pt}\begin{minipage}[c]{0.62\linewidth}
      {\popxbold\fontsize{12}{13}\selectfont\color{popcream}Signed\ \textbullet\ {\poplogo OnBuch\color{poporange}+}}\par\vspace{2pt}
      {\raggedright\popsemi\fontsize{8}{10.5}\selectfont\color{popline}\MakeUppercase{OnBuch+ teaching team}\\\MakeUppercase{Follows the official MINESEC syllabus}\par}
    \end{minipage}\hfill
    {\poplogo\fontsize{22}{22}\selectfont\color{popcream}OnBuch\color{poporange}+}
  \end{tcolorbox}%
}

% ============================================================
% Page de garde
% #1 titre · #2 matière · #3 niveau · #4 module · #5 n° leçon · #6 durée ·
% #7 description / objectifs (texte ou itemize)
% ============================================================
\newcommand{\pagegarde}[7]{%
  \thispagestyle{empty}
  \newgeometry{a4paper,margin=1.7cm,top=1.6cm,bottom=1.4cm}%
  \begin{tikzpicture}[remember picture,overlay]
    \fill[poporange!18] ([xshift=-3.2cm,yshift=-3.6cm]current page.north east) circle (4.6cm);
    \fill[poppurple!14] ([xshift=2.4cm,yshift=7.6cm]current page.south west) circle (3.8cm);
  \end{tikzpicture}%
  {\poplogo\fontsize{30}{30}\selectfont\color{popink}OnBuch\color{poporange}+}\hfill
  \raisebox{0.6ex}{\poppill{Signed course \textbullet\ Premium}{popink}{popcream}}\par
  \vspace{1pt}\popsquiggle{poppurple}\par
  \vspace{8pt}
  \begin{tcolorbox}[enhanced,colback=poporange,colframe=popink,boxrule=2.8pt,arc=13pt,
      drop shadow=popink,left=18pt,right=18pt,top=12pt,bottom=13pt,after skip=10pt]
    \poppill{#2 \textbullet\ #3}{popink}{popcream}\hspace{5pt}\poppill{#4}{white}{popink}\par\vspace{8pt}
    {\popdisplay\fontsize{14}{14}\selectfont\color{popink}CHAPTER #5}\par\vspace{4pt}
    {\raggedright\hyphenpenalty=10000\exhyphenpenalty=10000\popdisplay\fontsize{25}{27}\selectfont\color{popcream}\MakeUppercase{#1}\par}
    \vspace{8pt}
    {\popxbold\fontsize{9.5}{9.5}\selectfont\color{popink}\MakeUppercase{Suggested time: #6}}
  \end{tcolorbox}
  \begin{tcolorbox}[enhanced,colback=white,colframe=popink,boxrule=2.2pt,arc=10pt,
      drop shadow=popink,left=16pt,right=16pt,top=10pt,bottom=10pt,after skip=8pt]
    \poppill{In this chapter}{poppurple}{white}\par\vspace{5pt}
    \popsemi\fontsize{9.3}{12.6}\selectfont\color{popink}#7
  \end{tcolorbox}
  \noindent\begin{minipage}[t]{0.487\linewidth}
    \begin{tcolorbox}[enhanced,colback=popgreen,colframe=popink,boxrule=2.2pt,arc=10pt,
        drop shadow=popink,left=11pt,right=11pt,top=10pt,bottom=10pt,height=3.1cm,valign=center]
      \tcbox[colback=white,colframe=popink,boxrule=1.3pt,arc=5pt,boxsep=0pt,nobeforeafter,
        left=3pt,right=3pt,top=3pt,bottom=3pt,valign=center]{\qrcode[height=1.9cm]{https://play.google.com/store/apps/details?id=com.luvvix.onbuch&hl=en}}%
      \hspace{8pt}\begin{minipage}[c]{0.5\linewidth}
        {\popdisplay\fontsize{11}{12.5}\selectfont\color{white}\MakeUppercase{Download OnBuch}}\par\vspace{4pt}
        {\raggedright\popsemi\fontsize{8.4}{11}\selectfont\color{white}Free \textbullet\ Google Play\\AI tutor, past papers, results\par}
      \end{minipage}
    \end{tcolorbox}
  \end{minipage}\hfill
  \begin{minipage}[t]{0.487\linewidth}
    \begin{tcolorbox}[enhanced,colback=poppurple,colframe=popink,boxrule=2.2pt,arc=10pt,
        drop shadow=popink,left=11pt,right=11pt,top=10pt,bottom=10pt,height=3.1cm,valign=center]
      \tikz[baseline=-0.7ex]\node[circle,fill=white,draw=popink,line width=1.3pt,minimum size=1.6cm,inner sep=0]{\popdisplay\fontsize{24}{24}\selectfont\color{poppurple}+};%
      \hspace{8pt}\begin{minipage}[c]{0.6\linewidth}
        {\popdisplay\fontsize{11}{12.5}\selectfont\color{white}\MakeUppercase{Go OnBuch+}}\par\vspace{4pt}
        {\raggedright\popsemi\fontsize{8.4}{11}\selectfont\color{white}All signed courses, unlimited AI tutor, PDF sheets — OnBuch+ tab in the app\par}
      \end{minipage}
    \end{tcolorbox}
  \end{minipage}
  \vspace{8pt}
  \popcontenu
  \vfill
  \signatureonbuch
  \restoregeometry
  \newpage
}

% Rangée « Ce cours contient » (page de garde)
\newcommand{\poptile}[3]{% #1 couleur #2 gros texte #3 légende
  \begin{tcolorbox}[enhanced,colback=#1,colframe=popink,boxrule=2pt,arc=9pt,shadow={2.5pt}{-2.5pt}{0pt}{popink},
      left=6pt,right=6pt,top=9pt,bottom=9pt,halign=center,valign=center,height=1.75cm,width=\linewidth,nobeforeafter]
    {\popdisplay\fontsize{12.5}{13}\selectfont\color{white}\MakeUppercase{#2}}\par\vspace{3pt}
    {\centering\popsemi\fontsize{7.8}{9.5}\selectfont\color{white}#3\par}
  \end{tcolorbox}}
\newcommand{\popcontenu}{%
  \par\noindent{\popxboldls\fontsize{7.5}{7.5}\selectfont\color{popmuted}THIS SIGNED COURSE CONTAINS}\par\vspace{7pt}
  \noindent\begin{minipage}[t]{0.235\linewidth}\poptile{popblue}{Course}{complete, illustrated, step by step}\end{minipage}\hfill
  \begin{minipage}[t]{0.235\linewidth}\poptile{poporange}{Methods}{and worked examples}\end{minipage}\hfill
  \begin{minipage}[t]{0.235\linewidth}\poptile{poppink}{Exercises}{exam style + solutions}\end{minipage}\hfill
  \begin{minipage}[t]{0.235\linewidth}\poptile{popgreen}{Summary sheet}{the essentials to remember}\end{minipage}\par}

% Sommaire
\newtcolorbox{sommaire}{enhanced,colback=white,colframe=popink,boxrule=2.2pt,arc=10pt,
  drop shadow=popink,left=18pt,right=18pt,top=13pt,bottom=13pt,before skip=6pt,after skip=20pt,
  before upper={\poppill{Contents}{poppurple}{white}\par\vspace{9pt}\popsemi\fontsize{10.6}{14.5}\selectfont\color{popink}}}

% Signature de fin de cours — poussée tout en bas de la dernière page
\newcommand{\findecours}{%
  \par\vfill\nopagebreak
  \begin{center}\popsquiggle{poporange}\end{center}\vspace{4pt}
  \signatureonbuch
}
\AtBeginDocument{\def\llbracket{\mathopen{[\mkern-3.5mu[}}\def\rrbracket{\mathclose{]\mkern-3.5mu]}}} % intervalles d'entiers (glyphe ⟦ absent de la police)
% Glyphes absents de Fira Math : redéfinis à partir de glyphes disponibles
\AtBeginDocument{%
  \def\setminus{\mathbin{\backslash}}\let\smallsetminus\setminus
  \def\vdots{\mathord{\vbox{\baselineskip3.2pt\lineskiplimit0pt\kern4pt\hbox{.}\hbox{.}\hbox{.}}}}%
  \def\ddots{\mathinner{\mkern1mu\raise6.5pt\hbox{.}\mkern2mu\raise3.6pt\hbox{.}\mkern2mu\raise0.7pt\hbox{.}\mkern1mu}}%
  \def\triangle{\mathord{\scalebox{0.9}{$\symup{\Delta}$}}}%
  \def\nabla{\mathord{\raisebox{\depth}{\scalebox{1}[-1]{$\symup{\Delta}$}}}}%
  \def\checkmark{\popcheck{popdark}}%
  \def\star{\mathbin{\ast}}\def\bigstar{\mathord{\ast}}%
  \def\diamond{\mathbin{\scalebox{0.6}{$\Diamond$}}}%
  \def\bigcup{\mathop{\mathchoice{\raisebox{-0.3em}{\scalebox{1.7}{$\cup$}}}{\scalebox{1.25}{$\cup$}}{\cup}{\cup}}\displaylimits}%
  \def\bigcap{\mathop{\mathchoice{\raisebox{-0.3em}{\scalebox{1.7}{$\cap$}}}{\scalebox{1.25}{$\cap$}}{\cap}{\cap}}\displaylimits}%
  \def\lesseqqgtr{\mathrel{\lessgtr}}\def\bigoplus{\mathop{\oplus}\displaylimits}%
}
\providecommand{\ssi}{if and only if} % abréviation fréquente
\AtBeginDocument{\providecommand{\wideparen}{\overparen}} % arc de cercle

% Fonctions usuelles que les modèles utilisent sans que amsmath les fournisse
\providecommand{\arsinh}{\operatorname{arsinh}}\providecommand{\arcosh}{\operatorname{arcosh}}
\providecommand{\artanh}{\operatorname{artanh}}\providecommand{\arsech}{\operatorname{arsech}}
\providecommand{\arcsch}{\operatorname{arcsch}}\providecommand{\arcoth}{\operatorname{arcoth}}
\providecommand{\arcsinh}{\operatorname{arsinh}}\providecommand{\arccosh}{\operatorname{arcosh}}
\providecommand{\arctanh}{\operatorname{artanh}}\providecommand{\sech}{\operatorname{sech}}
\providecommand{\csch}{\operatorname{csch}}\providecommand{\arcsec}{\operatorname{arcsec}}
\providecommand{\arccsc}{\operatorname{arccsc}}\providecommand{\arccot}{\operatorname{arccot}}
\providecommand{\sgn}{\operatorname{sgn}}\providecommand{\lcm}{\operatorname{lcm}}
\providecommand{\Var}{\operatorname{Var}}\providecommand{\Cov}{\operatorname{Cov}}

%% FIGFIX-BEGIN v4 — correctifs globaux des figures (docs/onbuch-generation/tools/figfix)
\usepackage{adjustbox}\usepackage{etoolbox}
% toute figure TikZ/pgfplots plus large que la ligne est réduite à la largeur disponible
\BeforeBeginEnvironment{tikzpicture}{\begin{adjustbox}{max width=\linewidth}}
\AfterEndEnvironment{tikzpicture}{\end{adjustbox}}
% légendes pgfplots lisibles par-dessus les courbes
\pgfplotsset{every axis/.append style={legend style={fill=white,fill opacity=0.92,text opacity=1,draw=black!15,inner sep=3pt}}}
% étiquettes d'arêtes (syntaxe edge["..."]) avec fond blanc
\tikzset{every edge quotes/.append style={fill=white,inner sep=1.2pt}}
%% FIGFIX-END
\begin{document}

\pagegarde{Sampling and Estimation}{Further Mathematics}{Upper Sixth}{Upper Sixth — Further mathematics}{22}{13 periods}{This chapter develops the theory of sampling from finite and infinite populations, the behaviour of sample statistics, and the construction of point and interval estimates of population parameters. It equips the learner with the tools — random sampling, the Central Limit Theorem, efficient and pooled estimators, confidence intervals — needed for the GCE Advanced Level examination and for real statistical decision-making.
\vspace{4pt}
\begin{multicols}{2}\small
\begin{itemize}
\item By the end of this chapter I can distinguish sampling with replacement from sampling without replacement and count the number of possible samples in each case.
\item By the end of this chapter I can construct the sampling distribution of a statistic (mean, proportion) from a small finite population.
\item By the end of this chapter I can use random number tables to select a simple random sample and sample attributes.
\item By the end of this chapter I can state and apply the Central Limit Theorem to the distribution of the sample mean.
\item By the end of this chapter I can compare sampling methods (simple random, systematic, stratified, cluster, quota) and choose an appropriate one.
\item By the end of this chapter I can verify whether an estimator is unbiased and compute unbiased estimates of a population mean and variance.
\end{itemize}\end{multicols}}

\begin{sommaire}
\begin{multicols}{2}
\begin{enumerate}
\item Sampling from a Finite Population
\item Sampling Distributions and Random Sampling
\item The Sample Mean and the Central Limit Theorem
\item Sampling Methods
\item Point Estimation: Unbiased and Most Efficient Estimators
\item Pooled Estimation from Two Samples
\item Interval Estimation: Confidence Intervals
\item Integration activity
\item Methods and worked examples
\item Exercises
\item Solutions to the exercises
\item Summary sheet
\end{enumerate}
\end{multicols}
\end{sommaire}

\begin{objectifs}
\begin{itemize}
\item By the end of this chapter I can distinguish sampling with replacement from sampling without replacement and count the number of possible samples in each case.
\item By the end of this chapter I can construct the sampling distribution of a statistic (mean, proportion) from a small finite population.
\item By the end of this chapter I can use random number tables to select a simple random sample and sample attributes.
\item By the end of this chapter I can state and apply the Central Limit Theorem to the distribution of the sample mean.
\item By the end of this chapter I can compare sampling methods (simple random, systematic, stratified, cluster, quota) and choose an appropriate one.
\item By the end of this chapter I can verify whether an estimator is unbiased and compute unbiased estimates of a population mean and variance.
\item By the end of this chapter I can identify the most efficient estimator among unbiased estimators of the mean and variance.
\item By the end of this chapter I can pool information from two samples to estimate a common population mean and variance.
\item By the end of this chapter I can construct and interpret confidence intervals for a mean (variance known or unknown) and for the difference of two means.
\end{itemize}
\end{objectifs}

\begin{prerequis}
\begin{itemize}
\item Descriptive statistics from Lower Sixth: mean, variance and standard deviation of a data set
\item Probability: random variables, expectation E(X) and variance Var(X), including E(aX+b) and Var(aX+b)
\item The binomial and normal distributions, and standardising with Z = (X − μ)/σ
\item Use of standard normal tables (Φ(z)) and reading probabilities backwards
\item Basic combinatorics: arrangements and combinations (nCr) from Lower Sixth
\end{itemize}
\end{prerequis}

\newpage

\coursec{Sampling from a Finite Population}

Before the 2025 cocoa season, the Littoral Cooperative Union wants to estimate the average weight of a cocoa pod harvested by its $12\,000$ member farmers — but weighing every pod is impossible. In Yaoundé, the National Institute of Statistics must predict household electricity consumption from a survey of only $1\,500$ homes, and guarantee that the error stays small. How can a tiny sample speak reliably for a whole population? When does replacing a drawn ball in an urn change the answer, and when does it not? This chapter gives the mathematical machinery — sampling distributions, the Central Limit Theorem, unbiased estimators and confidence intervals — that turns a handful of observations into trustworthy statements about millions. We begin where everything begins: with the population, the sample, and the two fundamentally different ways of drawing one from the other.

\courssub{Population, Sample and Why We Sample}

Imagine the cooperative's warehouse in Douala. The \emph{population} is not the farmers — it is the collection of \emph{measurements} we are interested in: the weights of all cocoa pods harvested this season. Each pod is a \emph{sampling unit}, and the list that lets us locate pods (the register of farmers and their deliveries) is the \emph{sampling frame}.

\begin{definition}[Population, sample, sampling unit, sampling frame, parameter, statistic]
A \cle{population} is the complete collection of items (individuals, objects or measurements) about which we wish to draw conclusions. A \cle{sampling unit} is a single member of the population that can be selected. A \cle{sampling frame} is the list or device that identifies every sampling unit and from which the sample is actually drawn. A \cle{sample} is any subset of the population, of size $n$, chosen for observation. The number $N$ of units in the population is the \cle{population size}. A numerical characteristic of the population (e.g.\ mean $\mu$, variance $\sigma^2$) is called a \cle{parameter}; the corresponding numerical characteristic computed from the sample (e.g.\ sample mean $\bar{X}$, sample variance $S^2$) is called a \cle{statistic}.
\end{definition}

A \cle{census} observes \emph{every} unit of the population; a \cle{sample survey} observes only a sample. Why not always carry out a census? Three reasons dominate in practice:

\begin{itemize}
\item \textbf{Cost and time.} Weighing millions of cocoa pods, or interviewing every household in Cameroon, would cost far more than the information is worth and would finish long after the decision it was meant to inform.
\item \textbf{Destructive testing.} To test the lifetime of a batch of bulbs produced in a factory, or the breaking strength of cables for a bridge in Edea, the test \emph{destroys} the item. A census would leave nothing to sell or use — sampling is the only option.
\item \textbf{Accuracy.} Paradoxically, a well-run survey of $1\,500$ homes can be \emph{more} accurate than a sloppy census of millions, because resources can be concentrated on careful measurement of a few units.
\end{itemize}

The price we pay for sampling is \emph{uncertainty}: different samples give different answers. The whole of this chapter is about measuring and controlling that uncertainty.

\courssub{Sampling With and Without Replacement}

Take an urn containing $N = 4$ balls labelled $1, 2, 3, 4$ — our miniature "population". We draw a sample of size $n = 2$. There are two protocols:

\begin{itemize}
\item \cle{Sampling with replacement}: after noting the first ball, we put it back before the second draw. The same unit may appear twice in the sample.
\item \cle{Sampling without replacement}: the first ball is kept aside. Every unit appears at most once.
\end{itemize}

\begin{propriete}[Number of samples of size $n$ from a population of size $N$]
Let a population contain $N$ distinct units.
\begin{itemize}
\item \textbf{With replacement}, the number of \emph{ordered} samples of size $n$ is
\[ N^{n}. \]
\item \textbf{Without replacement}, the number of \emph{ordered} samples of size $n$ is
\[ N(N-1)\cdots(N-n+1) = \frac{N!}{(N-n)!}, \]
and the number of \emph{unordered} samples (where order does not matter) is
\[ \binom{N}{n} = \frac{N!}{n!\,(N-n)!}. \]
\end{itemize}
\end{propriete}

\textbf{Justification.} With replacement, each of the $n$ draws offers $N$ possibilities, independently: $N \times N \times \cdots \times N = N^n$. Without replacement, the first draw offers $N$ choices, the second $N-1$, and so on down to $N-n+1$; multiplying gives $N!/(N-n)!$ ordered samples. Since each unordered sample of $n$ distinct units can be arranged in $n!$ orders, we divide by $n!$ to obtain $\binom{N}{n}$. $\blacksquare$

For our urn ($N=4$, $n=2$): with replacement there are $4^2 = 16$ ordered samples; without replacement there are $\frac{4!}{2!} = 12$ ordered samples and $\binom{4}{2} = 6$ unordered samples.

\begin{popfigure}
\begin{center}
\begin{tikzpicture}[
  grow=right,
  level 1/.style={sibling distance=2.6cm, level distance=2.6cm},
  level 2/.style={sibling distance=0.62cm, level distance=2.6cm},
  every node/.style={font=\small},
  edge from parent/.style={draw=popmuted}]
% With replacement tree
\node[popblue] at (0,4.6) {\textbf{With replacement: 16 ordered samples}};
\node[circle,draw=popblue,fill=popblueL,inner sep=1.5pt] (r1) at (0,3.2) {start}
  child { node[circle,draw=popink,inner sep=1pt] {4}
    child { node {4} } child { node {3} } child { node {2} } child { node {1} } }
  child { node[circle,draw=popink,inner sep=1pt] {3}
    child { node {4} } child { node {3} } child { node {2} } child { node {1} } }
  child { node[circle,draw=popink,inner sep=1pt] {2}
    child { node {4} } child { node {3} } child { node {2} } child { node {1} } }
  child { node[circle,draw=popink,inner sep=1pt] {1}
    child { node {4} } child { node {3} } child { node {2} } child { node {1} } };
% Without replacement tree
\node[poporange] at (0,-1.6) {\textbf{Without replacement: 12 ordered samples}};
\node[circle,draw=poporange,fill=poporange!15,inner sep=1.5pt] (r2) at (0,-3.0) {start}
  child { node[circle,draw=popink,inner sep=1pt] {4}
    child { node {3} } child { node {2} } child { node {1} } }
  child { node[circle,draw=popink,inner sep=1pt] {3}
    child { node {4} } child { node {2} } child { node {1} } }
  child { node[circle,draw=popink,inner sep=1pt] {2}
    child { node {4} } child { node {3} } child { node {1} } }
  child { node[circle,draw=popink,inner sep=1pt] {1}
    child { node {4} } child { node {3} } child { node {2} } };
\end{tikzpicture}
\legende{Tree diagrams for two successive draws from an urn of 4 balls. With replacement, every branch has 4 twigs; without replacement, the drawn ball is removed, so every branch has only 3 twigs.}
\end{center}
\end{popfigure}

\begin{attention}[Ordered or unordered?]
In counting questions, read carefully whether the sample is \emph{ordered} or \emph{unordered}. From 4 items, the ordered samples of size 2 without replacement number $12$, but the unordered samples number only $6$ — the pairs $(1,2)$ and $(2,1)$ are the \emph{same} unordered sample. Writing $\binom{4}{2}=12$ or $\frac{4!}{2!}=6$ is a classic slip that loses marks.
\end{attention}

\begin{methode}[Enumerating all samples systematically]
To list \emph{all} samples of size 2 or 3 without omission or repetition:
\begin{enumerate}
\item Fix the first element at its smallest value; vary the last element(s) over all allowed values in increasing order.
\item Increase the first element by one and repeat, always keeping the labels in increasing order for unordered samples.
\item Count your list and check against the formula: $\binom{N}{n}$ unordered, or $\dfrac{N!}{(N-n)!}$ ordered (without replacement), or $N^n$ (with replacement).
\end{enumerate}
\end{methode}

\courssub{Independence of Draws and Variability of Sample Statistics}

The two protocols differ in a deep probabilistic way.

\textbf{With replacement}, the urn is restored before each draw, so the draws are \cle{independent} and identically distributed: $P(\text{second ball} = 3 \mid \text{first} = 3) = \tfrac{1}{4}$, exactly as if the first draw had never happened.

\textbf{Without replacement}, the draws are \emph{dependent}: $P(\text{second} = 3 \mid \text{first} = 3) = 0$, while $P(\text{second} = 3 \mid \text{first} = 1) = \tfrac{1}{3}$. Knowing the first draw changes the distribution of the second.

\begin{attention}[False independence]
When sampling \emph{without} replacement from a \emph{small} population, do not treat the draws as independent. The composition of the population changes after each draw, so probabilities must be updated. (If $N$ is huge compared with $n$, the dependence is negligible — this is why the two protocols give nearly identical answers for large populations.)
\end{attention}

Dependence has a measurable consequence: it \emph{reduces} the variability of the sample mean. Intuitively, without replacement the sample is forced to "spread out" over the population — it cannot repeat the same extreme value twice — so the sample mean is trapped closer to the population mean.

\begin{propriete}[Mean and variance of the sample mean from a finite population]
Let a population of size $N$ have mean $\mu$ and variance $\sigma^2$, and let $\bar{X}$ be the mean of a random sample of size $n$. Then, for both protocols,
\[ E(\bar{X}) = \mu. \]
With replacement:
\[ \operatorname{Var}(\bar{X}) = \frac{\sigma^2}{n}. \]
Without replacement (enrichment — proof not examinable at GCE):
\[ \operatorname{Var}(\bar{X}) = \frac{\sigma^2}{n}\cdot\frac{N-n}{N-1}. \]
The factor $\dfrac{N-n}{N-1} < 1$ is called the \cle{finite population correction} (fpc).
\end{propriete}

Note that when $n = N$ (a census without replacement), the fpc is $0$: the sample mean \emph{is} the population mean, with no variability at all — exactly as common sense demands.

\courssub{Worked Example: All Samples of Size 2 from a Population of 4}

\begin{exoresolu}[The sampling distribution of $\bar{X}$ for $N=4$, $n=2$]
A population consists of the four values $2, 4, 6, 8$ (e.g.\ the weights, in hundreds of grams, of four cocoa pods).
\begin{enumerate}
\item Compute the population mean $\mu$ and variance $\sigma^2$.
\item List all $\binom{4}{2}=6$ unordered samples of size 2 drawn without replacement, compute each sample mean, and hence find $E(\bar{X})$ and $\operatorname{Var}(\bar{X})$.
\item Repeat for sampling with replacement (all $16$ ordered samples) and compare with the formulas.
\end{enumerate}
\tcblower
\textbf{1. Population parameters.}
\[ \mu = \frac{2+4+6+8}{4} = 5, \qquad
\sigma^2 = \frac{(2-5)^2+(4-5)^2+(6-5)^2+(8-5)^2}{4} = \frac{9+1+1+9}{4} = 5. \]

\textbf{2. Without replacement.} The 6 samples, each with probability $\tfrac{1}{6}$:
\begin{center}
\begin{tabular}{|c|c|c|}
\hline
\rowcolor{popblueL} Sample & Values & Sample mean $\bar{x}$ \\
\hline
$S_1$ & $2, 4$ & $3$ \\
$S_2$ & $2, 6$ & $4$ \\
$S_3$ & $2, 8$ & $5$ \\
$S_4$ & $4, 6$ & $5$ \\
$S_5$ & $4, 8$ & $6$ \\
$S_6$ & $6, 8$ & $7$ \\
\hline
\end{tabular}
\end{center}
Hence
\[ E(\bar{X}) = \frac{3+4+5+5+6+7}{6} = \frac{30}{6} = 5 = \mu, \]
\[ \operatorname{Var}(\bar{X}) = \frac{(3-5)^2+(4-5)^2+(5-5)^2+(5-5)^2+(6-5)^2+(7-5)^2}{6} = \frac{4+1+0+0+1+4}{6} = \frac{5}{3}. \]
Check with the formula: $\dfrac{\sigma^2}{n}\cdot\dfrac{N-n}{N-1} = \dfrac{5}{2}\cdot\dfrac{2}{3} = \dfrac{5}{3}$. \checkmark

\textbf{3. With replacement.} The 16 ordered samples give means:
\begin{center}
\begin{tabular}{|c|cccccccc|}
\hline
\rowcolor{popblueL} $\bar{x}$ & $2$ & $3$ & $4$ & $5$ & $6$ & $7$ & $8$ & Total \\
\hline
Frequency & $1$ & $2$ & $3$ & $4$ & $3$ & $2$ & $1$ & $16$ \\
\hline
\end{tabular}
\end{center}
(e.g.\ $\bar{x} = 4$ arises from $(2,6), (6,2), (4,4)$). Then
\[ E(\bar{X}) = \frac{1(2)+2(3)+3(4)+4(5)+3(6)+2(7)+1(8)}{16} = \frac{80}{16} = 5 = \mu, \]
\[ \operatorname{Var}(\bar{X}) = \frac{1(9)+2(4)+3(1)+4(0)+3(1)+2(4)+1(9)}{16} = \frac{40}{16} = \frac{5}{2} = \frac{\sigma^2}{n}. \ \checkmark \]
\textbf{Comparison.} Both protocols give an unbiased sample mean ($E(\bar{X}) = \mu = 5$), but the variance is smaller without replacement: $\tfrac{5}{3} < \tfrac{5}{2}$, exactly the factor $\tfrac{N-n}{N-1} = \tfrac{2}{3}$. Sampling without replacement is \emph{more efficient}.
\end{exoresolu}

\begin{popfigure}
\begin{center}
\begin{tikzpicture}[font=\small]
\matrix[matrix of nodes, nodes={draw=popline, minimum width=1.5cm, minimum height=0.7cm, anchor=center}, row sep=-\pgflinewidth, column sep=-\pgflinewidth] (m) {
|[fill=popblueL]| Sample \& |[fill=popblueL]| Values \& |[fill=popblueL]| $\bar{x}$ \\
$S_1$ \& $2,\ 4$ \& $3$ \\
$S_2$ \& $2,\ 6$ \& $4$ \\
$S_3$ \& $2,\ 8$ \& $5$ \\
$S_4$ \& $4,\ 6$ \& $5$ \\
$S_5$ \& $4,\ 8$ \& $6$ \\
$S_6$ \& $6,\ 8$ \& $7$ \\
};
\node[below=0.15cm of m, text=popmuted] {$E(\bar{X}) = 5 = \mu$, \quad $\operatorname{Var}(\bar{X}) = \dfrac{5}{3}$};
\end{tikzpicture}
\legende{The $\binom{4}{2}=6$ unordered samples of size 2 from the population $\{2,4,6,8\}$, with their sample means. The means cluster around the population mean $\mu = 5$.}
\end{center}
\end{popfigure}

\begin{aretenir}[Key points]
\begin{itemize}
\item A \cle{census} observes the whole population; a \cle{sample} observes a subset — chosen because of cost, time, or destructive testing.
\item Number of samples of size $n$ from $N$ units: $N^n$ ordered with replacement; $\dfrac{N!}{(N-n)!}$ ordered and $\dbinom{N}{n}$ unordered without replacement.
\item Draws are independent \emph{with} replacement, dependent \emph{without} replacement.
\item In both cases $E(\bar{X}) = \mu$; with replacement $\operatorname{Var}(\bar{X}) = \dfrac{\sigma^2}{n}$; without replacement it is multiplied by the finite population correction $\dfrac{N-n}{N-1} < 1$.
\end{itemize}
\end{aretenir}

\begin{exercice}[Samples from a small population]
A population consists of the five values $1, 3, 5, 7, 9$.
\begin{enumerate}
\item How many unordered samples of size 2 can be drawn without replacement? List them all.
\item Compute the mean of each sample and verify that $E(\bar{X}) = \mu$.
\item Compute $\operatorname{Var}(\bar{X})$ directly and check the finite population correction formula.
\item (Extension) If sampling were done \emph{with} replacement, what would be $\operatorname{Var}(\bar{X})$? Compare.
\end{enumerate}
\end{exercice}

\begin{corrige}[Exercise 1]
\textbf{1.} $\binom{5}{2} = 10$ samples: $(1,3), (1,5), (1,7), (1,9), (3,5), (3,7), (3,9), (5,7), (5,9), (7,9)$.

\textbf{2.} $\mu = \dfrac{1+3+5+7+9}{5} = 5$. The sample means are $2, 3, 4, 5, 4, 5, 6, 6, 7, 8$, summing to $50$, so $E(\bar{X}) = \dfrac{50}{10} = 5 = \mu$. \checkmark

\textbf{3.} $\sigma^2 = \dfrac{16+4+0+4+16}{5} = 8$. Directly: $\operatorname{Var}(\bar{X}) = \dfrac{9+4+1+0+1+0+1+1+4+9}{10} = \dfrac{30}{10} = 3$. Formula: $\dfrac{\sigma^2}{n}\cdot\dfrac{N-n}{N-1} = \dfrac{8}{2}\cdot\dfrac{3}{4} = 3$. \checkmark

\textbf{4.} With replacement: $\operatorname{Var}(\bar{X}) = \dfrac{\sigma^2}{n} = \dfrac{8}{2} = 4$. The variance is larger ($4 > 3$) because the finite population correction $\frac{N-n}{N-1} = \frac{3}{4}$ is absent. This confirms that sampling without replacement yields a more precise estimator of the population mean.
\end{corrige}

\begin{savaistu}[Did you know?]
Modern opinion polls and market surveys in Cameroon and worldwide almost always sample \emph{without} replacement — interviewing the same household twice would waste resources. Yet when the population is very large (millions of units) and the sample is small (a few thousand), the finite population correction $\frac{N-n}{N-1}$ is practically $1$, so statisticians safely use the simpler "with replacement" formulas. The mathematics of this section explains exactly why that shortcut is legitimate.
\end{savaistu}

\coursec{Sampling Distributions and Random Sampling}

In the previous section we saw that a sample is only a small window on a population: different samples generally give different values of the sample mean, the sample variance, or the sample proportion. This simple observation raises a fundamental question: \emph{if the sample is chosen at random, how are the values of the sample statistic themselves distributed?} Answering this question is the key to everything that follows in this chapter, because it is only when we know how a statistic varies from sample to sample that we can use it to estimate a population parameter with confidence.

\courssub{Statistics as Random Variables and Sampling Distributions}

Imagine a bag containing the examination marks of every candidate in a school. You draw a sample of $n$ candidates and compute the mean mark. Repeat the experiment with a different sample and you will almost surely obtain a different mean. The sample mean is therefore not a fixed number: it is a \cle{random variable}, whose value depends on the outcome of the random selection.

\begin{definition}[Statistic and sampling distribution]
A \cle{statistic} is any quantity computed from a sample, such as the sample mean $\bar{X}$, the sample variance $S^2$, or a sample proportion $P$. Since the sample is chosen at random, a statistic is a random variable. The \cle{sampling distribution} of a statistic is the probability distribution of that statistic over \emph{all possible samples} of a given size $n$ drawn from the population.
\end{definition}

\begin{attention}[Sampling distribution $\neq$ population distribution]
A very common mistake is to confuse the \emph{population distribution} (the distribution of the individual values $X$ in the population) with the \emph{sampling distribution} of a statistic (the distribution of $\bar{X}$, $S^2$, $P$, \dots over all possible samples). They are completely different objects: the population distribution describes \emph{individuals}; the sampling distribution describes \emph{summaries of samples}. The sampling distribution is usually much less spread out than the population distribution.
\end{attention}

\courssub{The Sampling Distribution of the Sample Mean: a Complete Construction}

The best way to understand a sampling distribution is to build one by hand. Take the small finite population
\[
\{2,\ 4,\ 6,\ 8\},
\]
whose mean and variance are
\[
\mu = \frac{2+4+6+8}{4} = 5, \qquad
\sigma^2 = \frac{(2-5)^2+(4-5)^2+(6-5)^2+(8-5)^2}{4} = \frac{9+1+1+9}{4} = 5.
\]

Draw samples of size $n=2$ \textbf{with replacement}. There are $4\times 4 = 16$ equally likely ordered samples, listed below with their means:

\begin{center}
\begin{tabular}{|c|c||c|c|}
\hline
\rowcolor{popblueL} Sample & $\bar{x}$ & Sample & $\bar{x}$ \\
\hline
$(2,2)$ & $2$ & $(6,2)$ & $4$ \\
$(2,4)$ & $3$ & $(6,4)$ & $5$ \\
$(2,6)$ & $4$ & $(6,6)$ & $6$ \\
$(2,8)$ & $5$ & $(6,8)$ & $7$ \\
$(4,2)$ & $3$ & $(8,2)$ & $5$ \\
$(4,4)$ & $4$ & $(8,4)$ & $6$ \\
$(4,6)$ & $5$ & $(8,6)$ & $7$ \\
$(4,8)$ & $6$ & $(8,8)$ & $8$ \\
\hline
\end{tabular}
\end{center}

Grouping the 16 results by the value of $\bar{x}$ gives the sampling distribution of $\bar{X}$:

\begin{center}
\begin{tabular}{|c|ccccccc|}
\hline
\rowcolor{popblueL} $\bar{x}$ & $2$ & $3$ & $4$ & $5$ & $6$ & $7$ & $8$ \\
\hline
$P(\bar{X}=\bar{x})$ & $\tfrac{1}{16}$ & $\tfrac{2}{16}$ & $\tfrac{3}{16}$ & $\tfrac{4}{16}$ & $\tfrac{3}{16}$ & $\tfrac{2}{16}$ & $\tfrac{1}{16}$ \\
\hline
\end{tabular}
\end{center}

\begin{popfigure}
\begin{center}
\begin{tikzpicture}
\begin{axis}[
    ybar, bar width=14pt,
    width=11cm, height=6.5cm,
    xlabel={$\bar{x}$}, ylabel={$P(\bar{X}=\bar{x})$},
    xtick={2,3,4,5,6,7,8},
    ymin=0, ymax=0.30,
    ytick={0,0.0625,0.125,0.1875,0.25},
    yticklabels={$0$,$\frac{1}{16}$,$\frac{2}{16}$,$\frac{3}{16}$,$\frac{4}{16}$},
    axis lines=left,
    every axis plot/.append style={fill=popblue, draw=popdark},
]
\addplot coordinates {(2,0.0625) (3,0.125) (4,0.1875) (5,0.25) (6,0.1875) (7,0.125) (8,0.0625)};
\end{axis}
\end{tikzpicture}
\end{center}
\legende{Sampling distribution of $\bar{X}$ for samples of size 2 drawn with replacement from $\{2,4,6,8\}$. Note the symmetry about $\mu = 5$ and the triangular shape.}
\end{popfigure}

We can now verify two crucial facts directly from this enumerated distribution.

\textbf{Mean of the sampling distribution.}
\[
E(\bar{X}) = \frac{1}{16}\big(1\cdot 2 + 2\cdot 3 + 3\cdot 4 + 4\cdot 5 + 3\cdot 6 + 2\cdot 7 + 1\cdot 8\big) = \frac{80}{16} = 5 = \mu.
\]

\textbf{Variance of the sampling distribution.}
\[
E(\bar{X}^2) = \frac{1}{16}\big(1\cdot 4 + 2\cdot 9 + 3\cdot 16 + 4\cdot 25 + 3\cdot 36 + 2\cdot 49 + 1\cdot 64\big) = \frac{440}{16} = 27.5,
\]
\[
\operatorname{Var}(\bar{X}) = E(\bar{X}^2) - [E(\bar{X})]^2 = 27.5 - 25 = 2.5 = \frac{5}{2} = \frac{\sigma^2}{n}.
\]

Both checks succeed. This is no accident, as the following general result shows.

\begin{propriete}[Mean and variance of the sample mean]
If random samples of size $n$ are drawn \textbf{with replacement} (or from an infinite population) from a population with mean $\mu$ and variance $\sigma^2$, then
\[
E(\bar{X}) = \mu \qquad \text{and} \qquad \operatorname{Var}(\bar{X}) = \frac{\sigma^2}{n}.
\]
\end{propriete}

\textbf{Proof (examinable).} Write $\bar{X} = \dfrac{1}{n}(X_1 + X_2 + \dots + X_n)$, where each $X_i$ is the value of the $i$-th draw. Since sampling is with replacement, the $X_i$ are independent and each has $E(X_i) = \mu$ and $\operatorname{Var}(X_i) = \sigma^2$. Then, by linearity of expectation,
\[
E(\bar{X}) = \frac{1}{n}\big(E(X_1) + \dots + E(X_n)\big) = \frac{1}{n}(n\mu) = \mu.
\]
For the variance, independence allows us to add variances:
\[
\operatorname{Var}(\bar{X}) = \frac{1}{n^2}\big(\operatorname{Var}(X_1) + \dots + \operatorname{Var}(X_n)\big) = \frac{1}{n^2}(n\sigma^2) = \frac{\sigma^2}{n}. \qquad \blacksquare
\]

\begin{aretenir}[Two messages to remember]
\begin{itemize}
\item The sample mean is \emph{centred} on the population mean: on average, $\bar{X}$ hits the target $\mu$.
\item Averaging \emph{reduces variability}: $\operatorname{Var}(\bar{X}) = \sigma^2/n$, so larger samples give more reliable means. The standard deviation of $\bar{X}$, called the \cle{standard error}, is $\sigma/\sqrt{n}$.
\end{itemize}
\end{aretenir}

\courssub{Sampling Distribution of a Proportion (Sampling of Attributes)}

Often we are not interested in a measurement but in an \cle{attribute}: defective or not, vaccinated or not, in favour or against. Suppose a proportion $p$ of a large population has the attribute. We draw a random sample of size $n$ and count the number $X$ of individuals having the attribute. The \cle{sample proportion} is
\[
P = \frac{X}{n}.
\]

When sampling is with replacement (or the population is large enough that removal is negligible), each draw is an independent Bernoulli trial with probability $p$ of ``success'', so
\[
X \sim B(n, p).
\]

Consequently the sampling distribution of the proportion satisfies
\[
E(P) = E\!\left(\frac{X}{n}\right) = \frac{np}{n} = p, \qquad
\operatorname{Var}(P) = \frac{\operatorname{Var}(X)}{n^2} = \frac{np(1-p)}{n^2} = \frac{p(1-p)}{n}.
\]

\begin{aretenir}[Sampling distribution of a proportion]
For a sample of size $n$ from a population in which a proportion $p$ has a given attribute,
\[
E(P) = p, \qquad \operatorname{Var}(P) = \frac{p(1-p)}{n}, \qquad \text{standard error} = \sqrt{\frac{p(1-p)}{n}}.
\]
The sample proportion is centred on the true proportion, and its variability decreases as $n$ grows.
\end{aretenir}

\begin{exemplebox}[Proportion of defectives]
A machine produces items of which $10\%$ are defective. A random sample of $50$ items is taken. Then $X \sim B(50, 0.1)$ and the sample proportion $P = X/50$ satisfies
\[
E(P) = 0.1, \qquad \operatorname{Var}(P) = \frac{0.1 \times 0.9}{50} = 0.0018,
\]
so the standard error is $\sqrt{0.0018} \approx 0.0424$. The probability that the sample contains exactly $4$ defectives is
\[
P(X = 4) = \binom{50}{4}(0.1)^4(0.9)^{46} \approx 0.181.
\]
\end{exemplebox}

\courssub{Random Sampling and Random Number Tables}

All the theory above rests on one essential hypothesis: the sample must be a \cle{random sample}.

\begin{definition}[Random sample]
A sample of size $n$ from a population of size $N$ is a \cle{(simple) random sample} if every member of the population has the same probability of being selected, and (for sampling without replacement) every possible sample of size $n$ has the same probability of being chosen.
\end{definition}

How do we achieve this in practice? Drawing names from a hat is unreliable (papers fold differently, mixing is imperfect). The standard tool is a table of \cle{random numbers}: a list of digits in which each digit $0$--$9$ appears with equal probability, independently of the others.

\begin{methode}[Selecting a sample with a random number table]
\begin{enumerate}
\item \textbf{Number the population.} Assign to each of the $N$ members a distinct number with the same number of digits (e.g.\ $01$ to $40$ for a class of $40$).
\item \textbf{Choose a starting point} in the table (close your eyes and point, or use a stated row and column) and a reading direction (across rows, say).
\item \textbf{Read digits in blocks} of the chosen length (two digits at a time for a population numbered $01$--$40$).
\item \textbf{Reject} any block that is out of range ($00$, or $41$--$99$ here) and, when sampling \emph{without replacement}, any block that has \emph{already been selected}.
\item \textbf{Stop} when $n$ valid distinct numbers have been obtained, and match them to the population list.
\end{enumerate}
\end{methode}

\begin{attention}[Repeated numbers]
When sampling \textbf{without replacement}, a number that has already been drawn must be rejected and reading continues. A frequent exam error is to keep a repeated number (giving a sample that is too small or counting one individual twice) or, conversely, to reject repeats when sampling \emph{with} replacement. Always read the question carefully.
\end{attention}

\begin{experience}[Selecting 8 students from a class of 40]
A class list contains $40$ students numbered $01$ to $40$. We use the extract of a random number table below, starting at the beginning of the first row and reading pairs of digits from left to right.

\begin{popfigure}
\begin{center}
\begin{tikzpicture}[scale=0.62, every node/.style={font=\small}]
\foreach \row/\nums in {
  0/{5,4,2,7,1,3,3,9,0,8,2,2,1,7,4,0,6,1,3,5},
  1/{1,2,0,5,3,8,7,4,1,1,2,9,3,6,0,4,1,8,2,3},
  2/{9,7,4,4,2,2,0,1,6,6,3,3,1,0,5,5,8,8,4,9}}{
  \foreach [count=\c from 0] \d in \nums {
    \node[minimum size=0.62cm] at (\c*0.62, -\row*0.62) {\d};
  }
}
% Highlight the eight accepted pairs in row 0: 27, 13, 39, 08, 22, 17, 40, 35
\draw[poporange, very thick, rounded corners=2pt] (2*0.62-0.31, 0.31) rectangle (4*0.62-0.31, -0.31);
\draw[poporange, very thick, rounded corners=2pt] (4*0.62-0.31, 0.31) rectangle (6*0.62-0.31, -0.31);
\draw[poporange, very thick, rounded corners=2pt] (6*0.62-0.31, 0.31) rectangle (8*0.62-0.31, -0.31);
\draw[poporange, very thick, rounded corners=2pt] (8*0.62-0.31, 0.31) rectangle (10*0.62-0.31, -0.31);
\draw[poporange, very thick, rounded corners=2pt] (10*0.62-0.31, 0.31) rectangle (12*0.62-0.31, -0.31);
\draw[poporange, very thick, rounded corners=2pt] (12*0.62-0.31, 0.31) rectangle (14*0.62-0.31, -0.31);
\draw[poporange, very thick, rounded corners=2pt] (14*0.62-0.31, 0.31) rectangle (16*0.62-0.31, -0.31);
\draw[poporange, very thick, rounded corners=2pt] (18*0.62-0.31, 0.31) rectangle (20*0.62-0.31, -0.31);
\node[poporange, font=\footnotesize] at (10*0.62, 0.75) {accepted pairs};
\end{tikzpicture}
\end{center}
\legende{Extract of a random number table. Pairs are read left to right; valid pairs in the range $01$--$40$ are highlighted.}
\end{popfigure}

Reading the first row in pairs and applying the rules:
\begin{itemize}
\item $54 > 40$: reject (out of range);
\item $27, 13, 39, 08, 22, 17, 40$: all in range and distinct — accept;
\item $61 > 40$: reject;
\item $35$: in range and not yet selected — accept.
\end{itemize}
We now have eight distinct valid numbers, so we stop. The sample consists of the students whose list numbers are
\[
08,\ 13,\ 17,\ 22,\ 27,\ 35,\ 39,\ 40.
\]
Every student had the same chance of selection, so this is a genuine random sample.
\end{experience}

\begin{exoresolu}[Sampling with a random number table]
A school canteen manager wants to select a random sample of $5$ days from the $60$ school days of a term, numbered $01$ to $60$. Starting at the given point, the random digits read in pairs are:
\[
73,\ 12,\ 60,\ 94,\ 12,\ 07,\ 55,\ 41,\ 88,\ 03, \dots
\]
Which days form the sample if sampling is without replacement?
\tcblower
\textbf{Solution.} Read each pair and apply the rules:
\begin{itemize}
\item $73 > 60$: reject (out of range).
\item $12$: accept — day $12$.
\item $60$: accept — day $60$.
\item $94 > 60$: reject.
\item $12$: already selected — reject (sampling without replacement).
\item $07$: accept — day $07$.
\item $55$: accept — day $55$.
\item $41$: accept — day $41$.
\end{itemize}
We now have five distinct valid numbers, so we stop. The sample is
\[
\{07,\ 12,\ 41,\ 55,\ 60\}.
\]
Note that $88$ and $03$ were never needed since the sample was complete.
\end{exoresolu}

\courssub{Advantages and Limitations of Random Number Tables}

\begin{center}
\begin{tabularx}{\linewidth}{|X|X|}
\hline
\rowcolor{popblueL} \textbf{Advantages} & \textbf{Limitations} \\
\hline
Guarantees equal probability of selection: removes human bias. & Requires the whole population to be listed and numbered first (a \emph{sampling frame}). \\
\hline
Simple, cheap, and requires no technology. & Tedious for large populations or large samples. \\
\hline
The procedure is transparent and reproducible if the starting point is recorded. & Careless reading (skipping, mis-grouping digits) can destroy randomness. \\
\hline
\end{tabularx}
\end{center}

\begin{savaistu}[Random numbers in your pocket]
Modern scientific calculators and spreadsheets have a random number function (often labelled \texttt{Rand} or \texttt{ran\#}) that produces (pseudo-)random numbers, and these are accepted in examinations when tables are not provided. The principle is identical: generate a number, scale it to the range $1$ to $N$, reject out-of-range or repeated values. Computers generate \emph{pseudo}-random numbers by deterministic algorithms, but for sampling purposes they behave like true random numbers.
\end{savaistu}

\begin{exercice}
The population $\{1, 3, 5\}$ is sampled with replacement using samples of size $2$.
\begin{enumerate}
\item List all $9$ ordered samples and their means.
\item Construct the sampling distribution of $\bar{X}$.
\item Verify that $E(\bar{X}) = \mu$ and $\operatorname{Var}(\bar{X}) = \sigma^2/2$.
\end{enumerate}
\end{exercice}

\begin{corrige}[Exercise 1]
\begin{enumerate}
\item The $9$ ordered samples and their means are:
$(1,1)\!:1$, $(1,3)\!:2$, $(1,5)\!:3$, $(3,1)\!:2$, $(3,3)\!:3$, $(3,5)\!:4$, $(5,1)\!:3$, $(5,3)\!:4$, $(5,5)\!:5$.
\item Grouping by the value of $\bar{x}$:
\begin{center}
\begin{tabular}{|c|ccccc|}
\hline
\rowcolor{popblueL} $\bar{x}$ & $1$ & $2$ & $3$ & $4$ & $5$ \\
\hline
$P(\bar{X}=\bar{x})$ & $\tfrac{1}{9}$ & $\tfrac{2}{9}$ & $\tfrac{3}{9}$ & $\tfrac{2}{9}$ & $\tfrac{1}{9}$ \\
\hline
\end{tabular}
\end{center}
\item The population parameters are
\[
\mu = \frac{1+3+5}{3} = 3, \qquad \sigma^2 = \frac{(1-3)^2+(3-3)^2+(5-3)^2}{3} = \frac{8}{3}.
\]
From the sampling distribution:
\[
E(\bar{X}) = \tfrac{1}{9}(1) + \tfrac{2}{9}(2) + \tfrac{3}{9}(3) + \tfrac{2}{9}(4) + \tfrac{1}{9}(5) = \frac{27}{9} = 3 = \mu. \ \checkmark
\]
\[
E(\bar{X}^2) = \tfrac{1}{9}(1) + \tfrac{2}{9}(4) + \tfrac{3}{9}(9) + \tfrac{2}{9}(16) + \tfrac{1}{9}(25) = \frac{93}{9} = \frac{31}{3},
\]
\[
\operatorname{Var}(\bar{X}) = \frac{31}{3} - 9 = \frac{4}{3} = \frac{8/3}{2} = \frac{\sigma^2}{n}. \ \checkmark
\]
\end{enumerate}
\end{corrige}

\begin{exercice}
A factory's production contains $5\%$ defective items. A random sample of $40$ items is inspected and $X$ denotes the number of defectives.
\begin{enumerate}
\item State the distribution of $X$.
\item Write down the mean and variance of the sample proportion $P = X/40$.
\item Find the probability that the sample contains exactly one defective.
\end{enumerate}
\end{exercice}

\begin{corrige}[Exercise 2]
\begin{enumerate}
\item Each item is independently defective with probability $p = 0.05$, so $X \sim B(40, 0.05)$.
\item $E(P) = p = 0.05$ and
\[
\operatorname{Var}(P) = \frac{p(1-p)}{n} = \frac{0.05 \times 0.95}{40} = 0.0011875.
\]
\item $P(X = 1) = \dbinom{40}{1}(0.05)^1(0.95)^{39} = 40 \times 0.05 \times (0.95)^{39} \approx 0.271$.
\end{enumerate}
\end{corrige}

\begin{exercice}
Using the random digits $47,\ 83,\ 06,\ 92,\ 33,\ 06,\ 71,\ 19, \dots$, select a random sample of $4$ candidates from a list of $50$ candidates numbered $01$ to $50$, sampling without replacement. Justify each acceptance or rejection.
\end{exercice}

\begin{corrige}[Exercise 3]
Apply the rules pair by pair:
\begin{itemize}
\item $47$: in range — accept.
\item $83 > 50$: reject (out of range).
\item $06$: in range — accept.
\item $92 > 50$: reject.
\item $33$: in range — accept.
\item $06$: already selected — reject (sampling without replacement).
\item $71 > 50$: reject.
\item $19$: in range — accept.
\end{itemize}
Four distinct valid numbers have been obtained, so we stop. The sample is $\{06, 19, 33, 47\}$.
\end{corrige}

\coursec{The Sample Mean and the Central Limit Theorem}

In the previous section we defined the \cle{sampling distribution} of a statistic: the probability distribution of a quantity such as the sample mean $\bar{X}$, computed over all possible samples of a given size. We now focus on the most important statistic of all, the \cle{sample mean}, and we answer two questions. First, what is the distribution of $\bar{X}$ when the population is normal? Second --- and this is one of the most remarkable results in all of mathematics --- what happens when the population is \emph{not} normal? The answer to the second question is the \cle{Central Limit Theorem}, the engine that drives almost everything we shall do later with confidence intervals and hypothesis testing.

\courssub{Sampling from an Infinite or Very Large Population}

So far we have sampled from \emph{finite} populations (a class of 40 students, a lot of 500 bags of cement). Many real populations are effectively \cle{infinite}: all possible measurements of a physical quantity, all possible phone calls made on a network, all possible rolls of a die. Even when the population is finite but very large (all adults in Cameroon), sampling a few dozen individuals barely changes the composition of the population, so the draws behave \emph{as if} the population were infinite.

\begin{definition}[Random sample from an infinite population]
A \cle{random sample} of size $n$ from a population with mean $\mu$ and variance $\sigma^{2}$ is a collection of random variables $X_{1}, X_{2}, \dots, X_{n}$ which are \textbf{independent} and \textbf{identically distributed} (i.i.d.), each with
\[ E(X_{i}) = \mu \qquad \text{and} \qquad \operatorname{Var}(X_{i}) = \sigma^{2}. \]
The \cle{sample mean} is the random variable
\[ \bar{X} = \frac{X_{1} + X_{2} + \cdots + X_{n}}{n}. \]
\end{definition}

Two facts follow immediately from the linearity of expectation and the variance of a sum of \emph{independent} variables:
\begin{align*}
E(\bar{X}) &= \frac{1}{n}\bigl(E(X_{1}) + \cdots + E(X_{n})\bigr) = \frac{n\mu}{n} = \mu,\\[2pt]
\operatorname{Var}(\bar{X}) &= \frac{1}{n^{2}}\bigl(\operatorname{Var}(X_{1}) + \cdots + \operatorname{Var}(X_{n})\bigr) = \frac{n\sigma^{2}}{n^{2}} = \frac{\sigma^{2}}{n}.
\end{align*}

\begin{propriete}[Mean and variance of the sample mean]
For a random sample of size $n$ from \emph{any} population with mean $\mu$ and variance $\sigma^{2}$,
\[ E(\bar{X}) = \mu \qquad \text{and} \qquad \operatorname{Var}(\bar{X}) = \frac{\sigma^{2}}{n}. \]
This result is examinable and you should be able to derive it from the rules $E(aX+bY)=aE(X)+bE(Y)$ and $\operatorname{Var}(aX)=a^{2}\operatorname{Var}(X)$, using independence for the variance of the sum.
\end{propriete}

The intuition is simple: averages are \emph{less spread out} than individual values. One customer may spend 5 minutes or 50 minutes on a call, but the \emph{average} of 30 calls will rarely be extreme, because large and small values tend to cancel. The factor $1/n$ in the variance quantifies exactly this smoothing.

\begin{definition}[Standard error of the mean]
The \cle{standard error of the mean} is the standard deviation of the sampling distribution of $\bar{X}$:
\[ \operatorname{SE}(\bar{X}) = \frac{\sigma}{\sqrt{n}}. \]
It measures the typical distance between a sample mean and the true population mean $\mu$.
\end{definition}

Note that the standard error decreases like $1/\sqrt{n}$: to \emph{halve} the standard error you must \emph{quadruple} the sample size. This square-root law explains why very large samples are needed for very precise surveys.

\courssub{The Sample Mean from a Normal Population}

When the population itself is normal, we can say much more: the sampling distribution of $\bar{X}$ is \emph{exactly} normal, for \emph{every} sample size, even $n = 2$.

\begin{propriete}[Sample mean of a normal population]
If $X \sim N(\mu, \sigma^{2})$ and $X_{1}, \dots, X_{n}$ is a random sample, then
\[ \bar{X} \sim N\!\left(\mu, \frac{\sigma^{2}}{n}\right) \qquad \text{exactly, for all } n. \]
Consequently the standardised variable
\[ Z = \frac{\bar{X} - \mu}{\sigma/\sqrt{n}} \sim N(0,1). \]
\end{propriete}

\textbf{Why is this true?} A fundamental property of the normal distribution (admitted at this level) is that any linear combination of independent normal random variables is itself normal. Since $\bar{X} = \frac{1}{n}X_{1} + \cdots + \frac{1}{n}X_{n}$ is such a linear combination, $\bar{X}$ is normal; and we have already proved that its mean is $\mu$ and its variance is $\sigma^{2}/n$. The distribution is therefore completely determined.

\begin{attention}[The classic trap: $\sigma$ versus $\sigma/\sqrt{n}$]
When standardising a \textbf{single observation} use $Z = \dfrac{X - \mu}{\sigma}$. When standardising a \textbf{sample mean} you must use $Z = \dfrac{\bar{X} - \mu}{\sigma/\sqrt{n}}$. Forgetting the $\sqrt{n}$ is the single most common error in this chapter at GCE. Always ask yourself first: \emph{am I dealing with one value or with an average of $n$ values?}
\end{attention}

\begin{exoresolu}[Bags of cement from a normal population]
A factory in Douala fills bags of cement whose masses are normally distributed with mean $\mu = 50.0\ \mathrm{kg}$ and standard deviation $\sigma = 0.8\ \mathrm{kg}$. A random sample of $16$ bags is taken. Find the probability that the sample mean mass is less than $49.7\ \mathrm{kg}$.
\tcblower
\textbf{Solution.} The population is normal, so $\bar{X} \sim N\!\left(50.0, \dfrac{0.8^{2}}{16}\right)$ exactly. The standard error is
\[ \frac{\sigma}{\sqrt{n}} = \frac{0.8}{\sqrt{16}} = 0.2\ \mathrm{kg}. \]
Standardising:
\[ P(\bar{X} < 49.7) = P\!\left(Z < \frac{49.7 - 50.0}{0.2}\right) = P(Z < -1.5) = 1 - \Phi(1.5) = 1 - 0.9332 = 0.0668. \]
So there is about a $6.7\%$ chance that the mean mass of the 16 bags falls below $49.7\ \mathrm{kg}$. Compare: for a \emph{single} bag, $P(X < 49.7) = P(Z < -0.375) \approx 0.354$ --- the average is far less likely to be low than an individual bag.
\end{exoresolu}

\courssub{The Central Limit Theorem}

What if the population is \emph{not} normal --- suppose it is uniform, or strongly skewed like the durations of phone calls? We can no longer claim that $\bar{X}$ is exactly normal. But something extraordinary happens as $n$ grows.

\begin{propriete}[The Central Limit Theorem (CLT)]
Let $X_{1}, X_{2}, \dots, X_{n}$ be a random sample (independent, identically distributed) from \textbf{any} population with finite mean $\mu$ and finite variance $\sigma^{2}$. Then as $n \to \infty$, the distribution of the standardised sample mean tends to the standard normal distribution:
\[ Z = \frac{\bar{X} - \mu}{\sigma/\sqrt{n}} \xrightarrow{\;d\;} N(0,1). \]
In practice, for large $n$ we use the approximation
\[ \bar{X} \approx N\!\left(\mu, \frac{\sigma^{2}}{n}\right). \]
\textbf{Rule of thumb:} the approximation is good enough for most examination purposes when $n \geq 30$; if the population is already symmetric, even $n \approx 10$ may suffice, and if the population is normal the result is exact for all $n$.
\end{propriete}

The full proof requires moment generating functions and is beyond the GCE syllabus, but the idea can be made plausible. The standardised sum $\dfrac{\sum X_{i} - n\mu}{\sigma\sqrt{n}}$ is built from many small independent contributions; no single term dominates, and the combined effect of many small independent fluctuations is known (from the theory of generating functions) to produce the bell-shaped curve. What matters for the examination is the \emph{statement}, the \emph{conditions} (random sample, finite variance, large $n$), and confident \emph{use} of the approximation.

\begin{attention}[The CLT applies to means and totals, not to single observations]
The CLT says nothing about individual values. If call durations are skewed, then a \emph{single} call duration remains skewed no matter how large your sample is. The theorem concerns $\bar{X}$ (or equivalently the total $T = \sum X_{i}$, since $T = n\bar{X}$). Writing ``by the CLT, $X \sim N(\mu,\sigma^{2})$'' for a single observation is a serious error.
\end{attention}

\begin{popfigure}
\begin{tikzpicture}
\begin{axis}[
  width=14cm, height=9.5cm,
  xmin=0, xmax=10, ymin=0, ymax=0.62,
  axis lines=left,
  xlabel={$x$}, ylabel={density},
  xtick={2,5,8}, ytick=\empty,
  legend style={font=\small, at={(0.98,0.97)}, anchor=north east, draw=popline},
  samples=200, domain=0:10, every axis plot/.append style={thick}
]
% Uniform population density: 1/6 = 0.1667 on [2,8]
\addplot[popblue] coordinates {(2,0) (2,0.1667) (8,0.1667) (8,0)};
\addlegendentry{Uniform population on $[2{,}8]$}
% n=2: exact triangular density. Var(X)=3, Var(Xbar)=1.5. Base [2,8], peak at 5, height 1/3.
\addplot[poporange] { (x<5)*((x-2)/9) + (x>=5)*((8-x)/9) };
\addlegendentry{$\overline X${,} $n=2$ (exact triangular)}
% n=5: approx normal, mu=5, sigma/sqrt(5)=sqrt(3/5)=0.7746
\addplot[popgreen] { 1/(0.7746*sqrt(max(0,2*pi)))*exp(-(x-5)^2/(2*0.7746^2)) };
\addlegendentry{$\overline X${,} $n=5$ (approx normal)}
% n=30: approx normal, mu=5, sigma/sqrt(30)=sqrt(0.1)=0.3162
\addplot[popdark] { 1/(0.3162*sqrt(max(0,2*pi)))*exp(-(x-5)^2/(2*0.3162^2)) };
\addlegendentry{$\overline X${,} $n=30$ (nearly normal)}
\end{axis}
\end{tikzpicture}
\legende{Sampling from a uniform population on $[2,8]$ ($\mu=5$, $\sigma^2=3$): the distribution of $\bar{X}$ is triangular for $n=2$ and rapidly becomes bell-shaped; by $n=30$ it is almost perfectly normal, centred on $\mu=5$ with shrinking spread $\sigma/\sqrt{n}$.}
\end{popfigure}

The same phenomenon occurs for skewed populations (such as exponential call durations): the skewness of $\bar{X}$ is divided by $\sqrt{n}$, so the sampling distribution symmetrises as $n$ grows. This is why the normal distribution appears everywhere in nature --- heights, measurement errors, examination marks are all \emph{sums of many small effects}.

\begin{popfigure}
\begin{tikzpicture}
\begin{axis}[
  width=13cm, height=8cm,
  xmin=-4, xmax=4, ymin=0, ymax=1.35,
  axis lines=left,
  xlabel={$\bar{x} - \mu$ (in units of $\sigma$)}, ylabel={density},
  xtick={-3,-2,-1,0,1,2,3}, ytick=\empty,
  legend style={font=\small, at={(0.98,0.97)}, anchor=north east, draw=popline},
  samples=150, domain=-4:4, every axis plot/.append style={thick}
]
\addplot[popmuted] {0.4*exp(-x^2/2)};
\addlegendentry{$n=1$: spread $\sigma$}
\addplot[popblue] {0.8*exp(-x^2/0.5)};
\addlegendentry{$n=4$: spread $\sigma/2$}
\addplot[poporange] {1.28*exp(-x^2/0.2)};
\addlegendentry{$n=25$: spread $\sigma/5$}
\end{axis}
\end{tikzpicture}
\legende{The standard error $\sigma/\sqrt{n}$ shrinks like $1/\sqrt{n}$: the sampling distribution of $\bar{X}$ becomes taller and narrower as $n$ increases, concentrating around $\mu$.}
\end{popfigure}

\courssub{A Four-Step Method for Probability Questions about \texorpdfstring{$\bar{X}$}{the sample mean}}

\begin{methode}[Probabilities involving the sample mean]
\begin{enumerate}
\item \textbf{Identify the model.} Write down $\mu$, $\sigma$ and $n$. Decide whether the population is normal (exact result) or non-normal with $n \geq 30$ (CLT approximation). State $\bar{X} \sim N(\mu, \sigma^{2}/n)$, exactly or approximately.
\item \textbf{Compute the standard error} $\sigma/\sqrt{n}$.
\item \textbf{Standardise} the boundary value(s): $z = \dfrac{\bar{x} - \mu}{\sigma/\sqrt{n}}$.
\item \textbf{Read the normal table}: use $\Phi(z)$, with $\Phi(-z) = 1 - \Phi(z)$ and $P(a < Z < b) = \Phi(b) - \Phi(a)$.
\end{enumerate}
For a \textbf{total} $T = X_{1} + \cdots + X_{n}$, use $T \approx N(n\mu, n\sigma^{2})$, i.e. standardise with $\dfrac{t - n\mu}{\sigma\sqrt{n}}$.
\end{methode}

\begin{exoresolu}[Lengths of phone calls --- CLT with a skewed population]
The duration of a phone call on a certain network has mean $4.2$ minutes and standard deviation $3.5$ minutes; the distribution is strongly right-skewed. A researcher records a random sample of $50$ calls.
\begin{enumerate}
\item[(a)] Find the approximate probability that the mean duration of the 50 calls exceeds $5$ minutes.
\item[(b)] Find the approximate probability that the total duration of the 50 calls is less than $180$ minutes.
\end{enumerate}
\tcblower
\textbf{Solution.} Although the population is skewed, $n = 50 \geq 30$, so by the Central Limit Theorem
\[ \bar{X} \approx N\!\left(4.2,\ \frac{3.5^{2}}{50}\right), \qquad \frac{\sigma}{\sqrt{n}} = \frac{3.5}{\sqrt{50}} \approx 0.4950. \]
(a) Standardising the boundary $\bar{x} = 5$:
\[ z = \frac{5 - 4.2}{0.4950} \approx 1.62, \qquad P(\bar{X} > 5) \approx 1 - \Phi(1.62) = 1 - 0.9474 = 0.0526. \]
(b) For the total $T = \sum X_{i}$: $E(T) = 50 \times 4.2 = 210$ and $\operatorname{sd}(T) = \sigma\sqrt{n} = 3.5\sqrt{50} \approx 24.75$. Then
\[ P(T < 180) \approx P\!\left(Z < \frac{180 - 210}{24.75}\right) = P(Z < -1.21) = 1 - \Phi(1.21) = 1 - 0.8869 = 0.1131. \]
So there is about a $5.3\%$ chance the mean exceeds $5$ minutes and about an $11.3\%$ chance the total is below $180$ minutes. Note carefully: we could \emph{not} have answered these questions for a single call, because the shape of the individual distribution is unknown.
\end{exoresolu}

\courssub{Bridge: Sample Proportions and the Normal Approximation to the Binomial}

You met the CLT in disguise in Lower Sixth. If $X \sim B(n, p)$ counts the number of successes in $n$ trials, then $X = Y_{1} + \cdots + Y_{n}$ where each $Y_{i} \sim B(1, p)$ has mean $p$ and variance $p(1-p)$. The \cle{sample proportion} $\hat{P} = X/n$ is therefore a \emph{sample mean} of Bernoulli variables, and the CLT gives, for large $n$,
\[ \hat{P} \approx N\!\left(p,\ \frac{p(1-p)}{n}\right), \qquad\text{equivalently}\qquad X \approx N\bigl(np,\ np(1-p)\bigr), \]
which is exactly the normal approximation to the binomial (valid when $np > 5$ and $n(1-p) > 5$, with a continuity correction since $X$ is discrete). The CLT is the single principle unifying all these approximations.

\begin{savaistu}[Why ``normal''?]
The bell curve was studied by de Moivre (1733) as an approximation to the binomial, then by Gauss in the theory of measurement errors. The name ``Central Limit Theorem'' was popularised by P\'olya in 1920: ``central'' because of its central \emph{importance} in probability theory. It is the reason quality control in factories, opinion polling before elections, and medical trials all rely on normal-based calculations even when individual measurements are not normal.
\end{savaistu}

\begin{exercice}[GCE practice]
The mass of a loaf of bread produced by a bakery has mean $500\ \mathrm{g}$ and standard deviation $12\ \mathrm{g}$, with unknown distribution. An inspector takes a random sample of $36$ loaves.
\begin{enumerate}
\item State the approximate distribution of the sample mean mass, justifying your answer.
\item Find the probability that the sample mean lies between $497\ \mathrm{g}$ and $503\ \mathrm{g}$.
\item Find the probability that the total mass of the 36 loaves exceeds $18.15\ \mathrm{kg}$.
\end{enumerate}
\end{exercice}

\begin{corrige}[GCE practice]
\textbf{1.} The population distribution is unknown but $n = 36 \geq 30$, so by the CLT
\[ \bar{X} \approx N\!\left(500,\ \frac{12^{2}}{36}\right) = N(500, 4), \qquad \frac{\sigma}{\sqrt{n}} = \frac{12}{6} = 2\ \mathrm{g}. \]
\textbf{2.} Standardising both bounds: $z_{1} = \dfrac{497-500}{2} = -1.5$ and $z_{2} = \dfrac{503-500}{2} = 1.5$. Hence
\[ P(497 < \bar{X} < 503) = \Phi(1.5) - \Phi(-1.5) = 2\Phi(1.5) - 1 = 2(0.9332) - 1 = 0.8664. \]
\textbf{3.} For the total: $E(T) = 36 \times 500 = 18000\ \mathrm{g}$, $\operatorname{sd}(T) = 12\sqrt{36} = 72\ \mathrm{g}$. With $18.15\ \mathrm{kg} = 18150\ \mathrm{g}$:
\[ P(T > 18150) = P\!\left(Z > \frac{18150 - 18000}{72}\right) = P(Z > 2.08) \approx 1 - 0.9812 = 0.0188. \]
\end{corrige}

\begin{aretenir}[Key points of this section]
\begin{itemize}
\item A random sample from an infinite population is modelled by i.i.d. variables $X_{1}, \dots, X_{n}$ with $E(X_{i}) = \mu$, $\operatorname{Var}(X_{i}) = \sigma^{2}$.
\item Always: $E(\bar{X}) = \mu$ and $\operatorname{Var}(\bar{X}) = \sigma^{2}/n$; the \cle{standard error} is $\sigma/\sqrt{n}$.
\item If the population is normal, $\bar{X} \sim N(\mu, \sigma^{2}/n)$ \textbf{exactly} for all $n$.
\item \textbf{Central Limit Theorem:} for any population with finite variance, $\bar{X} \approx N(\mu, \sigma^{2}/n)$ when $n$ is large (rule of thumb $n \geq 30$).
\item Standardise with $Z = \dfrac{\bar{X}-\mu}{\sigma/\sqrt{n}}$; for totals use $T \approx N(n\mu, n\sigma^{2})$.
\item The CLT applies to means and totals, never to single observations; the sample proportion $\hat{P}$ is a special case (normal approximation to the binomial).
\end{itemize}
\end{aretenir}

\coursec{Sampling Methods}

In the previous section we established a remarkable fact: whatever the shape of the population, the sample mean $\bar{X}$ of a random sample of size $n$ behaves, for large $n$, like a normal variable with mean $\mu$ and variance $\sigma^{2}/n$. But the Central Limit Theorem has a silent hypothesis hidden inside the words ``\cle{random sample}'': every formula we have proved --- $E(\bar{X})=\mu$, $\mathrm{Var}(\bar{X})=\sigma^{2}/n$, the confidence intervals to come --- collapses if the sample was not properly drawn. A sample of $50$ students obtained by interviewing the $50$ students who happen to be in the library at noon is \emph{not} a random sample of the school, and no theorem will rescue it.

This section is therefore devoted to a very practical question: \emph{how do we actually select the $n$ units?} We survey the classical \cle{sampling methods}, explain how each is carried out, and analyse its strengths, weaknesses and traps. This is also favourite material for GCE essay questions, where you may be asked to \emph{choose and justify} a method for a concrete situation.

\courssub{Probability and Non-Probability Sampling: an Overview}

All sampling methods fall into two great families.

\begin{definition}[Probability and non-probability sampling]
A \cle{probability sampling} method is one in which every unit of the population has a \textbf{known, non-zero probability} of being selected, determined by a deliberate random mechanism (lottery, random number tables, random number generator). A \cle{non-probability sampling} method is one in which the selection depends on convenience, judgement or the goodwill of respondents, so that selection probabilities are unknown or zero for some units.
\end{definition}

Only probability sampling allows us to \emph{quantify} the error: it is precisely because selection probabilities are known that we can compute $\mathrm{Var}(\bar{X})$ and build confidence intervals. Non-probability methods (quota, convenience, volunteer samples) are cheaper and faster, and are widely used in market research, but their results cannot be justified by the theory of this chapter. The main probability methods we study are \textbf{simple random}, \textbf{systematic}, \textbf{stratified} and \textbf{cluster} sampling; the main non-probability method you must know is \textbf{quota sampling}.

\courssub{Simple Random Sampling (SRS)}

Lesson 149 of the syllabus distinguishes between sampling \emph{with} and \emph{without} replacement. Both are simple random sampling provided the selection probabilities are equal at every draw.

\begin{definition}[Simple random sampling]
A \cle{simple random sample} (SRS) of size $n$ from a population of size $N$ is a sample drawn in such a way that \textbf{every possible sample of $n$ units has the same probability} of being chosen.
\begin{itemize}
\item \textbf{SRS without replacement (SRSWOR):} The $n$ units are distinct. There are $\binom{N}{n}$ possible samples, each with probability $1/\binom{N}{n}$. Every unit has inclusion probability $n/N$.
\item \textbf{SRS with replacement (SRSWR):} Units are drawn one by one and returned to the population before the next draw. There are $N^n$ possible ordered samples, each with probability $1/N^n$. Units may appear more than once.
\end{itemize}
Unless stated otherwise, ``SRS'' in this course means \textbf{without replacement}, which is the standard in survey practice because it gives more information per unit sampled.
\end{definition}

There are two classical implementations for SRSWOR.

\begin{enumerate}
\item \textbf{The lottery method.} Write the name (or number) of each of the $N$ units on identical slips of paper, mix thoroughly, and draw $n$ slips without replacement. This is exactly how a fair tombola works.
\item \textbf{Random number tables / generators.} Give each unit a serial number from $1$ to $N$ (the list of all units is called the \cle{sampling frame}). Read off random digits from a random number table (or a calculator's \texttt{Rand} function), group them into numbers of the appropriate length, keep those between $1$ and $N$, ignore repetitions, and stop when $n$ distinct numbers have been obtained.
\end{enumerate}

\begin{exemplebox}[Drawing an SRS with random numbers]
A school has $N = 800$ students and we want an SRS of $n = 20$ students. Number the students $001$ to $800$ using the class lists (the sampling frame). Reading three-digit groups from a random number table we obtain, say,
\[ 743,\ 092,\ 815,\ 377,\ 641,\ 092,\ 018,\ \dots \]
We keep $743$, $092$, $377$, $641$, $018$; we reject $815$ (greater than $800$) and the second $092$ (already chosen). We continue until $20$ distinct numbers are obtained, and interview the corresponding students.
\end{exemplebox}

\textbf{Strengths and weaknesses.} The SRS is the gold standard of fairness: it is free of selection bias and all the theory of sampling distributions applies directly. Its weaknesses are practical: it requires a \emph{complete and up-to-date sampling frame}, and the selected units may be scattered over a wide area, making data collection expensive (imagine interviewing $100$ randomly chosen households spread over the whole North-West Region).

\courssub{Systematic Sampling}

When the population is already arranged in a list, there is a quicker way to spread the sample evenly.

\begin{definition}[Systematic sampling]
To draw a \cle{systematic sample} of size $n$ from an ordered population of size $N$: compute the \cle{sampling interval} $k = N/n$ (rounded to the nearest integer if necessary), choose a \textbf{random start} $r$ uniformly in $\{1, 2, \dots, k\}$, and select the units numbered
\[ r,\ r+k,\ r+2k,\ \dots,\ r+(n-1)k. \]
If $N$ is not an exact multiple of $n$, the actual sample size may be $n$ or $n\pm 1$ depending on the rounding of $k$ and the value of $r$.
\end{definition}

\begin{methode}[Carrying out systematic sampling]
\begin{enumerate}
\item Obtain an ordered list (frame) of the $N$ units.
\item Compute the sampling interval $k = N/n$; if it is not an integer, round it (the sample size may then differ slightly from $n$).
\item Draw a random number $r$ between $1$ and $k$ inclusive (this random start is essential: it is what makes the method a probability method).
\item Select every $k$-th unit starting from $r$: $r, r+k, r+2k, \dots$
\end{enumerate}
\end{methode}

\begin{exemplebox}[Systematic sampling on a production line]
A bakery in Bamenda produces $N = 2000$ loaves per day and quality control must inspect $n = 50$ loaves. The sampling interval is $k = 2000/50 = 40$. A random number between $1$ and $40$ is drawn, say $r = 17$. The inspector then examines loaves number $17, 57, 97, 137, \dots, 1977$ as they leave the oven. Only \emph{one} random draw was needed instead of fifty.
\end{exemplebox}

Systematic sampling is fast, cheap, and spreads the sample evenly over the list --- which often makes it \emph{more} representative than a pure SRS when the list is arranged in a meaningful order (e.g., by geographical location). But it hides a serious trap.

\begin{attention}[The periodicity trap]
If the list has a \textbf{hidden periodic pattern} whose period coincides with (or is a multiple of) the sampling interval $k$, the systematic sample can be \emph{heavily biased}. Classic example: houses are numbered along a street, odd numbers on one side, even on the other; sampling every $k = 10$ houses starting from an odd number selects \emph{only} odd-numbered houses --- one side of the street only. Similarly, if a machine produces a defective item every $40$ items and we sample every $40$ items, we either catch all the defective items or none of them, depending on the random start. Always check the frame for periodic structure before using systematic sampling.
\end{attention}

\courssub{Stratified Sampling}

Suppose we want to estimate the average height of students in a school. Boys and girls differ systematically in height; a sample that happens to contain $90\%$ of boys would overestimate the mean. The remedy is to \emph{force} the sample to respect the known structure of the population.

\begin{definition}[Stratified random sampling]
In \cle{stratified sampling}, the population is first divided into $L$ non-overlapping, exhaustive subgroups called \cle{strata} which are \textbf{homogeneous within themselves} but different from one another (e.g.\ boys/girls, urban/rural, Form classes). A simple random sample is then drawn \textbf{independently within every stratum}. With \cle{proportional allocation}, the number sampled from stratum $h$ is
\[ n_h = n \times \frac{N_h}{N}, \]
where $N_h$ is the size of stratum $h$, $N = \sum_{h=1}^L N_h$ and $n$ the total sample size. (Other allocations, e.g. Neyman allocation, exist but proportional allocation is the standard for the GCE syllabus.)
\end{definition}

\begin{methode}[Proportional allocation in stratified sampling]
\begin{enumerate}
\item Identify a variable known to influence the quantity studied, and divide the population into strata $1, 2, \dots, L$ of sizes $N_1, \dots, N_L$.
\item For each stratum compute $n_h = n \cdot N_h / N$ and round to an integer (adjusting so that $\sum_{h=1}^L n_h = n$).
\item Draw an independent SRS of size $n_h$ within each stratum $h$.
\item Combine the results: the overall mean is estimated by $\bar{x}_{\mathrm{st}} = \dfrac{1}{N}\displaystyle\sum_{h=1}^L N_h\,\bar{x}_h$.
\end{enumerate}
\end{methode}

\begin{exoresolu}[Stratified sampling of a school]
A college has $N = 1200$ students: $N_1 = 480$ in the first cycle (Forms 1--5) and $N_2 = 720$ in the second cycle (Lower and Upper Sixth). We want a stratified sample of $n = 100$ students with proportional allocation. Compute the stratum sample sizes.
\tcblower
\textbf{Solution.} Proportional allocation gives
\[ n_1 = 100 \times \frac{480}{1200} = 40, \qquad n_2 = 100 \times \frac{720}{1200} = 60. \]
We check $40 + 60 = 100$. We then draw an SRS of $40$ students from the first-cycle list and an independent SRS of $60$ students from the senior list. Every stratum is represented in the sample in exactly its population proportion, so no group can be over- or under-represented by bad luck.
\end{exoresolu}

\textbf{Why stratification gains precision.} The variance of the stratified estimator depends only on the variability \emph{within} strata. If the strata are well chosen (homogeneous inside), this within-stratum variability is small, so the stratified estimate is \emph{more precise} than an SRS of the same total size. Stratification also guarantees that small but important subgroups appear in the sample.

\courssub{Cluster Sampling}

Stratified sampling needs a frame of \emph{all} units and sends the investigator into every stratum. Sometimes the opposite is wanted: for cost reasons we prefer to concentrate the fieldwork in a few places.

\begin{definition}[Cluster sampling]
In \cle{cluster sampling}, the population is divided into groups called \cle{clusters} (usually natural or geographical: villages, schools, markets, classes). A simple random sample of \textbf{whole clusters} is drawn, and \textbf{every unit} in the selected clusters is observed (one-stage cluster sampling) or a further sample is taken inside them (two-stage sampling). Ideally clusters are \textbf{heterogeneous within themselves}, each being a miniature of the whole population.
\end{definition}

\begin{exemplebox}[Cluster sampling of households]
To study household expenditure in a division, an SRS of households would scatter the interviewers over hundreds of kilometres. Instead, the division is divided into $60$ villages (clusters); $8$ villages are chosen at random, and \emph{all} households in those $8$ villages are interviewed. Travel costs fall dramatically.
\end{exemplebox}

The price to pay is precision: units inside the same cluster tend to resemble one another (people in the same village have similar incomes and habits), so a cluster sample of $n$ persons carries less information than an SRS of $n$ persons. Cluster sampling trades \emph{statistical precision for lower cost}.

\begin{attention}[Stratified or cluster? Do not confuse them]
The two methods both divide the population into groups, but the logic is \textbf{opposite}:
\begin{itemize}
\item \textbf{Stratified:} groups homogeneous \emph{inside}; we sample \textbf{within every stratum} (some units from each group). Aim: precision.
\item \textbf{Cluster:} groups heterogeneous \emph{inside}; we sample \textbf{some whole clusters} (all units from a few groups). Aim: lower cost.
\end{itemize}
In an exam, ask yourself: ``Are all groups represented in the sample?'' If yes $\to$ stratified. If only some groups are fully taken $\to$ cluster.
\end{attention}

\begin{popfigure}
\begin{tikzpicture}[scale=0.62, every node/.style={transform shape}]
% Population grid 8x5 with two types: circles (type A) and squares (type B)
% Panel 1: SRS
\begin{scope}[shift={(0,0)}]
\foreach \x in {0,...,7} {\foreach \y in {0,...,4} {
  \pgfmathparse{Mod(\x+\y,2)==0?1:0}
  \ifnum\pgfmathresult=1 \node[circle,draw=popmuted,inner sep=1.6pt] at (\x*0.5,\y*0.5) {};
  \else \node[rectangle,draw=popmuted,inner sep=1.8pt] at (\x*0.5,\y*0.5) {}; \fi}}
\foreach \p in {(0.5,2),(1.5,0.5),(2.5,1.5),(3,0),(3.5,2),(1,1),(2,0.5),(0,1.5)} \node[star,star points=5,fill=poporange,inner sep=1.1pt] at \p {};
\node[align=center,font=\small\bfseries,text=popink] at (1.75,-0.7) {Simple random};
\end{scope}
% Panel 2: Systematic
\begin{scope}[shift={(5.2,0)}]
\foreach \x in {0,...,7} {\foreach \y in {0,...,4} {
  \pgfmathparse{Mod(\x+\y,2)==0?1:0}
  \ifnum\pgfmathresult=1 \node[circle,draw=popmuted,inner sep=1.6pt] at (\x*0.5,\y*0.5) {};
  \else \node[rectangle,draw=popmuted,inner sep=1.8pt] at (\x*0.5,\y*0.5) {}; \fi}}
\foreach \p in {(0.5,2),(1.5,2),(2.5,2),(3.5,2),(0.5,0.5),(1.5,0.5),(2.5,0.5),(3.5,0.5)} \node[star,star points=5,fill=popblue,inner sep=1.1pt] at \p {};
\node[align=center,font=\small\bfseries,text=popink] at (1.75,-0.7) {Systematic};
\end{scope}
% Panel 3: Stratified
\begin{scope}[shift={(10.4,0)}]
\foreach \x in {0,...,7} {\foreach \y in {0,...,4} {
  \pgfmathparse{Mod(\x+\y,2)==0?1:0}
  \ifnum\pgfmathresult=1 \node[circle,draw=popmuted,inner sep=1.6pt] at (\x*0.5,\y*0.5) {};
  \else \node[rectangle,draw=popmuted,inner sep=1.8pt] at (\x*0.5,\y*0.5) {}; \fi}}
\foreach \p in {(0,2),(1.5,1),(3,0.5),(2.5,2)} \node[star,star points=5,fill=popgreen,inner sep=1.1pt] at \p {};
\foreach \p in {(0.5,0.5),(2,1.5),(3.5,1),(1,2)} \node[diamond,fill=popgreen,inner sep=1.1pt] at \p {};
\node[align=center,font=\small\bfseries,text=popink] at (1.75,-0.7) {Stratified};
\node[align=center,font=\scriptsize,text=popmuted] at (1.75,-1.15) {both types sampled};
\end{scope}
% Panel 4: Cluster
\begin{scope}[shift={(15.6,0)}]
\draw[rounded corners,draw=popline] (-0.2,-0.2) rectangle (1.7,2.2);
\draw[rounded corners,draw=popline] (1.8,-0.2) rectangle (3.7,2.2);
\foreach \x in {0,...,3} {\foreach \y in {0,...,4} {
  \pgfmathparse{Mod(\x+\y,2)==0?1:0}
  \ifnum\pgfmathresult=1 \node[circle,draw=popmuted,inner sep=1.6pt] at (\x*0.5,\y*0.5) {};
  \else \node[rectangle,draw=popmuted,inner sep=1.8pt] at (\x*0.5,\y*0.5) {}; \fi}}
\foreach \x in {4,...,7} {\foreach \y in {0,...,4} {
  \pgfmathparse{Mod(\x+\y,2)==0?1:0}
  \ifnum\pgfmathresult=1 \node[circle,draw=popmuted,inner sep=1.6pt] at (\x*0.5,\y*0.5) {};
  \else \node[rectangle,draw=popmuted,inner sep=1.8pt] at (\x*0.5,\y*0.5) {}; \fi}}
\foreach \x in {0,...,3} {\foreach \y in {0,...,4} \node[star,star points=5,fill=poppurple,inner sep=1.0pt] at (\x*0.5,\y*0.5) {};}
\node[align=center,font=\small\bfseries,text=popink] at (1.75,-0.7) {Cluster};
\node[align=center,font=\scriptsize,text=popmuted] at (1.75,-1.15) {one whole cluster taken};
\end{scope}
\end{tikzpicture}
\legende{The same population (two types of units) sampled four ways. SRS: units scattered at random. Systematic: evenly spaced units. Stratified: samples drawn from \emph{both} types. Cluster: \emph{all} units of one randomly chosen cluster.}
\end{popfigure}

\courssub{Quota Sampling (a Non-Probability Method)}

\begin{definition}[Quota sampling]
In \cle{quota sampling}, the interviewer is instructed to interview fixed numbers (\emph{quotas}) of people in given categories --- e.g.\ $30$ men and $30$ women, of whom $20$ are under $25$ and $40$ are $25$ or over --- but is free to choose \emph{which} individuals fill each quota (typically in the street, at a market entrance, \dots).
\end{definition}

Quota sampling \emph{mimics} stratified sampling in that the sample respects population proportions, but there is \textbf{no random selection at any stage}: the interviewer chooses the respondents, and unconsciously picks approachable, available, cooperative people. The selection probabilities are unknown, so sampling errors \emph{cannot be computed} and the theory of this chapter does not apply. It is fast and cheap, which is why market-research and opinion-poll companies use it, but it is exposed to hidden bias.

\begin{savaistu}[A famous failure of non-random sampling]
In 1936, the American magazine \emph{Literary Digest} polled more than two million people --- drawn from telephone directories and car registration lists --- and confidently predicted a landslide defeat for Franklin D. Roosevelt. Roosevelt won comfortably. The frame (telephone and car owners, during the Great Depression) over-represented the wealthy, and the huge sample size could not compensate for the biased selection. The lesson is permanent: \textbf{a large biased sample is worse than a small random one}. This is exactly the kind of example that earns marks in a GCE essay on sampling.
\end{savaistu}

\courssub{Summary Comparison of Probability Sampling Methods}

\begin{center}
\begin{tabularx}{\linewidth}{|l|X|X|X|X|}\hline
\rowcolor{popblueL}
\textbf{Method} & \textbf{Frame needed} & \textbf{Selection} & \textbf{Precision} & \textbf{Main use} \\ \hline
Simple Random (SRS) & Complete list of units & Equal probability, independent & High (benchmark) & Small, accessible populations \\ \hline
Systematic & Ordered list & Fixed interval $k$, random start & High if list order unrelated to variable; low if periodic & Production lines, ordered registers \\ \hline
Stratified & List + stratum membership & SRS \emph{within each stratum} & Highest (if strata homogeneous) & Heterogeneous populations with known subgroups \\ \hline
Cluster & List of clusters only & SRS of clusters, census inside & Lower (intra-cluster correlation) & Large geographical areas, no unit frame \\ \hline
\end{tabularx}
\end{center}

\courssub{Choosing a Method: Worked Scenarios}

\begin{exoresolu}[Matching a sampling method to a situation]
For each situation, propose a suitable sampling method and justify the choice.
\begin{enumerate}
\item A school of $900$ students: estimate the proportion of students who eat at the canteen, using $90$ students.
\item A region: estimate the mean household income from $400$ households, with a small budget for transport.
\item A market: estimate the mean mass of $50\ \mathrm{kg}$ bags of garri sold by $120$ traders, testing $60$ bags.
\end{enumerate}
\tcblower
\textbf{Solution.}
\begin{enumerate}
\item The class lists provide a complete frame, so a \textbf{simple random sample} of $90$ students (random numbers from $001$ to $900$) is ideal. If canteen use is known to differ between first and second cycle, a \textbf{stratified} sample with proportional allocation ($n_h = 90\,N_h/900$) would be even more precise.
\item No list of all households exists and transport is expensive: use \textbf{cluster sampling}. Divide the region into villages or quarters, choose a few clusters at random, and survey households within the chosen clusters. Precision is sacrificed for feasibility and cost.
\item Number the traders $1$ to $120$; the sampling interval is $k = 120/60 = 2$. Choose a random start $r \in \{1,2\}$ and test one bag from every second trader: this is \textbf{systematic sampling}, quick and spread over the whole market. (Check first that trader numbering has no periodic pattern, e.g.\ alternating wholesalers and retailers.)
\end{enumerate}
\end{exoresolu}

\courssub{Sources of Bias}

Whatever the method, three classic defects can destroy representativeness.

\begin{itemize}
\item \textbf{Non-response bias.} Selected individuals who refuse or cannot be reached often differ systematically from respondents (busy people, hostile people). Replacing them by ``the neighbour who was at home'' reintroduces selection bias. Good practice: call-backs and follow-up visits.
\item \textbf{Volunteer (self-selected) samples.} Radio phone-ins and online polls only reach people \emph{motivated enough} to respond; they over-represent strong opinions. A volunteer sample is never a random sample.
\item \textbf{Faulty sampling frames.} If the list is incomplete or outdated (new quarters missing, deceased persons still listed, telephone directories in a country where many have no landline), some units have probability zero of selection and the bias cannot be removed by increasing $n$.
\end{itemize}

\begin{attention}[What randomness does --- and does not --- protect against]
Random selection protects against \emph{selection bias} and allows us to compute sampling errors. It does \emph{not} protect against non-response, lying respondents, defective measuring instruments or a faulty frame. A perfect random mechanism applied to a bad frame still gives a bad sample.
\end{attention}

\begin{aretenir}[Key points]
\begin{itemize}
\item Probability methods (SRS, systematic, stratified, cluster) give known selection probabilities and allow error computation; non-probability methods (quota, convenience, volunteer) do not.
\item SRS can be with or without replacement; without replacement is standard (more efficient).
\item Systematic sampling: interval $k = N/n$, random start $r \in \{1,\dots,k\}$; beware of periodic lists.
\item Stratified: homogeneous strata, sample \emph{within every stratum}, proportional allocation $n_h = nN_h/N$; gains precision.
\item Cluster: heterogeneous clusters, take \emph{some whole clusters}; saves cost, loses precision.
\item A large biased sample is worse than a small random one: no sample size cures a faulty frame or self-selection.
\end{itemize}
\end{aretenir}

\begin{exercice}[Choosing and defending a method]
A cocoa cooperative has $N = 2400$ member farmers spread over $40$ villages, $60$ farmers per village on average. The cooperative wants to estimate the mean yield per farmer from a sample of $120$ farmers.
\begin{enumerate}
\item Describe how to obtain a systematic sample of $120$ farmers from the membership register.
\item Describe a cluster sampling plan using villages, and give one advantage and one disadvantage compared with (a).
\item Explain why interviewing the $120$ farmers who attend the annual general meeting would be unsatisfactory.
\end{enumerate}
\end{exercice}

\begin{corrige}[Exercise 1]
\begin{enumerate}
\item $k = 2400/120 = 20$. Draw a random start $r$ between $1$ and $20$, then select farmers numbered $r, r+20, r+40, \dots$ from the register. One random draw suffices and the sample is spread over the whole register.
\item Choose, say, $2$ villages at random out of $40$ and interview all their farmers (about $120$). Advantage: fieldwork concentrated in two places, so much cheaper. Disadvantage: farmers in the same village resemble one another (same soils, same techniques), so the estimate is less precise than a register-wide sample of the same size.
\item Attendance at the meeting is voluntary: attendees are the most active, most prosperous or most concerned farmers. This is a self-selected (volunteer) sample with unknown selection probabilities; it will very likely overestimate the mean yield, and no sampling theory can correct it.
\end{enumerate}
\end{corrige}

\coursec{Point Estimation: Unbiased and Most Efficient Estimators}

In the previous section we studied \emph{how} to select a sample from a population: simple random sampling, systematic sampling, stratified sampling and so on. But a sample is never an end in itself. The cocoa exporter in Douala who weighs 40 bags out of a consignment of 5000 does not care about those 40 bags for their own sake: he wants to \emph{draw a conclusion about the whole consignment}. This step --- using the numbers computed from a sample to make a statement about an unknown feature of the population --- is called \cle{estimation}, and it is the heart of this section. We shall ask two natural questions: what makes an estimator \emph{good}, and among several good estimators, which one is \emph{best}?

\courssub{Parameters, Statistics, Estimators and Estimates}

\begin{definition}[Parameter and statistic]
A \cle{parameter} is a numerical characteristic of a \textbf{population}, usually unknown: the population mean $\mu$, the population variance $\sigma^{2}$, a population proportion $p$. A \cle{statistic} is any quantity computed from a \textbf{sample}: the sample mean $\bar{X}$, the sample variance, the sample proportion. Statistics are used to estimate parameters.
\end{definition}

\begin{definition}[Estimator and estimate]
Let $\theta$ be an unknown population parameter. An \cle{estimator} of $\theta$, written $\hat{\theta}$, is a \textbf{random variable} defined as a function of the sample observations $X_{1}, X_{2}, \dots, X_{n}$. An \cle{estimate} is the numerical \textbf{value} taken by $\hat{\theta}$ for one particular observed sample.
\end{definition}

For example, $\hat{\mu} = \bar{X} = \dfrac{1}{n}\sum_{i=1}^{n} X_{i}$ is an \emph{estimator} (a random variable, with its own sampling distribution); if the 40 cocoa bags give $\bar{x} = 62.4$ kg, then $62.4$ is an \emph{estimate} (a single number).

\begin{attention}[Estimator or estimate?]
The estimator is a \emph{rule} (a random variable, capital letters); the estimate is the \emph{result of applying the rule} to one sample (a number, lower case). Confusing the two in an examination costs marks: $\bar{X}$ varies from sample to sample, $\bar{x}$ does not.
\end{attention}

\courssub{Bias of an Estimator}

Since $\hat{\theta}$ is random, it will sometimes overshoot $\theta$, sometimes undershoot it. What we demand, at minimum, is that it should be right \emph{on average}.

\begin{definition}[Bias and unbiased estimator]
The \cle{bias} of an estimator $\hat{\theta}$ of a parameter $\theta$ is
\[
\operatorname{Bias}(\hat{\theta}) = E(\hat{\theta}) - \theta.
\]
The estimator $\hat{\theta}$ is \cle{unbiased} if $\operatorname{Bias}(\hat{\theta}) = 0$, that is
\[
E(\hat{\theta}) = \theta.
\]
\end{definition}

Intuitively: an unbiased estimator hits the target \emph{on average}. If we repeated the sampling process many times, the mean of all the estimates would coincide with the true parameter.

\begin{methode}[Checking whether an estimator is unbiased]
\begin{enumerate}
\item Write the estimator $\hat{\theta}$ explicitly in terms of the sample variables $X_{1},\dots,X_{n}$.
\item Compute $E(\hat{\theta})$ using the linearity of expectation and the facts $E(X_{i})=\mu$, $\operatorname{Var}(X_{i})=\sigma^{2}$.
\item Compare with $\theta$: if $E(\hat{\theta})=\theta$, the estimator is unbiased; otherwise compute the bias $E(\hat{\theta})-\theta$ and, if possible, \emph{correct} the estimator by adjusting it so that the bias vanishes.
\end{enumerate}
\end{methode}

\courssub{The Sample Mean is an Unbiased Estimator of $\mu$}

\begin{propriete}[Unbiasedness of $\bar{X}$]
Let $X_{1}, \dots, X_{n}$ be a random sample from a population with mean $\mu$. Then
\[
E(\bar{X}) = \mu,
\]
so $\bar{X}$ is an \textbf{unbiased estimator} of $\mu$. (Proof examinable.)
\end{propriete}

\textbf{Proof.} Each $X_{i}$ is drawn from the population, so $E(X_{i}) = \mu$ for every $i$. By linearity of expectation,
\begin{align*}
E(\bar{X}) &= E\!\left(\frac{1}{n}\sum_{i=1}^{n} X_{i}\right) = \frac{1}{n}\sum_{i=1}^{n} E(X_{i}) = \frac{1}{n}\,(n\mu) = \mu. \qquad \blacksquare
\end{align*}

\courssub{Estimating the Variance: Why We Divide by $n-1$}

A natural first guess for estimating $\sigma^{2}$ is the ``average squared deviation from the sample mean'', $\hat{\sigma}^{2} = \dfrac{1}{n}\sum_{i=1}^{n}(X_{i}-\bar{X})^{2}$. Surprisingly, this estimator is \textbf{biased}.

\begin{propriete}[Bias of the variance estimator with divisor $n$]
For a random sample from a population with variance $\sigma^{2}$,
\[
E\!\left(\frac{1}{n}\sum_{i=1}^{n}(X_{i}-\bar{X})^{2}\right) = \frac{n-1}{n}\,\sigma^{2}.
\]
Consequently the \cle{unbiased sample variance} is
\[
s^{2} = \frac{1}{n-1}\sum_{i=1}^{n}(X_{i}-\bar{X})^{2}, \qquad \text{with } E(s^{2}) = \sigma^{2}.
\]
(Proof examinable.)
\end{propriete}

\textbf{Proof.} Write $(X_{i}-\bar{X}) = (X_{i}-\mu) - (\bar{X}-\mu)$. Squaring and summing:
\begin{align*}
\sum_{i=1}^{n}(X_{i}-\bar{X})^{2} &= \sum_{i=1}^{n}(X_{i}-\mu)^{2} - 2(\bar{X}-\mu)\sum_{i=1}^{n}(X_{i}-\mu) + n(\bar{X}-\mu)^{2} \\
&= \sum_{i=1}^{n}(X_{i}-\mu)^{2} - n(\bar{X}-\mu)^{2},
\end{align*}
since $\sum (X_{i}-\mu) = n(\bar{X}-\mu)$. Now take expectations. By definition $E[(X_{i}-\mu)^{2}] = \sigma^{2}$, and since $E(\bar{X})=\mu$ we have $E[(\bar{X}-\mu)^{2}] = \operatorname{Var}(\bar{X}) = \dfrac{\sigma^{2}}{n}$. Hence
\begin{align*}
E\!\left(\sum_{i=1}^{n}(X_{i}-\bar{X})^{2}\right) &= n\sigma^{2} - n\cdot\frac{\sigma^{2}}{n} = (n-1)\sigma^{2},
\end{align*}
and dividing by $n$ gives $\dfrac{n-1}{n}\sigma^{2} \neq \sigma^{2}$. Dividing instead by $n-1$ removes the bias: $E(s^{2}) = \sigma^{2}$. $\blacksquare$

The intuitive explanation is beautiful: the deviations $(X_{i}-\bar{X})$ are measured from $\bar{X}$, which is itself computed \emph{from the same data} and therefore sits, on average, slightly closer to the points than the true $\mu$ does. The squared deviations are thus systematically a little too small, and dividing by $n-1$ instead of $n$ inflates the result just enough to compensate.

\begin{attention}[The divisor trap]
When \textbf{estimating a population variance from a sample}, always divide by $n-1$, not by $n$. Dividing by $n$ underestimates $\sigma^{2}$ by a factor $\frac{n-1}{n}$. Note the contrast: if the data \emph{are} the whole population, the population variance uses divisor $N$; if the data are a \emph{sample} used to estimate $\sigma^{2}$, use $n-1$.
\end{attention}

\begin{exoresolu}[Unbiased estimates from raw data]
A random sample of $n=5$ sacks of garri from a market in Bafoussam gives masses (kg): $48,\ 52,\ 49,\ 51,\ 50$. Obtain unbiased estimates of the population mean and variance.
\tcblower
\textbf{Solution.} The unbiased estimate of $\mu$ is
\[
\bar{x} = \frac{48+52+49+51+50}{5} = \frac{250}{5} = 50 \text{ kg}.
\]
The deviations from $50$ are $-2, 2, -1, 1, 0$, with squares $4,4,1,1,0$, summing to $10$. The unbiased estimate of $\sigma^{2}$ uses the divisor $n-1 = 4$:
\[
s^{2} = \frac{10}{5-1} = 2.5 \text{ kg}^{2}.
\]
(Dividing by $n=5$ would give $2.0$, a biased underestimate.)
\end{exoresolu}

\begin{exoresolu}[Unbiased estimates from summarised data]
A sample of $n=20$ students gives $\sum x = 1460$ and $\sum x^{2} = 107\,300$. Find unbiased estimates of the population mean and variance.
\tcblower
\textbf{Solution.} First, $\bar{x} = \dfrac{1460}{20} = 73$. Using the identity $\sum(x_{i}-\bar{x})^{2} = \sum x^{2} - n\bar{x}^{2}$:
\[
\sum(x_{i}-\bar{x})^{2} = 107\,300 - 20(73)^{2} = 107\,300 - 106\,580 = 720.
\]
Hence the unbiased estimate of the variance is
\[
s^{2} = \frac{720}{20-1} = \frac{720}{19} \approx 37.9.
\]
\end{exoresolu}

\courssub{Efficiency: Choosing the Best Unbiased Estimator}

Unbiasedness alone is not enough: many estimators can be unbiased for the same parameter. For instance, a single observation $X_{1}$ is unbiased for $\mu$ (since $E(X_{1})=\mu$), and so is the sample median of a symmetric population, and so is $\bar{X}$. How do we choose? We prefer the estimator whose values are \emph{least scattered} around $\theta$ --- the one with the smallest variance.

\begin{definition}[Most efficient estimator]
Among all \textbf{unbiased} estimators of a parameter $\theta$, the \cle{most efficient estimator} is the one with the \textbf{smallest variance}. If $\hat{\theta}_{1}$ and $\hat{\theta}_{2}$ are both unbiased and $\operatorname{Var}(\hat{\theta}_{1}) < \operatorname{Var}(\hat{\theta}_{2})$, then $\hat{\theta}_{1}$ is said to be \emph{more efficient} than $\hat{\theta}_{2}$.
\end{definition}

\textbf{Comparison of rival estimators of $\mu$.} For a random sample of size $n$ from a population with variance $\sigma^{2}$:

\begin{center}
\begin{tabularx}{\linewidth}{|l|X|X|}
\hline
\rowcolor{popblueL} \textbf{Estimator} & \textbf{Unbiased?} & \textbf{Variance} \\
\hline
Single observation $X_{1}$ & Yes, $E(X_{1})=\mu$ & $\sigma^{2}$ \\
\hline
Mean of the first $k$ observations ($k<n$) & Yes & $\dfrac{\sigma^{2}}{k}$ \\
\hline
Sample median (symmetric population) & Yes & larger than $\dfrac{\sigma^{2}}{n}$ (about $\dfrac{\pi\sigma^{2}}{2n}$ for a normal population) \\
\hline
Full sample mean $\bar{X}$ & Yes & $\dfrac{\sigma^{2}}{n}$ \\
\hline
\end{tabularx}
\end{center}

Since $\operatorname{Var}(\bar{X}) = \dfrac{\sigma^{2}}{n}$ is the smallest of these variances, the full sample mean is the \textbf{most efficient} of these estimators. Discarding data (using one observation or a subsample) can only inflate the variance.

\begin{attention}[Unbiased does not mean efficient]
An unbiased estimator is \emph{not} automatically the most efficient one. $X_{1}$ is unbiased for $\mu$ but is a poor estimator: its variance $\sigma^{2}$ is $n$ times larger than that of $\bar{X}$. Always check \textbf{both} properties: first $E(\hat{\theta}) = \theta$, then compare variances.
\end{attention}

\begin{popfigure}
\begin{center}
\begin{tikzpicture}
\begin{axis}[
  width=12cm, height=6.5cm,
  domain=0:10, samples=200,
  axis lines=middle,
  xlabel={value of the estimator $\hat{\theta}$},
  ylabel={density},
  xtick={4,6}, xticklabels={$\theta$,$\theta+b$},
  ytick=\empty,
  legend style={at={(0.97,0.95)}, anchor=north east, font=\small},
  every axis plot/.append style={very thick}
]
\addplot[popblue] {1.6*exp(-((x-4)^2)/0.5)};
\addlegendentry{$\hat{\theta}_{1}$: unbiased{,} small variance (efficient)}
\addplot[poporange, dashed] {0.8*exp(-((x-6)^2)/4)};
\addlegendentry{$\hat{\theta}_{2}$: biased{,} large variance}
\draw[popdark, thick] (axis cs:4,0) -- (axis cs:4,1.75);
\node[popdark, anchor=south, font=\small] at (axis cs:4,1.78) {true $\theta$};
\end{axis}
\end{tikzpicture}
\end{center}
\legende{Two rival estimators of $\theta$. $\hat{\theta}_{1}$ is centred on $\theta$ (unbiased) and tightly concentrated (efficient); $\hat{\theta}_{2}$ is centred away from $\theta$ (bias $b$) and widely spread.}
\end{popfigure}

\courssub{Unbiased Estimation of a Population Proportion}

Suppose each member of a population either has an attribute (``smoker'', ``defective item'', ``voter in favour'') or not, and the unknown population proportion is $p$. In a random sample of size $n$ (with replacement, or from a large population), let $X$ be the number of sampled members having the attribute. Then $X \sim B(n, p)$, and the natural estimator is the \cle{sample proportion} $\hat{p} = \dfrac{X}{n}$.

\begin{propriete}[Mean and variance of the sample proportion]
The sample proportion $\hat{p} = X/n$ is an unbiased estimator of $p$, with
\[
E(\hat{p}) = p, \qquad \operatorname{Var}(\hat{p}) = \frac{p(1-p)}{n}.
\]
\end{propriete}

\textbf{Proof.} Since $X \sim B(n,p)$, $E(X) = np$ and $\operatorname{Var}(X) = np(1-p)$. Therefore
\[
E(\hat{p}) = \frac{E(X)}{n} = \frac{np}{n} = p, \qquad
\operatorname{Var}(\hat{p}) = \frac{\operatorname{Var}(X)}{n^{2}} = \frac{np(1-p)}{n^{2}} = \frac{p(1-p)}{n}. \qquad \blacksquare
\]

Note that $\operatorname{Var}(\hat{p}) \to 0$ as $n$ grows: larger samples give more precise estimates of a proportion. The square root of this variance, $\sqrt{p(1-p)/n}$, is the \cle{standard error} of the proportion.

\begin{exoresolu}[Estimating a proportion]
In a random sample of $n = 200$ packets of salt from a factory in Douala, $14$ are underweight. Give an unbiased estimate of the proportion $p$ of underweight packets, and estimate the variance of the estimator.
\tcblower
\textbf{Solution.} The unbiased estimate is
\[
\hat{p} = \frac{X}{n} = \frac{14}{200} = 0.07.
\]
Since $p$ is unknown, we estimate the variance of $\hat{p}$ by replacing $p$ with $\hat{p}$:
\[
\widehat{\operatorname{Var}}(\hat{p}) = \frac{\hat{p}(1-\hat{p})}{n} = \frac{0.07 \times 0.93}{200} \approx 0.000326,
\]
so the standard error is about $\sqrt{0.000326} \approx 0.018$.
\end{exoresolu}

\begin{aretenir}[Key points of this section]
\begin{itemize}
\item An \textbf{estimator} $\hat{\theta}$ is a random variable; an \textbf{estimate} is its value on one sample.
\item $\operatorname{Bias}(\hat{\theta}) = E(\hat{\theta}) - \theta$; $\hat{\theta}$ is \textbf{unbiased} when $E(\hat{\theta}) = \theta$.
\item $E(\bar{X}) = \mu$: the sample mean is an unbiased estimator of the population mean.
\item $E\!\left(\frac{1}{n}\sum(X_{i}-\bar{X})^{2}\right) = \frac{n-1}{n}\sigma^{2}$: biased. The \textbf{unbiased} sample variance is $s^{2} = \dfrac{\sum(X_{i}-\bar{X})^{2}}{n-1}$.
\item The \textbf{most efficient} unbiased estimator is the one with the smallest variance; for $\mu$, $\bar{X}$ with $\operatorname{Var}(\bar{X}) = \sigma^{2}/n$ beats a single observation, a subsample mean or the median.
\item For a proportion, $\hat{p} = X/n$ satisfies $E(\hat{p}) = p$ and $\operatorname{Var}(\hat{p}) = \dfrac{p(1-p)}{n}$.
\end{itemize}
\end{aretenir}

\begin{exercice}[Checking unbiasedness and efficiency]
A population has mean $\mu$ and variance $\sigma^{2}$. A random sample $X_{1}, X_{2}, X_{3}$ of size $3$ is taken. Consider the two estimators
\[
T_{1} = \frac{X_{1}+X_{2}+X_{3}}{3}, \qquad T_{2} = \frac{X_{1}+2X_{2}+X_{3}}{4}.
\]
\begin{enumerate}
\item Show that both $T_{1}$ and $T_{2}$ are unbiased estimators of $\mu$.
\item Compute $\operatorname{Var}(T_{1})$ and $\operatorname{Var}(T_{2})$, and state which estimator is more efficient.
\item Is $T_{3} = \dfrac{X_{1}+X_{2}}{3}$ unbiased for $\mu$? Justify.
\end{enumerate}
\end{exercice}

\begin{corrige}[Exercise 1]
\begin{enumerate}
\item $E(T_{1}) = \dfrac{1}{3}(E(X_{1})+E(X_{2})+E(X_{3})) = \dfrac{3\mu}{3} = \mu$. Also $E(T_{2}) = \dfrac{1}{4}(\mu + 2\mu + \mu) = \dfrac{4\mu}{4} = \mu$. Both are unbiased.
\item By independence, $\operatorname{Var}(T_{1}) = \dfrac{1}{9}(3\sigma^{2}) = \dfrac{\sigma^{2}}{3}$, and $\operatorname{Var}(T_{2}) = \dfrac{1}{16}(\sigma^{2} + 4\sigma^{2} + \sigma^{2}) = \dfrac{6\sigma^{2}}{16} = \dfrac{3\sigma^{2}}{8}$. Since $\dfrac{\sigma^{2}}{3} < \dfrac{3\sigma^{2}}{8}$, $T_{1}$ (the sample mean) is more efficient.
\item $E(T_{3}) = \dfrac{2\mu}{3} \neq \mu$: $T_{3}$ is biased, with bias $-\dfrac{\mu}{3}$.
\end{enumerate}
\end{corrige}

\begin{exercice}[Unbiased estimates from data]
The masses (g) of a random sample of $6$ bread loaves from a bakery in Yaound\'e are: $498,\ 502,\ 505,\ 497,\ 500,\ 504$. Compute unbiased estimates of the population mean and variance.
\end{exercice}

\begin{corrige}[Exercise 2]
$\bar{x} = \dfrac{498+502+505+497+500+504}{6} = \dfrac{3006}{6} = 501$ g. The deviations from $501$ are $-3, 1, 4, -4, -1, 3$; their squares are $9, 1, 16, 16, 1, 9$, summing to $52$. Hence $s^{2} = \dfrac{52}{6-1} = 10.4\ \mathrm{g}^{2}$ (dividing by $6$ would give the biased value $8.67$).
\end{corrige}

\begin{exercice}[Constructing an unbiased estimator]
A random sample $X_{1}, X_{2}$ of size $2$ is taken from a population with mean $\mu$ and variance $\sigma^{2}$. The estimator $W = aX_{1} + bX_{2}$ is proposed for $\mu$, where $a$ and $b$ are constants.
\begin{enumerate}
\item Show that $W$ is unbiased for $\mu$ if and only if $a + b = 1$.
\item Assuming $a + b = 1$, show that $\operatorname{Var}(W) = (a^{2} + b^{2})\sigma^{2}$, and deduce the values of $a$ and $b$ that make $W$ the most efficient unbiased estimator of this form. Comment on your answer.
\end{enumerate}
\end{exercice}

\begin{corrige}[Exercise 3]
\begin{enumerate}
\item $E(W) = aE(X_{1}) + bE(X_{2}) = (a+b)\mu$. This equals $\mu$ for all $\mu$ if and only if $a + b = 1$.
\item By independence, $\operatorname{Var}(W) = a^{2}\sigma^{2} + b^{2}\sigma^{2} = (a^{2}+b^{2})\sigma^{2}$. With $b = 1 - a$, we minimise $f(a) = a^{2} + (1-a)^{2} = 2a^{2} - 2a + 1$: $f'(a) = 4a - 2 = 0$ gives $a = \tfrac{1}{2}$, hence $b = \tfrac{1}{2}$. The most efficient choice is $W = \dfrac{X_{1}+X_{2}}{2} = \bar{X}$: once again, the sample mean is the best unbiased linear combination of the observations.
\end{enumerate}
\end{corrige}

\begin{savaistu}[Why $n-1$?]
The divisor $n-1$ is called the number of \emph{degrees of freedom}: once $\bar{x}$ is fixed, the $n$ deviations $(x_{i}-\bar{x})$ satisfy $\sum(x_{i}-\bar{x}) = 0$, so only $n-1$ of them are free to vary. This idea, popularised by the statistician R. A. Fisher in the 1920s, will reappear when we build confidence intervals with the $t$-distribution later in this chapter.
\end{savaistu}

\coursec{Pooled Estimation from Two Samples}

In the previous section we learnt how to estimate a population parameter from \emph{one} sample, and how to choose the \cle{most efficient} estimator among several unbiased candidates. In practice, however, an investigator often has \emph{two} samples available, taken from the same population (or from two populations known to share a common parameter). Think of the same Biology test marked separately in Upper Sixth A and Upper Sixth B of a school in Yaound\'e, or bags of cement weighed from two production lines of the same factory in Douala. If the two samples carry information about the \emph{same} unknown quantity, it would be wasteful to use only one of them. The natural question is: \emph{how do we combine the two samples into a single, better estimate?} This is exactly what \cle{pooled estimation} achieves.

\courssub{The Context: Two Samples, One Parameter}

\begin{definition}[Pooled estimation]
Let $X_1, X_2, \dots, X_{n_1}$ be a random sample of size $n_1$ and $Y_1, Y_2, \dots, Y_{n_2}$ an independent random sample of size $n_2$, both drawn from populations having a \emph{common} parameter $\theta$ (mean $\mu$, variance $\sigma^2$, or proportion $p$). A \cle{pooled estimator} of $\theta$ is a single estimator constructed by combining the information from both samples.
\end{definition}

The key assumption is that the parameter is \emph{common}. We may pool the two class means to estimate the mean mark of the whole Upper Sixth only if both classes are drawn from the same achievement population. Pooling the mean height of boys with the mean height of girls to estimate ``the mean height of students'' would be meaningless: the parameters are different.

\begin{attention}[When pooling is not allowed]
Never pool two samples when there is evidence that the populations have \emph{different} parameters. Pooling two means that estimate different quantities $\mu_1 \neq \mu_2$ produces a number that estimates neither. In an examination, look for phrases such as ``the two machines are identical'', ``the same population'', or ``a common variance may be assumed'' before pooling.
\end{attention}

\courssub{Pooled Estimate of a Common Population Mean}

Suppose both samples come from a population with mean $\mu$. The first sample gives mean $\bar{x}_1$, the second gives $\bar{x}_2$. Each is unbiased for $\mu$. How should we combine them? A simple average $\tfrac{1}{2}(\bar{x}_1 + \bar{x}_2)$ treats both samples equally, even if one is much larger than the other. Intuition says the bigger sample deserves more \emph{weight}, because its mean is more reliable. If we simply merge the two samples into one big sample of size $n_1 + n_2$, the combined mean is the total of all observations divided by the total count.

\begin{propriete}[Pooled estimator of the common mean]
If two independent samples of sizes $n_1$ and $n_2$, with means $\bar{x}_1$ and $\bar{x}_2$, are drawn from populations with common mean $\mu$, the pooled estimator is
\[
\hat{\mu} = \bar{x}_p = \frac{n_1\bar{x}_1 + n_2\bar{x}_2}{n_1 + n_2}.
\]
It is a \cle{weighted average} of the two sample means, with weights proportional to the sample sizes, and it is \textbf{unbiased} for $\mu$.
\end{propriete}

\textbf{Proof of unbiasedness (examinable).} Since each sample mean is unbiased, $E(\bar{X}_1) = E(\bar{X}_2) = \mu$. By linearity of expectation,
\[
E(\bar{X}_p) = \frac{n_1 E(\bar{X}_1) + n_2 E(\bar{X}_2)}{n_1 + n_2} = \frac{n_1\mu + n_2\mu}{n_1+n_2} = \frac{(n_1+n_2)\mu}{n_1+n_2} = \mu. \qquad \blacksquare
\]

\textbf{Why is the pooled mean more efficient?} Assume the common variance is $\sigma^2$ and the samples are independent. Then
\[
\operatorname{Var}(\bar{X}_p) = \frac{n_1^2 \operatorname{Var}(\bar{X}_1) + n_2^2 \operatorname{Var}(\bar{X}_2)}{(n_1+n_2)^2} = \frac{n_1^2 \cdot \frac{\sigma^2}{n_1} + n_2^2 \cdot \frac{\sigma^2}{n_2}}{(n_1+n_2)^2} = \frac{\sigma^2}{n_1+n_2}.
\]
Since $\dfrac{\sigma^2}{n_1+n_2} < \dfrac{\sigma^2}{n_1}$ and $< \dfrac{\sigma^2}{n_2}$, the pooled mean has a \emph{smaller variance} than either separate mean: it is \cle{more efficient}. Among all weighted averages $w\bar{X}_1 + (1-w)\bar{X}_2$, minimising the variance with respect to $w$ gives precisely $w = \dfrac{n_1}{n_1+n_2}$: the pooled mean is the \emph{best} linear combination. More data, less uncertainty --- exactly as expected.

\courssub{Pooled Estimate of a Common Population Variance}

Suppose now that two populations share a common variance $\sigma^2$ (a very common assumption when comparing two means). Each sample gives an unbiased estimate $s_1^2$ and $s_2^2$ of $\sigma^2$, where
\[
s_1^2 = \frac{\sum (x - \bar{x}_1)^2}{n_1 - 1}, \qquad s_2^2 = \frac{\sum (y - \bar{y}_2)^2}{n_2 - 1}.
\]
Each sum of squared deviations carries $n_i - 1$ \cle{degrees of freedom} (one is lost because the sample mean was estimated from the same data). To combine, we add the sums of squares and divide by the \emph{total} degrees of freedom.

\begin{propriete}[Pooled estimator of the common variance]
The pooled estimator of a common population variance $\sigma^2$ is
\[
s_p^2 = \frac{\sum (x - \bar{x}_1)^2 + \sum (y - \bar{y}_2)^2}{n_1 + n_2 - 2} = \frac{(n_1-1)s_1^2 + (n_2-1)s_2^2}{n_1 + n_2 - 2}.
\]
It is a weighted average of the two sample variances with weights equal to their degrees of freedom, and it is \textbf{unbiased} for $\sigma^2$.
\end{propriete}

\textbf{Proof of unbiasedness (instructive).} Recall that $E(s_i^2) = \sigma^2$, i.e.\ $E\big[\sum (x-\bar{x}_1)^2\big] = (n_1-1)\sigma^2$. Hence
\[
E(s_p^2) = \frac{(n_1-1)\sigma^2 + (n_2-1)\sigma^2}{n_1+n_2-2} = \frac{(n_1+n_2-2)\sigma^2}{n_1+n_2-2} = \sigma^2. \qquad \blacksquare
\]

\begin{attention}[Do not average the variances directly]
A frequent error is to compute $\tfrac{1}{2}(s_1^2 + s_2^2)$. This is only correct when $n_1 = n_2$. The correct weights are the \textbf{degrees of freedom} $n_1 - 1$ and $n_2 - 1$, not the sample sizes, and certainly not equal weights. Likewise, the denominator is $n_1 + n_2 - 2$, \emph{not} $n_1 + n_2$ and not $n_1 + n_2 - 1$: two means were estimated, so two degrees of freedom are lost.
\end{attention}

\begin{methode}[Recovering $\sum(x-\bar{x})^2$ from a quoted variance --- check the divisor!]
Examination questions often summarise a sample by giving ``a variance of $12$'' without saying which divisor was used. Before pooling, you must recover the sum of squared deviations:
\begin{enumerate}
\item \textbf{Identify the divisor.} If the variance was computed with divisor $n$ (sometimes denoted $\hat{\sigma}^2$ or ``the variance of the data''), then $\sum(x-\bar{x})^2 = n \times \text{(quoted variance)}$.
\item If the divisor was $n-1$ (unbiased sample variance $s^2$), then $\sum(x-\bar{x})^2 = (n-1) \times s^2$.
\item \textbf{Add} the two sums of squared deviations and \textbf{divide} by $n_1 + n_2 - 2$.
\item As a check, convert everything to the unbiased form first: $s^2 = \dfrac{n}{n-1}\hat{\sigma}^2$, then apply $s_p^2 = \dfrac{(n_1-1)s_1^2 + (n_2-1)s_2^2}{n_1+n_2-2}$.
\end{enumerate}
\end{methode}

\courssub{Pooled Estimate of a Common Proportion}

The same logic applies to proportions. If $x_1$ successes are observed in $n_1$ trials and $x_2$ successes in $n_2$ independent trials, all with the same success probability $p$, then merging the trials gives $x_1 + x_2$ successes out of $n_1 + n_2$ trials.

\begin{propriete}[Pooled estimator of a common proportion]
\[
\hat{p} = \frac{x_1 + x_2}{n_1 + n_2} = \frac{n_1\hat{p}_1 + n_2\hat{p}_2}{n_1+n_2},
\]
where $\hat{p}_i = x_i/n_i$. It is unbiased, since $E(\hat{p}) = \dfrac{n_1 p + n_2 p}{n_1+n_2} = p$.
\end{propriete}

For example, if a quality inspector finds $6$ defective items in $100$ bags from machine A and $9$ defective in $150$ bags from an identical machine B, the pooled estimate of the common defect rate is $\hat{p} = \dfrac{6+9}{100+150} = \dfrac{15}{250} = 0.06$.

\courssub{Summary Flow Diagram}

\begin{popfigure}
\begin{tikzpicture}[node distance=0.55cm and 1.1cm,
box/.style={draw=popblue, thick, rounded corners, fill=popblueL, text width=3.6cm, align=center, minimum height=1.1cm, font=\small},
sumbox/.style={draw=poppurple, thick, rounded corners, fill=white, text width=3.8cm, align=center, minimum height=1.3cm, font=\small},
resbox/.style={draw=popgreen, thick, rounded corners, fill=popgreenL, text width=4.6cm, align=center, minimum height=1.3cm, font=\small},
arr/.style={-{stealth}, thick, popdark}]
\node[box, align=center] (s1){Sample 1\\ size $n_1$};
\node[box, below=of s1, align=center] (s2){Sample 2\\ size $n_2$};
\node[sumbox, right=of s1] (m1) {$\bar{x}_1$,\quad $\sum(x-\bar{x}_1)^2$};
\node[sumbox, right=of s2] (m2) {$\bar{x}_2$,\quad $\sum(y-\bar{y}_2)^2$};
\node[resbox, right=2.2cm of $(m1)!0.5!(m2)$, align=center] (pool){Pooled estimates\\[2pt] $\hat{\mu}=\dfrac{n_1\bar{x}_1+n_2\bar{x}_2}{n_1+n_2}$\\[4pt] $s_p^2=\dfrac{(n_1-1)s_1^2+(n_2-1)s_2^2}{n_1+n_2-2}$};
\draw[arr] (s1) -- (m1);
\draw[arr] (s2) -- (m2);
\draw[arr] (m1.east) -- (pool.170);
\draw[arr] (m2.east) -- (pool.190);
\end{tikzpicture}
\legende{From two independent samples to pooled estimates of the common mean $\mu$ and common variance $\sigma^2$.}
\end{popfigure}

\courssub{Worked Examples}

\begin{exoresolu}[Pooled mean and pooled variance from raw summaries]
Two independent random samples are taken from the same population of masses (in grams) of packets of sugar produced by a factory in Bafoussam.

Sample 1: $n_1 = 10$, $\bar{x}_1 = 502.4$, $\sum(x-\bar{x}_1)^2 = 58.5$.

Sample 2: $n_2 = 15$, $\bar{x}_2 = 501.8$, $\sum(y-\bar{y}_2)^2 = 84.0$.

Find unbiased pooled estimates of the population mean and variance.
\tcblower
\textbf{Pooled mean:}
\[
\hat{\mu} = \frac{n_1\bar{x}_1 + n_2\bar{x}_2}{n_1+n_2} = \frac{10(502.4) + 15(501.8)}{25} = \frac{5024 + 7527}{25} = \frac{12551}{25} = 502.04\ \mathrm{g}.
\]
\textbf{Pooled variance:} the sums of squared deviations are given directly, so we add them and divide by $n_1+n_2-2 = 23$:
\[
s_p^2 = \frac{58.5 + 84.0}{10+15-2} = \frac{142.5}{23} \approx 6.196\ \mathrm{g}^2.
\]
Hence $\hat{\mu} = 502.04\ \mathrm{g}$ and $s_p^2 \approx 6.20\ \mathrm{g}^2$ (so $s_p \approx 2.49\ \mathrm{g}$). Note that the pooled mean lies between the two sample means, closer to $\bar{x}_2$ because sample 2 is larger.
\end{exoresolu}

\begin{exoresolu}[GCE-style: variances quoted with different divisors]
A random sample of $8$ seedlings from nursery A has mean height $24.5\ \mathrm{cm}$ and variance (computed with divisor $n$) $9.6\ \mathrm{cm}^2$. An independent random sample of $12$ seedlings from nursery B has mean height $25.1\ \mathrm{cm}$ and unbiased sample variance $s_2^2 = 13.2\ \mathrm{cm}^2$. Assuming the two nurseries produce seedlings with a common mean height and common variance, obtain unbiased pooled estimates of the mean and the variance.
\tcblower
\textbf{Step 1: recover the sums of squared deviations.}

For sample A the divisor was $n_1 = 8$:
\[
\sum(x-\bar{x}_1)^2 = n_1 \times 9.6 = 8 \times 9.6 = 76.8.
\]
For sample B the divisor was $n_2 - 1 = 11$:
\[
\sum(y-\bar{y}_2)^2 = (n_2-1)s_2^2 = 11 \times 13.2 = 145.2.
\]
\textbf{Step 2: pooled mean.}
\[
\hat{\mu} = \frac{8(24.5) + 12(25.1)}{20} = \frac{196 + 301.2}{20} = \frac{497.2}{20} = 24.86\ \mathrm{cm}.
\]
\textbf{Step 3: pooled variance.}
\[
s_p^2 = \frac{76.8 + 145.2}{8+12-2} = \frac{222.0}{18} \approx 12.33\ \mathrm{cm}^2.
\]
\textbf{Check:} converting sample A to the unbiased form, $s_1^2 = \dfrac{8}{7}\times 9.6 \approx 10.97$; then $\dfrac{7(10.97)+11(13.2)}{18} = \dfrac{76.8+145.2}{18}$, the same result. The naive average $\tfrac12(9.6+13.2) = 11.4$ would be wrong on two counts: mixed divisors and equal weights.
\end{exoresolu}

\begin{exercice}[Two shifts in a bakery]
The masses (g) of loaves from the morning shift of a bakery in Bamenda give: $n_1 = 12$, $\bar{x}_1 = 798$, $s_1^2 = 36$ (unbiased). The evening shift gives: $n_2 = 18$, $\bar{x}_2 = 803$, $s_2^2 = 50$ (unbiased). Assuming a common mean and common variance, find pooled estimates of both.
\end{exercice}

\begin{corrige}[Exercise 1]
Pooled mean: $\hat{\mu} = \dfrac{12(798) + 18(803)}{30} = \dfrac{9576 + 14454}{30} = \dfrac{24030}{30} = 801\ \mathrm{g}$.

Pooled variance: $s_p^2 = \dfrac{11(36) + 17(50)}{12+18-2} = \dfrac{396 + 850}{28} = \dfrac{1246}{28} = 44.5\ \mathrm{g}^2$.
\end{corrige}

\begin{aretenir}[Key points]
\begin{itemize}
\item Pool only when the parameter is genuinely \cle{common} to both populations.
\item Pooled mean: $\hat{\mu} = \dfrac{n_1\bar{x}_1 + n_2\bar{x}_2}{n_1+n_2}$ --- unbiased, and more efficient than either mean alone since $\operatorname{Var}(\hat{\mu}) = \dfrac{\sigma^2}{n_1+n_2}$.
\item Pooled variance: $s_p^2 = \dfrac{(n_1-1)s_1^2 + (n_2-1)s_2^2}{n_1+n_2-2}$ --- weight by \textbf{degrees of freedom}, denominator $n_1+n_2-2$.
\item Pooled proportion: $\hat{p} = \dfrac{x_1+x_2}{n_1+n_2}$.
\item When a variance is quoted, always check whether the divisor was $n$ or $n-1$ before recovering $\sum(x-\bar{x})^2$.
\end{itemize}
\end{aretenir}

\begin{savaistu}[Why $n_1 + n_2 - 2$?]
The pooled variance is the engine behind the \emph{two-sample $t$-test}, one of the most used statistical procedures in agriculture, medicine and industry --- for instance when IRAD researchers compare the yields of two maize varieties grown under the same conditions. The divisor $n_1+n_2-2$ reflects the total degrees of freedom: $n_1 + n_2$ data values, minus the two sample means that had to be estimated first. Respecting degrees of freedom is what keeps the estimator unbiased.
\end{savaistu}

\coursec{Interval Estimation: Confidence Intervals}

In the previous section we learned how to combine two independent samples to obtain a single, more precise \emph{point estimate} of a population mean or variance. A point estimate, however, is just one number — it gives no indication of the uncertainty attached to it. In practice, a decision-maker (a quality-control engineer at a brewery in Douala, a public-health official estimating malaria prevalence in the Far North, or a GCE examiner setting grade boundaries) needs a \emph{range} of plausible values for the parameter. This is the purpose of \textbf{interval estimation}.

\courssub{From Point Estimation to Interval Estimation}

Suppose we take a random sample of size $n$ from a population and compute the sample mean $\bar{x}$. We know that $\bar{x}$ is an unbiased estimator of the population mean $\mu$, but we also know that a different sample would give a different $\bar{x}$. The \textbf{sampling distribution} of $\bar{x}$ quantifies this variability. If we can describe that distribution, we can construct an interval that, in a well-defined sense, \emph{captures} the true $\mu$ with a prescribed probability.

\begin{definition}[Confidence Interval and Confidence Level]
A \cle{confidence interval (CI)} for a population parameter $\theta$ is a random interval $(L, U)$, computed from the sample data, such that
\[
P(L < \theta < U) = 1 - \alpha
\]
where $\alpha \in (0,1)$ is the \cle{significance level}. The value $1-\alpha$ (often expressed as a percentage: $90\%$, $95\%$, $99\%$) is called the \cle{confidence level}.
\end{definition}

\begin{attention}[The Correct Interpretation]
Once a particular sample is drawn and the numbers $L$ and $U$ are calculated, the interval $(L, U)$ is \textbf{fixed}. The parameter $\theta$ is also fixed (though unknown). It is \textbf{incorrect} to say ``There is a $95\%$ probability that $\mu$ lies in this interval.'' The correct interpretation is:
\begin{quote}
``If we were to repeat the sampling process a very large number of times and construct a $95\%$ confidence interval each time, approximately $95\%$ of those intervals would contain the true value of $\mu$.''
\end{quote}
The confidence level refers to the \emph{long-run performance of the method}, not to the probability attached to a single realised interval.
\end{attention}

\begin{popfigure}
\centering
\begin{tikzpicture}
\begin{axis}[
    width=0.9\linewidth,
    height=6cm,
    domain=-4:4,
    samples=200,
    axis lines=middle,
    xlabel=$z$,
    ylabel={$f(z)$},
    xtick={-1.96,0,1.96},
    xticklabels={$-1.96$,$0$,$1.96$},
    ytick=\empty,
    xmin=-4, xmax=4,
    ymin=0, ymax=0.45,
    clip=false
]
% full curve
\addplot [popblue, thick] {exp(-x^2/2)/sqrt(max(0,2*pi))};
% shaded central 95%
\addplot [popblueL, fill=popblueL, opacity=0.4, domain=-1.96:1.96] {exp(-x^2/2)/sqrt(max(0,2*pi))} \closedcycle;
% tail labels
\node[popdark, anchor=north] at (axis cs:-3,0.02) {$2.5\%$};
\node[popdark, anchor=north] at (axis cs:3,0.02) {$2.5\%$};
\node[popdark, anchor=south] at (axis cs:0,0.38) {$95\%$};
% vertical lines at critical values
\draw [poporange, dashed, thick] (axis cs:-1.96,0) -- (axis cs:-1.96,0.058);
\draw [poporange, dashed, thick] (axis cs:1.96,0) -- (axis cs:1.96,0.058);
\end{axis}
\end{tikzpicture}
\legende{Standard normal density with the central $95\%$ shaded. The critical values $\pm z_{0.025} = \pm 1.96$ leave $2.5\%$ in each tail.}
\end{popfigure}

\courssub{General Structure of a Confidence Interval}

All the confidence intervals in this syllabus share the same architecture:

\[
\textbf{Point Estimate} \;\pm\; \textbf{(Critical Value)} \times \textbf{(Standard Error of the Estimate)}
\]

\begin{itemize}
\item The \textbf{point estimate} is the unbiased estimator of the parameter (e.g. $\bar{x}$ for $\mu$, $\bar{x}_1-\bar{x}_2$ for $\mu_1-\mu_2$, $\hat{p}$ for $p$).
\item The \textbf{critical value} comes from the sampling distribution of the estimator. For large samples or known variance we use the standard normal quantile $z_{\alpha/2}$.
\item The \textbf{standard error (SE)} is the standard deviation of the sampling distribution of the estimator (e.g. $\sigma/\sqrt{n}$ for $\bar{x}$, $\sqrt{\sigma_1^2/n_1+\sigma_2^2/n_2}$ for $\bar{x}_1-\bar{x}_2$).
\end{itemize}

\courssub{Confidence Interval for $\mu$ when $\sigma^2$ is Known}

\begin{propriete}[CI for $\mu$, $\sigma^2$ Known]
Let $X_1,\dots,X_n$ be a random sample from a population with \textbf{known} variance $\sigma^2$. If the population is normal, or $n$ is large enough for the Central Limit Theorem to apply, the sampling distribution of $\bar{X}$ is $N(\mu,\,\sigma^2/n)$. A $(1-\alpha)100\%$ confidence interval for $\mu$ is
\[
\boxed{\bar{x} \;\pm\; z_{\alpha/2}\,\frac{\sigma}{\sqrt{n}}}
\]
where $z_{\alpha/2}$ is the upper $\alpha/2$ quantile of the standard normal distribution.
\end{propriete}

\begin{aretenir}[Common Critical Values]
\begin{center}
\begin{tabular}{|c|c|c|}\hline
\rowcolor{popblueL}
Confidence Level & $\alpha$ & $z_{\alpha/2}$ \\ \hline
$90\%$ & $0.10$ & $1.645$ \\
$95\%$ & $0.05$ & $1.96$ \\
$99\%$ & $0.01$ & $2.576$ \\ \hline
\end{tabular}
\end{center}
These values are read from the standard normal tables provided in the examination.
\end{aretenir}

\begin{exoresolu}[CI for Mean with Known Variance]
The weights of bags of cement produced by a factory in Douala are normally distributed with a standard deviation of $\SI{1.2}{kg}$. A random sample of $25$ bags has a mean weight of $\SI{50.4}{kg}$. Construct a $95\%$ confidence interval for the true mean weight $\mu$.
\tcblower
\textbf{Solution:}
\begin{itemize}
\item Given: $\sigma = 1.2$, $n = 25$, $\bar{x} = 50.4$, confidence level $95\% \Rightarrow \alpha = 0.05$, $z_{0.025} = 1.96$.
\item Standard error: $\displaystyle \frac{\sigma}{\sqrt{n}} = \frac{1.2}{\sqrt{25}} = 0.24$.
\item Margin of error: $E = z_{\alpha/2} \times \text{SE} = 1.96 \times 0.24 = 0.4704$.
\item Confidence interval: $50.4 \pm 0.4704$, i.e. $(49.9296,\; 50.8704)$.
\end{itemize}
\textbf{Interpretation:} We are $95\%$ confident that the true mean weight of all bags lies between $\SI{49.93}{kg}$ and $\SI{50.87}{kg}$.
\end{exoresolu}

\courssub{Confidence Interval for $\mu$ when $\sigma^2$ is Unknown (Large Samples)}

In most practical situations the population variance $\sigma^2$ is unknown. For \textbf{large samples} ($n \ge 30$ is the usual rule of thumb in the GCE syllabus), the Central Limit Theorem guarantees that
\[
\frac{\bar{X} - \mu}{S/\sqrt{n}} \approx N(0,1)
\]
where $S$ is the sample standard deviation (the unbiased estimator of $\sigma$, i.e. $S = \sqrt{\frac{n}{n-1}}\,s$ with $s$ the usual sample standard deviation). Replacing $\sigma$ by $S$ (or $s$, the difference is negligible for large $n$) in the formula gives an approximate confidence interval.

\begin{propriete}[CI for $\mu$, $\sigma^2$ Unknown, Large $n$]
For a large random sample ($n \ge 30$) from any population with finite variance, an approximate $(1-\alpha)100\%$ confidence interval for $\mu$ is
\[
\boxed{\bar{x} \;\pm\; z_{\alpha/2}\,\frac{s}{\sqrt{n}}}
\]
where $s$ is the sample standard deviation.
\end{propriete}

\begin{attention}[Small Samples — Beyond the Syllabus]
If $n < 30$ and $\sigma$ is unknown, the exact sampling distribution of $\frac{\bar{X}-\mu}{S/\sqrt{n}}$ is the \textbf{Student's $t$ distribution} with $n-1$ degrees of freedom. The critical value $t_{\alpha/2,\,n-1}$ replaces $z_{\alpha/2}$. This topic is \textbf{not required} for the GCE Advanced Level Further Mathematics syllabus, but you may encounter it in university courses. In GCE questions, if $n < 30$ and $\sigma$ is unknown, the question will either state that the population is normal and provide the $t$-value, or it will be a large-sample scenario.
\end{attention}

\begin{exoresolu}[CI for Mean with Unknown Variance, Large Sample]
A random sample of $40$ students from a large secondary school in Yaoundé gave a mean mathematics score of $62.5$ with a standard deviation of $12.3$. Find a $99\%$ confidence interval for the population mean score.
\tcblower
\textbf{Solution:}
\begin{itemize}
\item $n = 40 \ge 30$ (large sample), $\bar{x} = 62.5$, $s = 12.3$, $99\% \Rightarrow \alpha = 0.01$, $z_{0.005} = 2.576$.
\item Standard error: $\displaystyle \frac{s}{\sqrt{n}} = \frac{12.3}{\sqrt{40}} \approx 1.945$.
\item Margin of error: $E = 2.576 \times 1.945 \approx 5.01$.
\item CI: $62.5 \pm 5.01 \Rightarrow (57.49,\; 67.51)$.
\end{itemize}
\end{exoresolu}

\courssub{Determining the Sample Size for a Required Interval Width}

Often we want to design a study so that the confidence interval has a specified maximum width $W$ (or margin of error $E = W/2$). For a mean with known $\sigma$ (or a good estimate from a pilot study), the width is
\[
W = 2\,z_{\alpha/2}\,\frac{\sigma}{\sqrt{n}}.
\]
Solving for $n$:
\[
n = \left(\frac{2\,z_{\alpha/2}\,\sigma}{W}\right)^2 = \left(\frac{z_{\alpha/2}\,\sigma}{E}\right)^2.
\]
Always \textbf{round up} to the next integer, because a smaller $n$ would give a wider interval.

\begin{methode}[Finding the Required Sample Size for a Mean]
\begin{enumerate}
\item Identify the desired confidence level and obtain $z_{\alpha/2}$.
\item Specify the maximum margin of error $E$ (half the desired width).
\item Obtain a planning value for $\sigma$ (from a pilot study, past data, or a conservative guess).
\item Compute $n = \left(\dfrac{z_{\alpha/2}\,\sigma}{E}\right)^2$.
\item Round \textbf{up} to the nearest whole number.
\end{enumerate}
\end{methode}

\begin{exoresolu}[Sample Size Determination]
A food-processing company wants to estimate the mean weight of tinned tomatoes to within $\SI{2}{g}$ with $95\%$ confidence. Past data suggest $\sigma \approx \SI{8}{g}$. What sample size is needed?
\tcblower
$z_{0.025}=1.96$, $E=2$, $\sigma=8$.
\[
n = \left(\frac{1.96 \times 8}{2}\right)^2 = (7.84)^2 = 61.4656 \Rightarrow n = 62 \text{ (round up)}.
\]
\end{exoresolu}

\courssub{Confidence Interval for the Difference of Two Means $\mu_1 - \mu_2$}

We now have two independent random samples:
\begin{itemize}
\item Sample 1: size $n_1$, mean $\bar{x}_1$, variance $\sigma_1^2$ (known or estimated by $s_1^2$).
\item Sample 2: size $n_2$, mean $\bar{x}_2$, variance $\sigma_2^2$ (known or estimated by $s_2^2$).
\end{itemize}
The point estimate of $\mu_1-\mu_2$ is $\bar{x}_1-\bar{x}_2$. Its sampling distribution has mean $\mu_1-\mu_2$ and variance $\frac{\sigma_1^2}{n_1}+\frac{\sigma_2^2}{n_2}$. For large samples (or normal populations with known variances) the standardised statistic is approximately $N(0,1)$.

\begin{propriete}[CI for $\mu_1-\mu_2$, Variances Known or Large Samples]
An approximate $(1-\alpha)100\%$ confidence interval for $\mu_1-\mu_2$ is
\[
\boxed{(\bar{x}_1-\bar{x}_2) \;\pm\; z_{\alpha/2}\,\sqrt{\frac{\sigma_1^2}{n_1}+\frac{\sigma_2^2}{n_2}}}
\]
If $\sigma_1^2,\sigma_2^2$ are unknown but both samples are large ($n_1,n_2 \ge 30$), replace $\sigma_1^2,\sigma_2^2$ by the sample variances $s_1^2,s_2^2$.
\end{propriete}

\begin{attention}[Do Not Confuse the Standard Errors!]
\begin{itemize}
\item For \textbf{one mean}: $\text{SE} = \dfrac{\sigma}{\sqrt{n}}$ (or $\dfrac{s}{\sqrt{n}}$).
\item For \textbf{difference of two means}: $\text{SE} = \sqrt{\dfrac{\sigma_1^2}{n_1}+\dfrac{\sigma_2^2}{n_2}}$ (or $\sqrt{\dfrac{s_1^2}{n_1}+\dfrac{s_2^2}{n_2}}$).
\end{itemize}
A common error is to use $\sqrt{\frac{\sigma_1^2}{n_1}+\frac{\sigma_2^2}{n_2}}$ when the question asks for a CI for a \emph{single} mean, or vice versa. Read the wording carefully: ``mean of population A'' vs ``difference between the means of population A and population B''.
\end{attention}

\begin{exoresolu}[CI for Difference of Two Means]
Two varieties of maize (Variety A and Variety B) are tested in the Northwest Region. Yields (in tonnes/hectare) from independent random samples are:
\begin{center}
\begin{tabular}{|c|c|c|c|}\hline
\rowcolor{popblueL}
Variety & Sample size & Sample mean & Sample variance \\ \hline
A & 35 & 4.2 & 0.64 \\
B & 40 & 3.8 & 0.49 \\ \hline
\end{tabular}
\end{center}
Construct a $95\%$ confidence interval for the difference in mean yields $\mu_A-\mu_B$.
\tcblower
\textbf{Solution:}
Both samples are large ($n_1=35, n_2=40$), so we use sample variances.
$\bar{x}_A-\bar{x}_B = 4.2-3.8 = 0.4$.
$\text{SE} = \sqrt{\frac{0.64}{35}+\frac{0.49}{40}} = \sqrt{0.0182857+0.01225} = \sqrt{0.0305357} \approx 0.1747$.
$z_{0.025}=1.96$.
Margin of error $E = 1.96 \times 0.1747 \approx 0.3424$.
CI: $0.4 \pm 0.3424 \Rightarrow (0.0576,\; 0.7424)$.
\textbf{Interpretation:} We are $95\%$ confident that Variety A yields between $0.06$ and $0.74$ tonnes/hectare more than Variety B. Since the interval does \textbf{not} contain $0$, there is evidence at the $5\%$ level that the mean yields differ.
\end{exoresolu}

\courssub{Confidence Interval for a Population Proportion (Frequent GCE Extension)}

Although the syllabus focuses on means, GCE papers often include a question on a proportion. For a large sample, the sample proportion $\hat{p} = X/n$ (where $X$ is the number of successes) is approximately normal:
\[
\hat{p} \sim N\left(p,\; \frac{p(1-p)}{n}\right).
\]
Replacing the unknown $p$ in the standard error by $\hat{p}$ gives the approximate CI.

\begin{propriete}[CI for a Proportion $p$, Large Sample]
If $n\hat{p} \ge 5$ and $n(1-\hat{p}) \ge 5$, an approximate $(1-\alpha)100\%$ confidence interval for $p$ is
\[
\boxed{\hat{p} \;\pm\; z_{\alpha/2}\,\sqrt{\frac{\hat{p}(1-\hat{p})}{n}}}
\]
\end{propriete}

\begin{exoresolu}[CI for a Proportion]
In a survey of $500$ voters in a constituency, $285$ support Candidate X. Find a $90\%$ confidence interval for the true proportion of supporters.
\tcblower
$\hat{p} = 285/500 = 0.57$. $n\hat{p}=285\ge5$, $n(1-\hat{p})=215\ge5$ (conditions satisfied).
$z_{0.05}=1.645$.
$\text{SE} = \sqrt{\frac{0.57\times0.43}{500}} = \sqrt{0.0004902} \approx 0.02214$.
$E = 1.645 \times 0.02214 \approx 0.0364$.
CI: $0.57 \pm 0.0364 \Rightarrow (0.5336,\; 0.6064)$.
\end{exoresolu}

\courssub{Interpreting and Using Confidence Intervals}

\begin{aretenir}[Key Interpretations]
\begin{enumerate}
\item \textbf{Contains a claimed value?} If a hypothesised value $\mu_0$ lies \textbf{inside} a $95\%$ CI, we do \textbf{not} reject the hypothesis $\mu=\mu_0$ at the $5\%$ significance level. If it lies \textbf{outside}, we reject it.
\item \textbf{Comparing two populations:} A CI for $\mu_1-\mu_2$ that does \textbf{not} contain $0$ indicates a statistically significant difference at the corresponding $\alpha$ level.
\item \textbf{Precision:} A narrower interval means a more precise estimate. Width decreases with larger $n$, smaller $\sigma$, or lower confidence level.
\end{enumerate}
\end{aretenir}

\begin{popfigure}
\centering
\begin{tikzpicture}[scale=0.85]
% True mean line
\draw[popdark, thick] (0,0) -- (0,21) node[above, font=\small] {True $\mu$};
% 20 simulated CIs
\foreach \i / \low / \high / \capture in {
  1/ -1.2/ 0.8/ 1,
  2/ -0.5/ 1.5/ 1,
  3/  0.2/ 2.2/ 1,
  4/ -1.8/ 0.2/ 1,
  5/ -0.9/ 1.1/ 1,
  6/  0.5/ 2.5/ 1,
  7/ -2.0/ 0.0/ 0, % misses
  8/ -0.3/ 1.7/ 1,
  9/  0.1/ 2.1/ 1,
  10/ -1.5/ 0.5/ 1,
  11/ -0.7/ 1.3/ 1,
  12/  0.8/ 2.8/ 1,
  13/ -1.0/ 1.0/ 1,
  14/ -0.2/ 1.8/ 1,
  15/  0.3/ 2.3/ 1,
  16/ -1.6/ 0.4/ 1,
  17/ -0.4/ 1.6/ 1,
  18/  0.6/ 2.6/ 1,
  19/ -1.3/ 0.7/ 1,
  20/ -0.8/ 1.2/ 1
}{
  \pgfmathsetmacro{\y}{20-\i+0.5}
  \ifnum\capture=1
    \draw[popgreen, thick] (\low,\y) -- (\high,\y);
    \fill[popgreen] (0,\y) circle (2pt); % true mean marker
  \else
    \draw[poppink, thick] (\low,\y) -- (\high,\y);
    \fill[poppink] (0,\y) circle (2pt);
  \fi
  \node[font=\tiny, anchor=east] at (-2.3,\y) {Sample \i};
}
% Legend
\node[font=\small, anchor=west] at (3,20) {\textcolor{popgreen}{\rule{1em}{0.5em}} Captures $\mu$};
\node[font=\small, anchor=west] at (3,19) {\textcolor{poppink}{\rule{1em}{0.5em}} Misses $\mu$};
\node[font=\small, anchor=west] at (3,18) {19/20 intervals capture $\mu$};
\end{tikzpicture}
\legende{Simulation of twenty $95\%$ confidence intervals for the same true mean $\mu$ (vertical line). In the long run, about $95\%$ (here $19$ out of $20$) capture $\mu$. Any single interval either does or does not — we never know which.}
\end{popfigure}

\courssub{Synthesis: Choosing the Right Formula — A Decision Guide}

GCE questions often mix scenarios. Use this checklist to pick the correct formula.

\begin{methode}[Which Confidence Interval Formula?]
\begin{enumerate}
\item \textbf{How many populations?}
  \begin{itemize}
  \item \textbf{One population} $\rightarrow$ CI for $\mu$ or $p$.
  \item \textbf{Two populations} $\rightarrow$ CI for $\mu_1-\mu_2$ (or $p_1-p_2$, though rare).
  \end{itemize}
\item \textbf{What parameter?}
  \begin{itemize}
  \item Mean $\mu$ $\rightarrow$ use $\bar{x}$ and standard error involving $\sigma$ or $s$.
  \item Proportion $p$ $\rightarrow$ use $\hat{p}$ and $\sqrt{\hat{p}(1-\hat{p})/n}$.
  \end{itemize}
\item \textbf{Is the population variance known?}
  \begin{itemize}
  \item \textbf{Known $\sigma^2$} $\rightarrow$ use $\sigma$ in the formula, critical value $z$.
  \item \textbf{Unknown $\sigma^2$}:
    \begin{itemize}
    \item Large sample ($n\ge30$) $\rightarrow$ use $s$, critical value $z$ (approximate).
    \item Small sample ($n<30$) $\rightarrow$ \textbf{not in syllabus} (would need $t$).
    \end{itemize}
  \end{itemize}
\item \textbf{For two means:} Are both variances known? Both large? Use the corresponding SE formula.
\end{enumerate}
\end{methode}

\begin{exoresolu}[Past-Paper Style Walkthrough]
\textbf{Question:} A manufacturer claims that the mean lifetime of his light bulbs is $\SI{1200}{hours}$ with a standard deviation of $\SI{100}{hours}$. A consumer association tests a random sample of $36$ bulbs and finds a mean lifetime of $\SI{1180}{hours}$.
\begin{enumerate}
\item Construct a $99\%$ confidence interval for the true mean lifetime.
\item Does the interval support the manufacturer's claim? Explain.
\item What sample size would be needed to estimate the mean lifetime to within $\SI{10}{hours}$ with $99\%$ confidence?
\end{enumerate}
\tcblower
\textbf{Solution:}
\begin{enumerate}
\item \textbf{Identify:} One population, mean $\mu$, \textbf{known} $\sigma=100$, $n=36$, $\bar{x}=1180$, $99\% \Rightarrow z_{0.005}=2.576$.
  $\text{SE} = 100/\sqrt{36} = 16.667$.
  $E = 2.576 \times 16.667 \approx 42.93$.
  CI: $1180 \pm 42.93 \Rightarrow (1137.07,\; 1222.93)$.
\item \textbf{Interpretation:} The claimed value $1200$ lies \textbf{inside} the interval $(1137.07, 1222.93)$. Therefore, at the $1\%$ significance level, the sample data do \textbf{not} provide sufficient evidence to reject the manufacturer's claim.
\item \textbf{Sample size:} $E = 10$, $\sigma=100$, $z=2.576$.
  $n = \left(\frac{2.576 \times 100}{10}\right)^2 = (25.76)^2 = 663.58 \Rightarrow n = 664$.
\end{enumerate}
\end{exoresolu}

\begin{exercice}[Practice: Mixed Scenarios]
\begin{enumerate}
\item A random sample of $50$ bags of rice from a market in Bafoussam has mean weight $\SI{24.8}{kg}$ and standard deviation $\SI{0.6}{kg}$. Find a $95\%$ CI for the population mean weight.
\item Two machines fill bottles of juice. Machine A: $n_1=40$, $\bar{x}_1=502$ ml, $s_1=4$ ml. Machine B: $n_2=35$, $\bar{x}_2=498$ ml, $s_2=5$ ml. Construct a $90\%$ CI for the difference in mean fills $\mu_A-\mu_B$. Is there evidence that Machine A fills more on average?
\item In a survey of $400$ households, $132$ have access to pipe-borne water. Find a $99\%$ CI for the true proportion of households with access.
\item A researcher wants a $95\%$ CI for a population mean to have width at most $4$ units. A pilot study gives $s=12$. What minimum sample size is required?
\end{enumerate}
\end{exercice}

\begin{corrige}[Exercise 1]
$\bar{x}=24.8$, $s=0.6$, $n=50\ge30$, $z_{0.025}=1.96$.
$\text{SE}=0.6/\sqrt{50}\approx0.08485$.
$E=1.96\times0.08485\approx0.1663$.
CI: $24.8\pm0.1663 \Rightarrow (24.6337,\;24.9663)$.
\end{corrige}

\begin{corrige}[Exercise 2]
$\bar{x}_A-\bar{x}_B=4$. $\text{SE}=\sqrt{4^2/40+5^2/35}=\sqrt{0.4+0.7143}=\sqrt{1.1143}\approx1.0556$.
$z_{0.05}=1.645$. $E=1.645\times1.0556\approx1.736$.
CI: $4\pm1.736 \Rightarrow (2.264,\;5.736)$.
Since the interval is entirely positive (does not contain $0$), there is evidence at the $10\%$ level that Machine A fills more on average.
\end{corrige}

\begin{corrige}[Exercise 3]
$\hat{p}=132/400=0.33$. $n\hat{p}=132\ge5$, $n(1-\hat{p})=268\ge5$.
$z_{0.005}=2.576$. $\text{SE}=\sqrt{0.33\times0.67/400}=\sqrt{0.00055275}\approx0.02351$.
$E=2.576\times0.02351\approx0.0606$.
CI: $0.33\pm0.0606 \Rightarrow (0.2694,\;0.3906)$.
\end{corrige}

\begin{corrige}[Exercise 4]
Width $W=4 \Rightarrow E=2$. $z_{0.025}=1.96$, $\sigma\approx s=12$.
$n=(1.96\times12/2)^2=(11.76)^2=138.30 \Rightarrow n=139$.
\end{corrige}

\begin{savaistu}[Confidence Intervals in Cameroonian Public Policy]
The National Institute of Statistics (INS) regularly publishes confidence intervals for key indicators: poverty rates, literacy rates, agricultural yields. For instance, Cameroon Demographic and Health Surveys report $95\%$ CIs for the national stunting rate among children under five. Policymakers use the \emph{width} of these intervals to decide whether a new nutrition programme in the East Region has had a measurable impact — if the CI for the difference in stunting rates before/after the programme excludes zero, the programme is deemed effective.
\end{savaistu}

\begin{aretenir}[Summary of Key Formulas]
\begin{center}
\begin{tabular}{|l|l|}\hline
\rowcolor{popblueL}
\textbf{Parameter} & \textbf{$(1-\alpha)100\%$ CI (Large Samples / Known Variance)} \\ \hline
$\mu$ ($\sigma$ known) & $\bar{x} \pm z_{\alpha/2}\,\dfrac{\sigma}{\sqrt{n}}$ \\ \hline
$\mu$ ($\sigma$ unknown, $n\ge30$) & $\bar{x} \pm z_{\alpha/2}\,\dfrac{s}{\sqrt{n}}$ \\ \hline
$\mu_1-\mu_2$ (variances known / large) & $(\bar{x}_1-\bar{x}_2) \pm z_{\alpha/2}\,\sqrt{\dfrac{\sigma_1^2}{n_1}+\dfrac{\sigma_2^2}{n_2}}$ \\ \hline
$p$ (large $n$) & $\hat{p} \pm z_{\alpha/2}\,\sqrt{\dfrac{\hat{p}(1-\hat{p})}{n}}$ \\ \hline
\end{tabular}
\end{center}
\end{aretenir}

This concludes the chapter on Sampling and Estimation. You now have the tools to design samples, derive sampling distributions, compute point estimates (including pooled estimates), and construct and interpret confidence intervals for means, differences of means, and proportions — the core inferential techniques examined at GCE Advanced Level Further Mathematics.

\newpage

\coursec{Integration activity}

\courssub{Aims of the activity}

This integration activity mobilises, in a single authentic study, all the competences developed in the chapter \emph{Sampling and Estimation}. By the end of the activity, each student should be able to:

\begin{itemize}
\item use a \cle{random number table} to select a simple random sample from a finite population, sampling \cle{without replacement};
\item compute a \cle{sample mean} and an \cle{unbiased sample variance} $s^2 = \dfrac{1}{n-1}\sum (x_i - \bar{x})^2$;
\item observe empirically the \cle{sampling distribution of the sample mean} $\bar{X}$, and compare its mean and variance with the theoretical results $E(\bar{X}) = \mu$ and $\mathrm{Var}(\bar{X}) = \dfrac{\sigma^2}{n}\cdot\dfrac{N-n}{N-1}$ (finite population correction);
\item justify a \cle{stratified sampling} plan and contrast it with simple random sampling;
\item produce \cle{pooled estimates} of a population mean and variance from two independent samples;
\item construct and interpret a \cle{95\% confidence interval} for a population mean, and explain the operational meaning of ``95\% confidence'';
\item determine the \cle{sample size} needed to estimate a mean to within a prescribed margin of error.
\end{itemize}

\begin{aretenir}[Competence targeted]
\emph{Integrate sampling techniques, properties of estimators and confidence intervals to design, carry out and report a small-scale statistical study, and to advise a real decision-maker.}
\end{aretenir}

\courssub{Situation}

The \textbf{Mungo Valley Cocoa Cooperative}, based in the Littoral Region, groups together \textbf{600 cocoa farmers}. The cooperative buys fermented cocoa pods by weight and must calibrate its grading scales. To do so, the management needs a reliable estimate of the \textbf{mean weight $\mu$ (in grams) of a mature pod} in the current harvest, together with a measure of the variability of pod weights.

Measuring all pods is impossible: the harvest involves hundreds of thousands of pods, and weighing destroys the pod (it must be broken open). The cooperative therefore hires your class as a team of \textbf{consultant statisticians}.

The cooperative's storekeeper has prepared a \textbf{sampling frame}: a numbered list of $N = 100$ pods, labelled $01$ to $100$, representative of the harvest. The true weight of each pod is known only to the storekeeper, who has sealed the complete list in an \textbf{envelope} kept by the teacher. What is known from previous seasons (long-run records of the cooperative) is summarised below.

\begin{center}
\rowcolors{1}{popblueL}{white}
\begin{tabularx}{\linewidth}{|l|X|}
\hline
\rowcolor{popblueL} \textbf{Known information} & \textbf{Value} \\
\hline
Size of the study population (frame) & $N = 100$ pods \\
\hline
Historical standard deviation of pod weights & $\sigma \approx 20$ g \\
\hline
Farm structure of the cooperative & 420 small farms ($\le 2$ ha), 180 large farms ($> 2$ ha) \\
\hline
True mean $\mu$ and true variance $\sigma^2$ & \emph{hidden in the envelope} \\
\hline
\end{tabularx}
\end{center}

Your class is divided into \textbf{ten working groups}. Each group will draw its own simple random sample of $n = 10$ pods from the frame, using the random number table provided by the teacher (two-digit blocks, reading rows from left to right, ignoring $00$, numbers above $100$ and repetitions, since sampling is \emph{without replacement}). The teacher then reveals, for each selected number only, the weight of the corresponding pod.

For the pooling part of the study, the results of \textbf{Group B} are given in advance: a sample of $n_2 = 12$ pods gave a mean of $\bar{x}_2 = 308$ g and an unbiased sample variance $s_2^2 = 380\ \mathrm{g}^2$.

At the end of the study, each group must hand in a \textbf{written report} containing its calculations, its confidence interval, and a recommendation to the cooperative on the sample size required to estimate the mean pod weight to within $5$ g at the $95\%$ confidence level.

\begin{attention}[Ground rules]
Sampling is done \cle{without replacement}: a pod weighed once cannot be weighed again. No group may communicate its sample to another group before the class pooling phase. All means are in grams; all variances in $\mathrm{g}^2$.
\end{attention}

\courssub{Tasks}

\begin{enumerate}
\item \textbf{Random selection.} Explain why the draw must be done without replacement here. Using the random number table, describe precisely the procedure your group followed to obtain its 10 pod numbers, and state why this procedure guarantees that every sample of 10 pods has the same probability $\dfrac{1}{\binom{100}{10}}$ of being selected.
\item \textbf{Point estimates.} For your sample, compute the sample mean $\bar{x}$ and the \emph{unbiased} sample variance $s^2 = \dfrac{1}{n-1}\sum_{i=1}^{n}(x_i - \bar{x})^2$. Explain why we divide by $n-1$ and not by $n$.
\item \textbf{Sampling distribution of $\bar{X}$.} The teacher collects the ten group means on the board. Compute the mean of the ten means and their variance. After the envelope is opened, compare these empirical values with $\mu$ and $\sigma^2/n$ (with the finite population correction). Which theoretical results does this illustrate? State the Central Limit Theorem and explain why $\bar{X}$ may be treated as approximately normal here.
\item \textbf{Stratification.} The cooperative suspects that large farms, which use more fertiliser, produce heavier pods. Propose a \emph{stratified} sampling plan for a total sample of 50 farmers' pods, with proportional allocation between small and large farms, and justify why this plan is likely to give a more precise estimator than simple random sampling.
\item \textbf{Pooled estimation.} Combine your group's sample (Group A: $n_1 = 10$) with Group B's results ($n_2 = 12$, $\bar{x}_2 = 308$, $s_2^2 = 380$) to obtain pooled estimates of $\mu$ and $\sigma^2$, assuming both samples come from the same population.
\item \textbf{Confidence interval.} Construct a $95\%$ confidence interval for $\mu$ from your group's sample alone, treating $\sigma^2$ as unknown (use the $t$-distribution with $9$ degrees of freedom, $t_{0.025} = 2.262$). Does your interval capture the true value $\mu$ revealed from the envelope?
\item \textbf{Interpretation of ``95\%''.} Out of the ten class intervals, count how many capture $\mu$. Is it a contradiction if one or two intervals fail? Explain what ``$95\%$ confidence'' really means.
\item \textbf{Advice to the cooperative.} Using $\sigma \approx 20$ g, determine the smallest sample size $n$ such that the $95\%$ confidence interval for $\mu$ has total width at most $10$ g (i.e.\ margin of error $5$ g). Write one sentence of advice to the management.
\end{enumerate}

\courssub{Model answer}

We give the full worked answer for \textbf{Group A}, whose random number table produced the pod numbers $07, 23, 41, 12, 88, 35, 59, 04, 77, 30$. The revealed weights (in grams) were:

\begin{center}
\begin{tabular}{|c|cccccccccc|}
\hline
\rowcolor{popblueL} Pod & 07 & 23 & 41 & 12 & 88 & 35 & 59 & 04 & 77 & 30 \\
\hline
Weight (g) & 312 & 298 & 345 & 287 & 330 & 305 & 318 & 290 & 340 & 315 \\
\hline
\end{tabular}
\end{center}

\textbf{1. Random selection.} A pod cannot be broken and weighed twice, so the draw is \cle{without replacement}. Reading two-digit blocks from the table, we keep the first ten distinct numbers between $01$ and $100$. Because each pod has the same chance of being picked at every stage and no pod can be picked twice, all $\binom{100}{10}$ possible samples are equally likely: the sampling is \emph{simple random sampling without replacement}.

\textbf{2. Point estimates.} The sum of the weights is $\sum x_i = 3140$, so
\[ \bar{x} = \frac{3140}{10} = 314\ \text{g}. \]
The squared deviations $(x_i - 314)^2$ are $4, 256, 961, 729, 256, 81, 16, 576, 676, 1$, with sum $\sum (x_i - \bar{x})^2 = 3556$. Hence the unbiased sample variance is
\[ s^2 = \frac{3556}{10 - 1} = \frac{3556}{9} \approx 395.1\ \mathrm{g}^2, \qquad s \approx 19.9\ \text{g}. \]
We divide by $n-1$ because the estimator $S^2 = \dfrac{1}{n-1}\sum (X_i - \bar{X})^2$ satisfies $E(S^2) = \sigma^2$: it is an \cle{unbiased estimator} of the population variance, whereas dividing by $n$ would underestimate $\sigma^2$ on average (the estimator $\frac{1}{n}\sum (X_i - \bar{X})^2$ has expectation $\frac{n-1}{n}\sigma^2$).

\textbf{3. Sampling distribution of $\bar{X}$.} The ten group means collected on the board were
\[ 314,\ 308,\ 321,\ 317,\ 311,\ 319,\ 306,\ 322,\ 313,\ 318. \]
Their mean is $\dfrac{3149}{10} = 314.9$ g and their (empirical) variance is about $29.2\ \mathrm{g}^2$. The envelope is opened: the true population parameters are
\[ \mu = 320\ \text{g}, \qquad \sigma^2 = 400\ \mathrm{g}^2 \quad (\sigma = 20\ \text{g}). \]
For sampling \emph{without replacement} from a finite population of size $N=100$, theory predicts
\[ E(\bar{X}) = \mu = 320, \qquad \mathrm{Var}(\bar{X}) = \frac{\sigma^2}{n}\cdot\frac{N-n}{N-1} = \frac{400}{10}\cdot\frac{90}{99} = 40 \times 0.9091 \approx 36.36\ \mathrm{g}^2. \]
Our empirical values ($314.9$ and $29.2$) are reasonably close: with only ten samples the agreement cannot be exact, but the mean of the means clusters around $\mu$ and the spread of the means is clearly \emph{smaller} than the spread of individual pods ($s \approx 20$ g versus a standard error $\approx \sqrt{36.36} \approx 6.0$ g). This illustrates the two key results: $\bar{X}$ is an unbiased estimator of $\mu$, and its variance is $\sigma^2/n$ multiplied by the finite population correction factor $(N-n)/(N-1)$. By the \cle{Central Limit Theorem}, since pod weights are themselves roughly symmetrically distributed, $\bar{X}$ is approximately $N\!\left(\mu, \dfrac{\sigma^2}{n}\cdot\dfrac{N-n}{N-1}\right)$, which justifies the normal and $t$ procedures used below.

\textbf{4. Stratification.} If pod weight is related to farm size, the population is \emph{heterogeneous} but can be split into two more homogeneous \cle{strata}. With proportional allocation for a sample of 50:
\[ n_{\text{small}} = 50 \times \frac{420}{600} = 35, \qquad n_{\text{large}} = 50 \times \frac{180}{600} = 15. \]
We then draw a simple random sample of 35 pods from small farms and 15 from large farms, and estimate $\mu$ by the weighted mean $\bar{x}_{\text{st}} = 0.7\,\bar{x}_{\text{small}} + 0.3\,\bar{x}_{\text{large}}$. Because variability \emph{within} each stratum is smaller than in the whole population, the stratified estimator has a smaller variance than a simple random sample of the same total size: it is more \cle{efficient}. (The variance of the stratified estimator is $\sum_{h} W_h^2 \frac{\sigma_h^2}{n_h}$ with $W_h = N_h/N$, which is less than $\frac{\sigma^2}{n}$ when strata means differ.)

\textbf{5. Pooled estimation.} Group A: $n_1 = 10$, $\bar{x}_1 = 314$, $s_1^2 = 395.1$. Group B: $n_2 = 12$, $\bar{x}_2 = 308$, $s_2^2 = 380$. The pooled estimate of the mean is
\[ \bar{x}_p = \frac{n_1 \bar{x}_1 + n_2 \bar{x}_2}{n_1 + n_2} = \frac{10(314) + 12(308)}{22} = \frac{6836}{22} \approx 310.7\ \text{g}. \]
The pooled estimate of the common variance (assuming equal population variances) is
\[ s_p^2 = \frac{(n_1 - 1)s_1^2 + (n_2 - 1)s_2^2}{n_1 + n_2 - 2} = \frac{9(395.1) + 11(380)}{20} = \frac{7735.9}{20} \approx 386.8\ \mathrm{g}^2. \]
Both pooled estimators are unbiased, and pooling uses all 22 observations, so they are more precise than either group's estimate alone.

\textbf{6. Confidence interval (variance unknown).} With $n = 10$, $\bar{x} = 314$, $s = 19.9$ and $t_{0.025}(9) = 2.262$, the standard error of the mean (with finite population correction) is
\[ \mathrm{SE} = \frac{s}{\sqrt{n}}\sqrt{\frac{N-n}{N-1}} = \frac{19.9}{\sqrt{10}}\sqrt{\frac{90}{99}} \approx 6.293 \times 0.9535 \approx 6.00\ \text{g}. \]
The $95\%$ confidence interval is
\[ \bar{x} \pm t_{0.025}\,\mathrm{SE} = 314 \pm 2.262 \times 6.00 = 314 \pm 13.6. \]
Thus the interval is $[300.4\ \text{g};\ 327.6\ \text{g}]$. Since $\mu = 320$ g lies inside it, Group A's interval \emph{captures} the true mean. (If the finite population correction were ignored, the interval would be $[299.8;\ 328.2]$, which also contains $\mu$.)

\textbf{7. Meaning of ``95\%''.} Nine of the ten class intervals contained $\mu = 320$ g; one did not. This is \emph{not} a contradiction: ``$95\%$ confidence'' means that the \emph{method} produces intervals capturing $\mu$ in about $95\%$ of all possible samples — roughly 9 or 10 out of 10 here, and about 95 out of 100 in the long run. Any single interval either contains $\mu$ or does not; the $95\%$ describes the reliability of the procedure, not a probability about this particular interval.

\textbf{8. Sample size for a margin of 5 g.} With $\sigma \approx 20$ g and $z_{0.025} = 1.96$, the margin of error for an infinite population (or a very large harvest) is $E = \dfrac{1.96\,\sigma}{\sqrt{n}}$. Requiring $E \le 5$:
\[ n \ge \left(\frac{1.96 \times 20}{5}\right)^2 = (7.84)^2 \approx 61.5. \]
The cooperative should therefore weigh a simple random sample of at least $\boxed{n = 62}$ pods. (If sampling were from the finite frame of $N=100$ pods, the required size would be $n = \frac{62}{1 + 61/100} \approx 39$, but the harvest is much larger than the frame.)

\textbf{Report conclusion (advice to the management).} \emph{``Our study estimates the mean pod weight at about 311--314 g with a standard deviation of about 20 g. To know the mean weight to within 5 g with 95\% confidence, a random sample of 62 pods suffices; if resources allow, a stratified sample (70\% small farms, 30\% large farms) will give the same precision at lower cost.''}

\begin{attention}[Common errors observed during the activity]
\begin{itemize}
\item Dividing by $n$ instead of $n-1$ when estimating $\sigma^2$ from a sample.
\item Using $s$ instead of $s/\sqrt{n}$ (or $s/\sqrt{n}\sqrt{(N-n)/(N-1)}$) in the confidence interval formula.
\item Using $1.96$ when $\sigma$ is unknown and $n$ is small: the $t$-distribution with $n-1$ degrees of freedom is required.
\item Saying ``there is a 95\% probability that $\mu$ is in this interval'': $\mu$ is fixed; it is the \emph{procedure} that succeeds 95\% of the time.
\item Forgetting that sampling from a finite population of pods is done without replacement, and neglecting the finite population correction when appropriate.
\end{itemize}
\end{attention}

\begin{savaistu}[Why this matters in Cameroon]
Cocoa is one of Cameroon's leading export crops, and cooperatives in the Littoral, South-West and Centre regions buy produce by weight and grade. The very same reasoning used in this classroom activity — random sampling, unbiased estimation and confidence intervals — is what professional quality-control statisticians use to certify lots of cocoa, coffee and cotton before export, so that neither the farmer nor the buyer is cheated by a misleading scale or an unrepresentative sample.
\end{savaistu}

\newpage

\coursec{Methods and worked examples}

\coursec{Sampling and Estimation}

\courssub{Sampling from a Finite Population}

\begin{experience}[Discovering Sampling Variability]
A small agricultural cooperative in the West Region has $N=5$ maize plots with yields (in tonnes/hectare): $2.0,\; 2.5,\; 3.0,\; 3.5,\; 4.0$. The manager wants to estimate the average yield by measuring only $n=2$ plots. 
\begin{enumerate}
\item If the two plots are chosen \emph{with replacement}, how many different ordered samples are possible? 
\item If chosen \emph{without replacement}, how many unordered samples are possible?
\item Compute the sample mean for each possible sample in the without-replacement case. Do all samples give the same mean?
\end{enumerate}
\tcblower
\textbf{Observations:}
\begin{itemize}
\item With replacement: $N^n = 5^2 = 25$ ordered samples.
\item Without replacement: $\binom{N}{n} = \binom{5}{2} = 10$ unordered samples.
\item The sample means vary: $2.25,\; 2.50,\; 2.75,\; 3.00,\; 3.25,\; 3.50,\; 3.75$. This variability is the essence of \cle{sampling distributions}.
\end{itemize}
\end{experience}

\begin{definition}[Population and Sample]
A \cle{finite population} consists of $N$ identifiable units $U_1,\dots,U_N$ with associated values $x_1,\dots,x_N$ of a variable $X$. The \cle{population mean} and \cle{population variance} are
\[
\mu = \frac{1}{N}\sum_{i=1}^N x_i,\qquad 
\sigma^2 = \frac{1}{N}\sum_{i=1}^N (x_i-\mu)^2.
\]
A \cle{sample} of size $n$ is a subset of $n$ units. A \cle{statistic} is any function of the sample values (e.g., sample mean $\bar{X}$, sample variance $S^2$). The \cle{sampling distribution} of a statistic is the probability distribution of that statistic over all possible samples of a given size drawn according to a specified sampling scheme.
\end{definition}

\begin{propriete}[Sampling Distribution of the Sample Mean -- Finite Population]
Let $\bar{X}$ be the sample mean from a sample of size $n$ drawn from a finite population of size $N$ with mean $\mu$ and variance $\sigma^2$.
\begin{itemize}
\item \textbf{With replacement} (or from an infinite population): 
\[
E(\bar{X}) = \mu,\qquad \mathrm{Var}(\bar{X}) = \frac{\sigma^2}{n}.
\]
\item \textbf{Without replacement}: 
\[
E(\bar{X}) = \mu,\qquad \mathrm{Var}(\bar{X}) = \frac{\sigma^2}{n}\left(1-\frac{n-1}{N-1}\right) = \frac{\sigma^2}{n}\left(\frac{N-n}{N-1}\right).
\]
The factor $\left(1-\frac{n-1}{N-1}\right)$ is the \cle{finite population correction (fpc)}. It reduces the variance because sampling without replacement gives more information per unit.
\end{itemize}
\textbf{Proof (examinable):} For without replacement, use the fact that $\mathrm{Cov}(X_i,X_j) = -\frac{\sigma^2}{N-1}$ for $i\neq j$ in a simple random sample. Then $\mathrm{Var}(\bar{X}) = \frac{1}{n^2}\left[n\sigma^2 + n(n-1)\left(-\frac{\sigma^2}{N-1}\right)\right] = \frac{\sigma^2}{n}\left(1-\frac{n-1}{N-1}\right)$.
\end{propriete}

\begin{methode}[Constructing the Sampling Distribution of $\bar{X}$ from a Finite Population]
\textbf{Objective:} Obtain the exact probability distribution of the sample mean for small $N$ and $n$.

\textbf{Steps:}
\begin{enumerate}
\item List the $N$ population values. Compute $\mu$ and $\sigma^2$.
\item Determine the sampling scheme (with/without replacement) and list all possible samples (ordered or unordered). Each sample is equally likely.
\item For each sample, compute the sample mean $\bar{x}$.
\item Tabulate the distinct values of $\bar{x}$ with their probabilities (frequency / total number of samples).
\item Verify $E(\bar{X})=\mu$ and $\mathrm{Var}(\bar{X})$ matches the formula.
\end{enumerate}
\end{methode}

\begin{exoresolu}[Sampling Distribution of the Mean -- Maize Yields]
A farm has $5$ plots with yields (t/ha): $2.0,\; 2.5,\; 3.0,\; 3.5,\; 4.0$.
\begin{enumerate}
\item Compute $\mu$ and $\sigma^2$.
\item List all samples of size $n=2$ \textbf{without replacement} and their means.
\item Construct the sampling distribution of $\bar{X}$ and verify $E(\bar{X})=\mu$ and $\mathrm{Var}(\bar{X}) = \frac{\sigma^2}{n}\left(1-\frac{n-1}{N-1}\right)$.
\end{enumerate}
\tcblower
\textbf{Detailed Solution}

\textbf{1. Population parameters:}
$N=5$, values: $2.0, 2.5, 3.0, 3.5, 4.0$.
\[
\mu = \frac{2.0+2.5+3.0+3.5+4.0}{5} = 3.0.
\]
\[
\sigma^2 = \frac{1}{5}\big[(2.0-3.0)^2+(2.5-3.0)^2+(3.0-3.0)^2+(3.5-3.0)^2+(4.0-3.0)^2\big]
= \frac{1}{5}[1.0+0.25+0+0.25+1.0] = 0.5.
\]

\textbf{2. All samples of size $2$ without replacement:}
$\binom{5}{2}=10$ equally likely samples (probability $1/10$ each):

\begin{center}
\begin{tabular}{|c|c|c|}
\hline
Sample & Values & $\bar{x}$ \\
\hline
1 & (2.0, 2.5) & 2.25 \\
2 & (2.0, 3.0) & 2.50 \\
3 & (2.0, 3.5) & 2.75 \\
4 & (2.0, 4.0) & 3.00 \\
5 & (2.5, 3.0) & 2.75 \\
6 & (2.5, 3.5) & 3.00 \\
7 & (2.5, 4.0) & 3.25 \\
8 & (3.0, 3.5) & 3.25 \\
9 & (3.0, 4.0) & 3.50 \\
10 & (3.5, 4.0) & 3.75 \\
\hline
\end{tabular}
\end{center}

\textbf{3. Sampling distribution of $\bar{X}$:}

\begin{center}
\begin{tabular}{|c|c|c|c|c|c|c|c|}
\hline
$\bar{x}$ & 2.25 & 2.50 & 2.75 & 3.00 & 3.25 & 3.50 & 3.75 \\
\hline
$P(\bar{X}=\bar{x})$ & 0.1 & 0.1 & 0.2 & 0.2 & 0.2 & 0.1 & 0.1 \\
\hline
\end{tabular}
\end{center}

\textbf{Check $E(\bar{X})$:}
\[
E(\bar{X}) = \sum \bar{x}\,P(\bar{X}=\bar{x}) = 0.1(2.25+2.50+2.75+3.00+3.25+3.50+3.75) + 0.1(2.75+3.00+3.25)
\]
Better: average of the $10$ sample means:
\[
E(\bar{X}) = \frac{1}{10}(2.25+2.50+2.75+3.00+2.75+3.00+3.25+3.25+3.50+3.75) = \frac{30.0}{10} = 3.0 = \mu.
\]

\textbf{Check $\mathrm{Var}(\bar{X})$:}
\[
E(\bar{X}^2) = \frac{1}{10}\sum \bar{x}_i^2 = \frac{1}{10}(5.0625+6.25+7.5625+9.0+7.5625+9.0+10.5625+10.5625+12.25+14.0625) = \frac{91.875}{10} = 9.1875.
\]
\[
\mathrm{Var}(\bar{X}) = E(\bar{X}^2) - [E(\bar{X})]^2 = 9.1875 - 9.0 = 0.1875.
\]

\textbf{Theoretical variance with fpc:}
\[
\frac{\sigma^2}{n}\left(1-\frac{n-1}{N-1}\right) = \frac{0.5}{2}\left(1-\frac{1}{4}\right) = 0.25 \times 0.75 = 0.1875.
\]
Matches exactly. \qed
\end{exoresolu}

\begin{popfigure}
\begin{tikzpicture}
\begin{axis}[
    width=0.9\linewidth,
    height=6cm,
    ybar=2pt,
    bar width=8pt,
    xlabel={Sample mean $\bar{x}$ (t/ha)},
    ylabel={Probability},
    xtick={2.25,2.50,2.75,3.00,3.25,3.50,3.75},
    xticklabel style={rotate=45, anchor=east},
    ymin=0, ymax=0.25,
    ymajorgrids=true,
    grid style=dashed,
]
\addplot[fill=popblue, draw=popdark] coordinates {
    (2.25,0.1) (2.50,0.1) (2.75,0.2) (3.00,0.2) (3.25,0.2) (3.50,0.1) (3.75,0.1)
};
\addplot[mark=*, only marks, red, thick] coordinates {(3.0,0)};
\end{axis}
\end{tikzpicture}
\legende{Sampling distribution of $\bar{X}$ for $n=2$ without replacement from the maize population. The red dot marks $\mu=3.0$.}
\end{popfigure}

\begin{attention}[Common Mistakes in Finite Population Sampling]
\begin{itemize}
\item Forgetting the finite population correction (fpc) when sampling without replacement. The fpc is $\frac{N-n}{N-1}$, \emph{not} $\frac{N-n}{N}$ (which is used for proportions in some textbooks; for means, the correct denominator is $N-1$).
\item Confusing ordered vs. unordered samples. For without replacement, use combinations $\binom{N}{n}$; for with replacement, use $N^n$.
\item Using the sample variance formula $s^2 = \frac{1}{n-1}\sum(x_i-\bar{x})^2$ to estimate $\sigma^2$ when the population is finite and the sample is a large fraction of it -- the fpc must be applied to the variance of the estimator as well.
\end{itemize}
\end{attention}

\courssub{Random Sampling and Sampling of Attributes}

\begin{definition}[Simple Random Sample and Random Number Tables]
A \cle{simple random sample (SRS)} of size $n$ from a population of size $N$ is a sample where every subset of $n$ units has the same probability $\frac{1}{\binom{N}{n}}$ of being selected. In practice, we use \cle{random number tables} or computer generators. Each unit is assigned a unique label from $1$ to $N$ (padding with leading zeros to match the number of digits of $N$). We then read the table in groups of that many digits, accepting numbers between $1$ and $N$ that have not appeared before (for without replacement).
\end{definition}

\begin{definition}[Sampling of Attributes]
When each unit either possesses a certain attribute (success) or not (failure), we have a population of \cle{attributes}. Let $p$ be the population proportion of successes. In a sample of size $n$, let $X$ be the number of successes. The \cle{sample proportion} is $\hat{p} = X/n$.
\begin{itemize}
\item For an infinite population or sampling with replacement: $X \sim \mathrm{Binomial}(n,p)$, so $E(\hat{p})=p$, $\mathrm{Var}(\hat{p}) = \frac{p(1-p)}{n}$.
\item For a finite population of size $N$ with $M$ successes ($p=M/N$), sampling without replacement: $X \sim \mathrm{Hypergeometric}(N,M,n)$, so $E(\hat{p})=p$, $\mathrm{Var}(\hat{p}) = \frac{p(1-p)}{n}\left(1-\frac{n-1}{N-1}\right)$.
\end{itemize}
For large $n$, $\hat{p}$ is approximately normal by the CLT.
\end{definition}

\begin{methode}[Selecting a Simple Random Sample Using Random Number Tables]
\textbf{Steps:}
\begin{enumerate}
\item Assign numbers $1$ to $N$ to each population unit (use leading zeros: e.g., $001$ to $200$ for $N=200$).
\item Choose a random starting point in the table. Read digits in groups of $d$ digits where $10^{d-1} \le N < 10^d$.
\item Accept a number if it is between $1$ and $N$ inclusive and has not been selected before (for without replacement). Skip numbers $>N$ or $=0$ (or $000$) and duplicates.
\item Continue until $n$ distinct units are selected.
\end{enumerate}
\end{methode}

\begin{exoresolu}[Estimating a Proportion Using Random Numbers]
A village has $200$ farmers (numbered $001$ to $200$). An SRS of $20$ farmers is selected using the random number table below (read row-wise in groups of three digits). It is found that $14$ of the $20$ sampled farmers use improved maize seeds.

\begin{center}
\begin{tabular}{cccccccccc}
842 & 103 & 197 & 025 & 118 & 042 & 201 & 156 & 089 & 134 \\
067 & 190 & 012 & 178 & 055 & 111 & 200 & 003 & 076 & 143 \\
091 & 162 & 034 & 187 & 001 & 123 & 045 & 068 & 099 & 110 \\
\end{tabular}
\end{center}

\begin{enumerate}
\item Identify the $20$ farmers selected.
\item Compute the point estimate $\hat{p}$ of the proportion using improved seeds.
\item Calculate the estimated standard error of $\hat{p}$ using the finite population correction.
\end{enumerate}
\tcblower
\textbf{Detailed Solution}

\textbf{1. Selection of 20 farmers:}
Read 3-digit groups sequentially, accepting numbers $001$--$200$, ignoring duplicates and numbers $>200$ or $000$.

\begin{center}
\begin{tabular}{|c|c|c|}
\hline
Order & Random number & Accepted? \\
\hline
1 & 842 & No ($>200$) \\
2 & 103 & \textbf{Yes} (\#1) \\
3 & 197 & \textbf{Yes} (\#2) \\
4 & 025 & \textbf{Yes} (\#3) \\
5 & 118 & \textbf{Yes} (\#4) \\
6 & 042 & \textbf{Yes} (\#5) \\
7 & 201 & No ($>200$) \\
8 & 156 & \textbf{Yes} (\#6) \\
9 & 089 & \textbf{Yes} (\#7) \\
10 & 134 & \textbf{Yes} (\#8) \\
11 & 067 & \textbf{Yes} (\#9) \\
12 & 190 & \textbf{Yes} (\#10) \\
13 & 012 & \textbf{Yes} (\#11) \\
14 & 178 & \textbf{Yes} (\#12) \\
15 & 055 & \textbf{Yes} (\#13) \\
16 & 111 & \textbf{Yes} (\#14) \\
17 & 200 & \textbf{Yes} (\#15) \\
18 & 003 & \textbf{Yes} (\#16) \\
19 & 076 & \textbf{Yes} (\#17) \\
20 & 143 & \textbf{Yes} (\#18) \\
21 & 091 & \textbf{Yes} (\#19) \\
22 & 162 & \textbf{Yes} (\#20) \\
\hline
\end{tabular}
\end{center}
The sample consists of farmers: $103, 197, 025, 118, 042, 156, 089, 134, 067, 190, 012, 178, 055, 111, 200, 003, 076, 143, 091, 162$.

\textbf{2. Point estimate $\hat{p}$:}
$n=20$, $X=14$ successes.
\[
\hat{p} = \frac{X}{n} = \frac{14}{20} = 0.7.
\]

\textbf{3. Estimated standard error with fpc:}
$N=200$, $n=20$, $\hat{p}=0.7$.
\[
\widehat{\mathrm{Var}}(\hat{p}) = \frac{\hat{p}(1-\hat{p})}{n}\left(1-\frac{n-1}{N-1}\right)
= \frac{0.7\times0.3}{20}\left(1-\frac{19}{199}\right)
= 0.0105 \times 0.904523 = 0.0094975.
\]
\[
\widehat{\mathrm{SE}}(\hat{p}) = \sqrt{0.0094975} \approx 0.09746.
\]
(Without fpc: $\mathrm{SE} = \sqrt{0.21/20} \approx 0.1025$; the fpc reduces the standard error by about $5\%$.) \qed
\end{exoresolu}

\begin{savaistu}[Random Number Tables in Cameroon]
Historically, GCE examinations provided printed random number tables (e.g., Fisher \& Yates, or Tippett tables). Today, candidates are expected to know the \emph{procedure} for using such tables; actual random numbers may be provided in the question. In modern practice, statistical software (R, Python, Excel) or calculator functions (\texttt{Ran\#}, \texttt{RanInt}) replace printed tables, but the principle of simple random sampling remains identical.
\end{savaistu}

\courssub{Sampling Methods}

\begin{definition}[Probability vs. Non-Probability Sampling]
\begin{itemize}
\item \cle{Probability sampling}: Every unit has a known, non-zero probability of selection. Allows statistical inference (confidence intervals, hypothesis tests). Examples: simple random, systematic, stratified, cluster, multi-stage.
\item \cle{Non-probability sampling}: Selection based on convenience, judgment, or quotas. Probabilities unknown. Inference is not statistically valid. Examples: quota, purposive, convenience, snowball.
\end{itemize}
The GCE syllabus focuses on probability sampling methods.
\end{definition}

\begin{propriete}[Main Probability Sampling Methods]
\begin{enumerate}
\item \textbf{Simple Random Sampling (SRS):} Each of the $\binom{N}{n}$ samples equally likely. Can be with or without replacement (usually without). Requires a complete sampling frame.
\item \textbf{Systematic Sampling:} Choose a random start $r$ between $1$ and $k$, then select every $k$-th unit where $k = \lfloor N/n \rfloor$. Efficient if the list is randomly ordered; biased if there is periodicity.
\item \textbf{Stratified Sampling:} Divide population into $H$ homogeneous strata of sizes $N_h$. Draw an SRS of size $n_h$ from each stratum (proportional allocation: $n_h = n\frac{N_h}{N}$; or optimal/Neyman allocation). Reduces variance if strata are internally homogeneous but differ from each other.
\item \textbf{Cluster Sampling:} Divide population into $M$ clusters (e.g., villages, schools). Select a random sample of $m$ clusters, then survey \emph{all} units in selected clusters (one-stage) or a subsample (two-stage). Cost-effective for geographically dispersed populations; increases variance due to intra-cluster correlation.
\item \textbf{Multi-stage Sampling:} Combination of the above (e.g., select regions, then districts within regions, then households within districts). Standard for large national surveys (e.g., Cameroon DHS, MICS).
\end{enumerate}
\end{propriete}

\begin{methode}[Stratified Sampling -- Estimation of Mean and Variance]
\textbf{Notation:} $H$ strata, stratum $h$ has size $N_h$, sample size $n_h$, sample mean $\bar{x}_h$, sample variance $s_h^2$. $N = \sum N_h$, $n = \sum n_h$.
\begin{itemize}
\item \textbf{Estimator of population mean:} $\bar{x}_{\text{st}} = \sum_{h=1}^H \frac{N_h}{N} \bar{x}_h$ (unbiased).
\item \textbf{Variance of $\bar{x}_{\text{st}}$ (without replacement within strata):}
\[
\mathrm{Var}(\bar{X}_{\text{st}}) = \sum_{h=1}^H \left(\frac{N_h}{N}\right)^2 \frac{s_h^2}{n_h}\left(1-\frac{n_h}{N_h}\right).
\]
\item \textbf{Proportional allocation:} $n_h = n \frac{N_h}{N}$. Then $\bar{x}_{\text{st}} = \bar{x}$ (overall sample mean) and variance simplifies.
\item \textbf{Optimal (Neyman) allocation:} $n_h \propto N_h s_h$ (minimizes variance for fixed cost).
\end{itemize}
\end{methode}

\begin{exoresolu}[Stratified Sampling -- Cocoa Yields]
A survey estimates mean cocoa yield in a region with two strata: Smallholdings ($N_1=800$) and Plantations ($N_2=200$). A stratified sample with proportional allocation gives:
Stratum 1: $n_1=80$, $\bar{x}_1=1.2$ t/ha, $s_1^2=0.16$.
Stratum 2: $n_2=20$, $\bar{x}_2=2.5$ t/ha, $s_2^2=0.36$.
\begin{enumerate}
\item Estimate the overall mean yield.
\item Estimate the variance of this estimator (ignore fpc for simplicity, or include it).
\end{enumerate}
\tcblower
\textbf{Solution:}
$N=1000$, $n=100$. Proportional allocation: $n_1=80$, $n_2=20$ (matches given).
\textbf{1.} $\bar{x}_{\text{st}} = \frac{800}{1000}\times1.2 + \frac{200}{1000}\times2.5 = 0.96 + 0.50 = 1.46$ t/ha.
\textbf{2.} With fpc:
$\mathrm{Var} = (0.8)^2 \frac{0.16}{80}(1-\frac{80}{800}) + (0.2)^2 \frac{0.36}{20}(1-\frac{20}{200})
= 0.64 \times 0.002 \times 0.9 + 0.04 \times 0.018 \times 0.9
= 0.001152 + 0.000648 = 0.0018$.
$\mathrm{SE} \approx 0.0424$ t/ha. \qed
\end{exoresolu}

\begin{attention}[Sampling Methods -- Exam Traps]
\begin{itemize}
\item \textbf{Systematic sampling} is \emph{not} equivalent to SRS unless the list is in random order. If the list has a cyclic pattern matching the sampling interval $k$, estimates can be severely biased.
\item \textbf{Cluster sampling} increases variance compared to SRS of the same size because units within a cluster tend to be similar (positive intra-cluster correlation). The \cle{design effect} (deff) quantifies this: $\mathrm{deff} = 1 + (b-1)\rho$ where $b$ is cluster size and $\rho$ is intra-cluster correlation.
\item \textbf{Stratified sampling} always reduces variance (or leaves it unchanged) compared to SRS of the same total size, provided strata are properly defined.
\item In exams, if asked to "describe how to select a systematic sample", you must mention: (1) determine $k = N/n$, (2) choose random start $r \in \{1,\dots,k\}$, (3) select units $r, r+k, r+2k, \dots$.
\end{itemize}
\end{attention}

\courssub{The Sample Mean and the Central Limit Theorem}

\begin{propriete}[Central Limit Theorem (CLT) -- Examinable Statement]
Let $X_1, X_2, \dots, X_n$ be a random sample (i.i.d.) from a population with mean $\mu$ and finite variance $\sigma^2$. Let $\bar{X} = \frac{1}{n}\sum_{i=1}^n X_i$. Then as $n \to \infty$, the distribution of the standardized sample mean
\[
Z = \frac{\bar{X} - \mu}{\sigma/\sqrt{n}}
\]
converges to the standard normal distribution $N(0,1)$. Equivalently, for large $n$,
\[
\bar{X} \stackrel{\text{approx}}{\sim} N\left(\mu,\; \frac{\sigma^2}{n}\right).
\]
\textbf{Practical rule:} $n \ge 30$ is usually sufficient for the approximation to be good, unless the population distribution is extremely skewed or has heavy tails. If the population is normal, $\bar{X}$ is exactly normal for any $n$.
\end{propriete}

\begin{popfigure}
\begin{tikzpicture}
\begin{axis}[
    width=0.9\linewidth,
    height=7cm,
    xlabel={$x$},
    ylabel={Density},
    legend pos=north east,
    domain=-4:4,
    samples=200,
    ymin=0, ymax=0.8,
]
% Population: exponential-like (skewed)
\addplot[thick, poporange, domain=0:4] {exp(-x)};
\addlegendentry{Population (skewed)}

% Sampling distribution of mean for n=1 (same as population)
% For n=5: approx normal with sd = 1/sqrt(5)
\addplot[thick, popblue, domain=-1:3] {1/sqrt(max(0,2*pi*0.2))*exp(-(x-1)^2/(2*0.2))};
\addlegendentry{$\overline{X}$ for $n=5$ (approx)}

% For n=30: sd = 1/sqrt(30) ≈ 0.183
\addplot[thick, popgreen, domain=0.5:1.5] {1/sqrt(max(0,2*pi*0.0333))*exp(-(x-1)^2/(2*0.0333))};
\addlegendentry{$\overline{X}$ for $n=30$ (approx)}

% Normal limit
\addplot[dashed, thick, popdark, domain=0.5:1.5] {1/sqrt(max(0,2*pi*0.01))*exp(-(x-1)^2/(2*0.01))};
\addlegendentry{Limit $N(1{,} 1/n)$ as $n\rightarrow \infty$}
\end{axis}
\end{tikzpicture}
\legende{Illustration of the CLT: as sample size $n$ increases, the sampling distribution of $\bar{X}$ becomes more concentrated around $\mu$ and approaches a normal shape, even if the population is skewed.}
\end{popfigure}

\begin{methode}[Using the CLT for Probabilities on $\bar{X}$]
\textbf{Steps:}
\begin{enumerate}
\item Identify $\mu$, $\sigma$ (or $s$ if $\sigma$ unknown and $n$ large), and $n$.
\item Check: $n \ge 30$ (or population normal). If finite population and $n/N > 0.05$, include fpc.
\item Write the approximate distribution: $\bar{X} \sim N(\mu, \sigma^2/n)$ (or with fpc).
\item Standardize: $Z = \frac{\bar{X}-\mu}{\sigma/\sqrt{n}}$ (or with fpc).
\item Use standard normal tables to find required probabilities.
\end{enumerate}
\end{methode}

\begin{exoresolu}[CLT Application -- Cement Bag Weights]
Cement bags from a Douala factory have mean weight $\mu = 50.0$ kg and standard deviation $\sigma = 0.5$ kg. Distribution shape unknown. A random sample of $n=50$ bags is weighed.
\begin{enumerate}
\item State the approximate distribution of $\bar{X}$.
\item Find $P(\bar{X} < 49.9)$.
\item Find $P(49.95 < \bar{X} < 50.05)$.
\end{enumerate}
\tcblower
\textbf{Detailed Solution}

\textbf{1.} $n=50 \ge 30$, so by CLT:
\[
\bar{X} \sim N\left(50.0,\; \frac{0.5^2}{50}\right) = N(50.0,\; 0.005).
\]
$\sigma_{\bar{X}} = 0.5/\sqrt{50} \approx 0.07071$ kg.

\textbf{2.} $P(\bar{X} < 49.9)$:
\[
Z = \frac{49.9 - 50.0}{0.07071} = -1.4142.
\]
$P(Z < -1.4142) = 1 - P(Z < 1.4142) \approx 1 - 0.9214 = 0.0786$.
(Using tables: $P(Z < -1.41) = 0.0793$, $P(Z < -1.42) = 0.0778$; interpolation $\approx 0.0786$.)

\textbf{3.} $P(49.95 < \bar{X} < 50.05)$:
\[
Z_1 = \frac{49.95-50.0}{0.07071} = -0.7071,\quad Z_2 = \frac{50.05-50.0}{0.07071} = 0.7071.
\]
$P(-0.7071 < Z < 0.7071) = 2P(Z < 0.7071) - 1$.
$P(Z < 0.71) \approx 0.7611$, $P(Z < 0.70) \approx 0.7580$. Interpolating: $P(Z < 0.7071) \approx 0.7600$.
Probability $\approx 2\times0.7600 - 1 = 0.5200$. \qed
\end{exoresolu}

\begin{aretenir}[Key Points on CLT]
\begin{itemize}
\item The CLT applies to the \emph{sample mean} (or sum), not to individual observations.
\item The approximation improves as $n$ increases. For highly skewed populations (e.g., exponential, lognormal), $n \ge 50$ or more may be needed.
\item For finite populations without replacement, replace $\sigma^2/n$ by $\frac{\sigma^2}{n}\left(1-\frac{n-1}{N-1}\right)$ in the standard error.
\item If $\sigma$ is unknown and $n \ge 30$, we can use the sample standard deviation $s$ in place of $\sigma$ (the error is negligible).
\end{itemize}
\end{aretenir}

\courssub{Point Estimation: Unbiased and Most Efficient Estimators}

\begin{definition}[Estimator, Estimate, Bias, MSE]
Let $\theta$ be a population parameter. A statistic $\hat{\Theta} = g(X_1,\dots,X_n)$ used to estimate $\theta$ is called an \cle{estimator}; its observed value is an \cle{estimate}.
\begin{itemize}
\item \textbf{Bias:} $\mathrm{Bias}(\hat{\Theta}) = E(\hat{\Theta}) - \theta$. If $\mathrm{Bias}=0$, $\hat{\Theta}$ is \cle{unbiased}.
\item \textbf{Mean Squared Error (MSE):} $\mathrm{MSE}(\hat{\Theta}) = E[(\hat{\Theta}-\theta)^2] = \mathrm{Var}(\hat{\Theta}) + [\mathrm{Bias}(\hat{\Theta})]^2$.
\item \textbf{Efficiency:} For two unbiased estimators $\hat{\Theta}_1, \hat{\Theta}_2$ of $\theta$, $\hat{\Theta}_1$ is more \cle{efficient} than $\hat{\Theta}_2$ if $\mathrm{Var}(\hat{\Theta}_1) < \mathrm{Var}(\hat{\Theta}_2)$. The \cle{relative efficiency} is $\frac{\mathrm{Var}(\hat{\Theta}_2)}{\mathrm{Var}(\hat{\Theta}_1)}$.
\item \textbf{MVUE:} The \cle{Minimum Variance Unbiased Estimator} has the smallest variance among all unbiased estimators.
\end{itemize}
\end{definition}

\begin{propriete}[Standard Unbiased Estimators]
For a random sample $X_1,\dots,X_n$ from a population with mean $\mu$ and variance $\sigma^2$:
\begin{itemize}
\item $\bar{X}$ is unbiased for $\mu$: $E(\bar{X}) = \mu$.
\item $S^2 = \frac{1}{n-1}\sum_{i=1}^n (X_i-\bar{X})^2$ is unbiased for $\sigma^2$: $E(S^2) = \sigma^2$.
\item $\hat{p} = X/n$ is unbiased for $p$ (binomial proportion).
\end{itemize}
\textbf{Proof that $S^2$ is unbiased (examinable):}
\[
\sum (X_i-\bar{X})^2 = \sum (X_i-\mu)^2 - n(\bar{X}-\mu)^2.
\]
Take expectations: $E[\sum (X_i-\mu)^2] = n\sigma^2$, $E[n(\bar{X}-\mu)^2] = n\frac{\sigma^2}{n} = \sigma^2$.
So $E[\sum (X_i-\bar{X})^2] = n\sigma^2 - \sigma^2 = (n-1)\sigma^2$, hence $E(S^2) = \sigma^2$.
\end{propriete}

\begin{methode}[Comparing Estimators for Efficiency]
\textbf{Steps:}
\begin{enumerate}
\item Compute $E(\hat{\Theta}_1)$ and $E(\hat{\Theta}_2)$. Check unbiasedness.
\item If both unbiased, compute $\mathrm{Var}(\hat{\Theta}_1)$ and $\mathrm{Var}(\hat{\Theta}_2)$ using independence and $\mathrm{Var}(aX) = a^2\mathrm{Var}(X)$.
\item The one with smaller variance is more efficient.
\item Relative efficiency of $\hat{\Theta}_2$ w.r.t. $\hat{\Theta}_1$ = $\frac{\mathrm{Var}(\hat{\Theta}_1)}{\mathrm{Var}(\hat{\Theta}_2)}$.
\end{enumerate}
\end{methode}

\begin{exoresolu}[Comparing Two Estimators of $\mu$]
Let $X_1,X_2,X_3$ be a random sample from a population with mean $\mu$, variance $\sigma^2$. Consider:
\[
\hat{\mu}_1 = \frac{X_1+X_2+X_3}{3},\qquad
\hat{\mu}_2 = \frac{X_1+2X_2+X_3}{4}.
\]
\begin{enumerate}
\item Show both are unbiased for $\mu$.
\item Find $\mathrm{Var}(\hat{\mu}_1)$ and $\mathrm{Var}(\hat{\mu}_2)$.
\item Which is more efficient? Compute relative efficiency of $\hat{\mu}_2$ w.r.t. $\hat{\mu}_1$.
\end{enumerate}
\tcblower
\textbf{Detailed Solution}

\textbf{1. Unbiasedness:}
$E(\hat{\mu}_1) = \frac{\mu+\mu+\mu}{3} = \mu$.
$E(\hat{\mu}_2) = \frac{\mu+2\mu+\mu}{4} = \mu$.
Both unbiased.

\textbf{2. Variances (using independence):}
$\mathrm{Var}(\hat{\mu}_1) = \frac{1}{9}(\sigma^2+\sigma^2+\sigma^2) = \frac{\sigma^2}{3}$.
$\mathrm{Var}(\hat{\mu}_2) = \frac{1}{16}(\sigma^2 + 4\sigma^2 + \sigma^2) = \frac{6\sigma^2}{16} = \frac{3\sigma^2}{8}$.

\textbf{3. Efficiency:}
$\frac{\sigma^2}{3} \approx 0.333\sigma^2 < \frac{3\sigma^2}{8} = 0.375\sigma^2$, so $\hat{\mu}_1$ is more efficient.
Relative efficiency of $\hat{\mu}_2$ w.r.t. $\hat{\mu}_1$:
\[
\frac{\mathrm{Var}(\hat{\mu}_1)}{\mathrm{Var}(\hat{\mu}_2)} = \frac{\sigma^2/3}{3\sigma^2/8} = \frac{8}{9} \approx 0.889.
\]
$\hat{\mu}_2$ has $88.9\%$ the efficiency of $\hat{\mu}_1$. \qed
\end{exoresolu}

\begin{savaistu}[Why Divide by $n-1$?]
The sample variance $S^2 = \frac{1}{n-1}\sum(X_i-\bar{X})^2$ uses $n-1$ (degrees of freedom) because $\bar{X}$ is estimated from the same data, causing $\sum(X_i-\bar{X})^2$ to underestimate $\sum(X_i-\mu)^2$ on average. The factor $n-1$ corrects this bias. If we used $n$, the estimator would be biased (though asymptotically unbiased).
\end{savaistu}

\courssub{Pooled Estimation from Two Samples}

\begin{definition}[Pooled Estimation]
When two independent samples come from populations with a \emph{common} parameter (mean or variance), we can combine (pool) them to get a more precise estimate.
\begin{itemize}
\item \textbf{Pooled mean} (common $\mu$, possibly different variances): $\bar{x}_p = \frac{n_1\bar{x}_1 + n_2\bar{x}_2}{n_1+n_2}$.
\item \textbf{Pooled variance} (common $\sigma^2$, normal populations): $s_p^2 = \frac{(n_1-1)s_1^2 + (n_2-1)s_2^2}{n_1+n_2-2}$.
\end{itemize}
\end{definition}

\begin{propriete}[Properties of Pooled Estimators]
\begin{itemize}
\item $\bar{X}_p$ is unbiased for $\mu$: $E(\bar{X}_p) = \mu$.
\item $\mathrm{Var}(\bar{X}_p) = \frac{n_1\sigma_1^2 + n_2\sigma_2^2}{(n_1+n_2)^2}$. If $\sigma_1^2=\sigma_2^2=\sigma^2$, then $\mathrm{Var}(\bar{X}_p) = \frac{\sigma^2}{n_1+n_2}$.
\item $S_p^2$ is unbiased for $\sigma^2$: $E(S_p^2) = \sigma^2$, with $n_1+n_2-2$ degrees of freedom.
\item If both populations are normal with common $\sigma^2$, then $\frac{(n_1+n_2-2)S_p^2}{\sigma^2} \sim \chi^2_{n_1+n_2-2}$.
\end{itemize}
\end{propriete}

\begin{exoresolu}[Pooled Estimation of Mean and Variance]
Two independent samples from normal populations with common mean $\mu$ but possibly different variances:
Sample A: $n_1=10$, $\bar{x}_1=25.4$, $s_1^2=4.2$.
Sample B: $n_2=15$, $\bar{x}_2=26.1$, $s_2^2=5.8$.
\begin{enumerate}
\item Pooled estimate of $\mu$.
\item Assuming common variance $\sigma^2$, pooled estimate $s_p^2$.
\item If true $\sigma^2=5.0$, variance of $\bar{X}_p$?
\end{enumerate}
\tcblower
\textbf{Solution:}

\textbf{1.} $\bar{x}_p = \frac{10\times25.4 + 15\times26.1}{25} = \frac{254+391.5}{25} = 25.82$.

\textbf{2.} $s_p^2 = \frac{9\times4.2 + 14\times5.8}{23} = \frac{37.8+81.2}{23} = \frac{119}{23} \approx 5.1739$.

\textbf{3.} With common $\sigma^2=5.0$: $\mathrm{Var}(\bar{X}_p) = \frac{\sigma^2}{n_1+n_2} = \frac{5.0}{25} = 0.2$. \qed
\end{exoresolu}

\courssub{Interval Estimation: Confidence Intervals}

\begin{definition}[Confidence Interval]
A \cle{confidence interval (CI)} for a parameter $\theta$ at confidence level $1-\alpha$ is an interval $(L, U)$ computed from the sample such that $P(L < \theta < U) = 1-\alpha$ (or $\ge 1-\alpha$ for discrete cases). The probability refers to the \emph{random interval} covering the fixed parameter in repeated sampling.
\textbf{Interpretation:} "We are $100(1-\alpha)\%$ confident that the true $\theta$ lies in this interval." This means that if we repeated the sampling many times, $100(1-\alpha)\%$ of the computed intervals would contain $\theta$.
\end{definition}

\begin{propriete}[Confidence Interval for $\mu$ -- Variance Known]
If $\sigma$ is known, or $n$ is large ($n\ge30$) so $s \approx \sigma$, and the population is normal or $n$ large (CLT):
\[
\bar{X} \pm z_{\alpha/2} \frac{\sigma}{\sqrt{n}} \quad\text{(infinite population or with replacement)},
\]
\[
\bar{X} \pm z_{\alpha/2} \frac{\sigma}{\sqrt{n}}\sqrt{1-\frac{n-1}{N-1}} \quad\text{(finite population without replacement)}.
\]
Common $z_{\alpha/2}$: $90\% \to 1.645$, $95\% \to 1.96$, $99\% \to 2.576$.
\end{propriete}

\begin{exoresolu}[CI for Mean -- Known Variance]
Electric bulb lifetimes: normal, $\sigma=100$ hours. Sample $n=40$, $\bar{x}=1200$ hours.
\begin{enumerate}
\item 95\% CI for $\mu$.
\item 99\% CI for $\mu$.
\item Interpretation.
\item Sample size needed for margin of error $\pm 20$ hours at 95\% confidence.
\end{enumerate}
\tcblower
\textbf{Solution:}

\textbf{1.} $z_{0.025}=1.96$, $\sigma/\sqrt{n} = 100/\sqrt{40} \approx 15.811$.
$E = 1.96 \times 15.811 = 30.99$.
95\% CI: $1200 \pm 30.99 \Rightarrow (1169.01,\; 1230.99)$ hours.

\textbf{2.} $z_{0.005}=2.576$.
$E = 2.576 \times 15.811 = 40.73$.
99\% CI: $1200 \pm 40.73 \Rightarrow (1159.27,\; 1240.73)$ hours.

\textbf{3.} We are 95\% confident that the true mean lifetime lies between 1169.01 and 1230.99 hours. In repeated sampling, 95\% of such intervals would contain $\mu$.

\textbf{4.} $n = \left(\frac{z_{\alpha/2}\sigma}{E}\right)^2 = \left(\frac{1.96\times100}{20}\right)^2 = 9.8^2 = 96.04 \to n=97$ (always round up). \qed
\end{exoresolu}

\begin{popfigure}
\begin{tikzpicture}
\begin{axis}[
    width=0.9\linewidth,
    height=5cm,
    xlabel={Lifetime (hours)},
    ylabel={Density},
    domain=1100:1300,
    samples=200,
    ymin=0,
    axis lines=middle,
    xtick={1169.01,1200,1230.99},
    xticklabels={1169.01,1200,1230.99},
]
% Normal curve for sample mean distribution
\addplot[thick, popblue] {1/(15.811*sqrt(max(0,2*pi)))*exp(-(x-1200)^2/(2*15.811^2))};
% Shaded CI
\addplot[fill=popblue!30, draw=none, domain=1169.01:1230.99] {1/(15.811*sqrt(max(0,2*pi)))*exp(-(x-1200)^2/(2*15.811^2))} \closedcycle;
% Mean line
\addplot[dashed, thick, popdark] coordinates {(1200,0) (1200,0.025)};
% CI bounds
\addplot[dashed, thick, poporange] coordinates {(1169.01,0) (1169.01,0.015)};
\addplot[dashed, thick, poporange] coordinates {(1230.99,0) (1230.99,0.015)};
\end{axis}
\end{tikzpicture}
\legende{95\% confidence interval for $\mu$: the shaded area under the sampling distribution of $\bar{X}$ corresponds to the central 95\% probability.}
\end{popfigure}

\begin{propriete}[Confidence Interval for $\mu$ -- Variance Unknown (t-Distribution)]
If the population is normal, $\sigma$ unknown, and $n$ is small ($n<30$), use the \cle{Student's $t$-distribution} with $\nu = n-1$ degrees of freedom:
\[
\bar{X} \pm t_{\alpha/2, n-1} \frac{S}{\sqrt{n}}.
\]
The $t$-distribution is symmetric, bell-shaped, with heavier tails than normal. As $\nu \to \infty$, $t \to N(0,1)$. For $n \ge 30$, $t$ and $z$ are very close, but using $t$ is always correct when $\sigma$ is estimated by $s$.
\end{propriete}

\begin{exoresolu}[CI for Mean -- Unknown Variance (t-Distribution)]
$n=15$ students' math scores (out of 100):
$62, 58, 71, 65, 59, 73, 68, 61, 66, 70, 64, 67, 69, 63, 72$.
Assume normality.
\begin{enumerate}
\item Compute $\bar{x}$ and $s$.
\item 95\% CI for $\mu$.
\item 90\% CI for $\mu$.
\item Effect of increasing $n$ to $60$ with same $\bar{x}, s$?
\end{enumerate}
\tcblower
\textbf{Solution:}

\textbf{1.} $\sum x_i = 988$, $\bar{x} = 988/15 = 65.8667$.
$\sum x_i^2 = 65384$.
$s^2 = \frac{1}{14}(65384 - 15\times65.8667^2) = \frac{1}{14}(65384 - 65076.6) = \frac{307.4}{14} = 21.957$.
$s = \sqrt{21.957} \approx 4.686$.

\textbf{2.} $\nu=14$, $t_{0.025,14}=2.145$.
$SE = s/\sqrt{15} = 4.686/3.873 = 1.210$.
$E = 2.145 \times 1.210 = 2.595$.
95\% CI: $65.867 \pm 2.595 \Rightarrow (63.272,\; 68.462)$.

\textbf{3.} $t_{0.05,14}=1.761$.
$E = 1.761 \times 1.210 = 2.131$.
90\% CI: $65.867 \pm 2.131 \Rightarrow (63.736,\; 67.998)$.

\textbf{4.} If $n=60$, $SE = 4.686/\sqrt{60} \approx 0.605$. $\nu=59$, $t_{0.025,59} \approx 2.001$.
$E \approx 2.001 \times 0.605 \approx 1.21$, roughly half the previous margin. Interval narrower. For $n=60$, normal approximation ($z=1.96$) is also acceptable. \qed
\end{exoresolu}

\begin{propriete}[Confidence Interval for $\mu_1 - \mu_2$ -- Two Independent Samples]
Let Sample 1: $n_1, \bar{x}_1, s_1$ (or $\sigma_1$); Sample 2: $n_2, \bar{x}_2, s_2$ (or $\sigma_2$). Independent.
\begin{enumerate}
\item \textbf{Variances known:} $(\bar{x}_1-\bar{x}_2) \pm z_{\alpha/2}\sqrt{\frac{\sigma_1^2}{n_1}+\frac{\sigma_2^2}{n_2}}$.
\item \textbf{Variances unknown, assumed equal (pooled $t$):}
$s_p^2 = \frac{(n_1-1)s_1^2+(n_2-1)s_2^2}{n_1+n_2-2}$, $\nu = n_1+n_2-2$.
$(\bar{x}_1-\bar{x}_2) \pm t_{\alpha/2,\nu}\, s_p\sqrt{\frac{1}{n_1}+\frac{1}{n_2}}$.
\item \textbf{Variances unknown, not assumed equal (Welch's $t$):}
$(\bar{x}_1-\bar{x}_2) \pm t_{\alpha/2,\nu}\, \sqrt{\frac{s_1^2}{n_1}+\frac{s_2^2}{n_2}}$,
with $\nu \approx \frac{(s_1^2/n_1 + s_2^2/n_2)^2}{(s_1^2/n_1)^2/(n_1-1) + (s_2^2/n_2)^2/(n_2-1)}$ (Welch--Satterthwaite formula).
\end{enumerate}
For large samples ($n_1,n_2 \ge 30$), $z$ can replace $t$.
\end{propriete}

\begin{exoresolu}[CI for Difference of Means -- Pooled t (Cassava Yields)]
Variety A ($n_1=12$): $12.1, 11.8, 12.5, 12.0, 11.9, 12.3, 12.2, 11.7, 12.4, 12.0, 11.9, 12.2$.
Variety B ($n_2=10$): $13.2, 12.9, 13.5, 13.0, 12.8, 13.3, 13.1, 12.7, 13.4, 13.0$.
Assume normal, common variance.
\begin{enumerate}
\item Sample means and variances.
\item Pooled variance $s_p^2$.
\item 95\% CI for $\mu_A - \mu_B$.
\item Interpretation.
\end{enumerate}
\tcblower
\textbf{Solution:}

\textbf{1. Variety A:} $\sum x = 145.0$, $\bar{x}_A = 12.0833$.
$\sum x^2 = 1752.74$, $s_A^2 = \frac{1}{11}(1752.74 - 12\times12.0833^2) = \frac{0.67}{11} = 0.0609$, $s_A \approx 0.2468$.

\textbf{Variety B:} $\sum x = 130.9$, $\bar{x}_B = 13.09$.
$\sum x^2 = 1714.09$, $s_B^2 = \frac{1}{9}(1714.09 - 10\times13.09^2) = \frac{0.61}{9} = 0.0678$, $s_B \approx 0.2604$.

\textbf{2.} $s_p^2 = \frac{11\times0.0609 + 9\times0.0678}{20} = \frac{0.6699+0.6102}{20} = 0.064005$, $s_p \approx 0.2530$.

\textbf{3.} $\nu=20$, $t_{0.025,20}=2.086$.
Point estimate: $\bar{x}_A - \bar{x}_B = -1.0067$.
$SE = s_p\sqrt{\frac{1}{12}+\frac{1}{10}} = 0.2530\sqrt{0.18333} = 0.1083$.
$E = 2.086 \times 0.1083 = 0.2259$.
95\% CI: $-1.0067 \pm 0.2259 \Rightarrow (-1.2326,\; -0.7808)$.

\textbf{4.} We are 95\% confident that $\mu_A - \mu_B$ is between $-1.23$ and $-0.78$ t/ha. Since the interval is entirely negative, Variety B has a significantly higher mean yield than Variety A at the 5\% level. \qed
\end{exoresolu}

\begin{attention}[Confidence Intervals -- Common Errors]
\begin{itemize}
\item Using $z$ instead of $t$ when $\sigma$ is unknown and $n<30$. Always use $t$ for small samples from normal populations with unknown variance.
\item Forgetting the finite population correction when $n/N > 0.05$ (or $>5\%$).
\item Misinterpreting the CI: "There is a 95\% probability that $\mu$ is in this interval" is \textbf{wrong}. $\mu$ is fixed; the interval is random. Correct: "95\% of such intervals constructed from repeated samples would contain $\mu$."
\item For two-sample CIs, using the pooled formula when variances are not equal. Check the assumption (e.g., $F$-test for equality of variances, or use Welch's interval by default if unsure).
\item Confusing the standard error of the \emph{mean} ($\sigma/\sqrt{n}$) with the standard deviation of the \emph{population} ($\sigma$).
\end{itemize}
\end{attention}

\courssub{Activities of Integration and Assessment}

\begin{exercice}[Integration Exercise 1 -- Sampling Distribution]
A population consists of the values $3, 5, 7, 9, 11$.
\begin{enumerate}
\item Find $\mu$ and $\sigma^2$.
\item List all samples of size $n=2$ without replacement and their means.
\item Construct the sampling distribution of $\bar{X}$.
\item Verify $E(\bar{X})=\mu$ and $\mathrm{Var}(\bar{X}) = \frac{\sigma^2}{n}\left(1-\frac{n-1}{N-1}\right)$.
\item Repeat for sampling \emph{with} replacement (list the distribution of $\bar{X}$ and its variance).
\end{enumerate}
\end{exercice}

\begin{exercice}[Integration Exercise 2 -- Random Sampling and Attributes]
A school has 500 students. A researcher wants to estimate the proportion who eat in the canteen. Using the random number table below (3-digit groups), select an SRS of 30 students. If 18 of them eat in the canteen, find $\hat{p}$ and its estimated standard error (with fpc).

\begin{center}
\begin{tabular}{cccccc}
123 & 456 & 001 & 499 & 234 & 501 \\
012 & 345 & 678 & 099 & 100 & 200 \\
300 & 400 & 500 & 002 & 003 & 004 \\
005 & 006 & 007 & 008 & 009 & 010 \\
011 & 012 & 013 & 014 & 015 & 016 \\
\end{tabular}
\end{center}
\end{exercice}

\begin{exercice}[Integration Exercise 3 -- CLT]
The time (minutes) to complete a task has mean $25$ and standard deviation $5$. Distribution unknown. A sample of $64$ workers is timed.
\begin{enumerate}
\item Approximate distribution of $\bar{X}$.
\item $P(\bar{X} > 26)$.
\item $P(24.5 < \bar{X} < 25.5)$.
\item What sample size is needed so that $P(|\bar{X}-\mu| < 0.5) \ge 0.95$?
\end{enumerate}
\end{exercice}

\begin{exercice}[Integration Exercise 4 -- Point Estimation]
Let $X_1,X_2,X_3,X_4$ be a random sample from a population with mean $\mu$, variance $\sigma^2$. Consider:
$\hat{\mu}_1 = \frac{X_1+X_2+X_3+X_4}{4}$,
$\hat{\mu}_2 = \frac{X_1+2X_2+3X_3+4X_4}{10}$,
$\hat{\mu}_3 = \frac{X_1+X_4}{2}$.
\begin{enumerate}
\item Which are unbiased?
\item Find the variance of each unbiased estimator.
\item Which is the most efficient?
\item If $\hat{\mu}_4 = \frac{X_1+X_2+X_3}{3}$ (ignoring $X_4$), compare its efficiency to $\hat{\mu}_1$.
\end{enumerate}
\end{exercice}

\begin{exercice}[Integration Exercise 5 -- Pooled Estimation]
Two independent samples from normal populations with common variance $\sigma^2$:
Sample 1: $n_1=8$, $\bar{x}_1=15.2$, $s_1^2=3.6$.
Sample 2: $n_2=12$, $\bar{x}_2=16.5$, $s_2^2=4.4$.
\begin{enumerate}
\item Pooled estimate of $\mu$ (assuming common mean).
\item Pooled estimate of $\sigma^2$.
\item 95\% CI for the common mean $\mu$.
\item 95\% CI for $\mu_1 - \mu_2$ (do not assume common mean).
\end{enumerate}
\end{exercice}

\begin{exercice}[Integration Exercise 6 -- Confidence Intervals]
\begin{enumerate}
\item A sample of $25$ items from a normal population gives $\bar{x}=50$, $s=8$. Find 99\% CI for $\mu$.
\item A sample of $100$ items gives $\bar{x}=50$, $s=8$. Find 99\% CI for $\mu$ (use $z$ or $t$? justify).
\item Two samples: $n_1=20, \bar{x}_1=100, s_1=10$; $n_2=25, \bar{x}_2=95, s_2=12$. Assume equal variances. Find 95\% CI for $\mu_1-\mu_2$.
\item Same data but do \emph{not} assume equal variances. Find 95\% CI using Welch's method.
\end{enumerate}
\end{exercice}

\begin{corrige}[Integration Exercise 1]
\textbf{1.} $\mu = 7$, $\sigma^2 = \frac{1}{5}[(3-7)^2+(5-7)^2+(7-7)^2+(9-7)^2+(11-7)^2] = \frac{16+4+0+4+16}{5} = 8$.
\textbf{2.} $\binom{5}{2}=10$ samples: $(3,5)\to4$, $(3,7)\to5$, $(3,9)\to6$, $(3,11)\to7$, $(5,7)\to6$, $(5,9)\to7$, $(5,11)\to8$, $(7,9)\to8$, $(7,11)\to9$, $(9,11)\to10$.
\textbf{3.} Distribution: $\bar{x}=4,5,6,7,8,9,10$ with probs $0.1,0.1,0.2,0.2,0.2,0.1,0.1$.
\textbf{4.} $E(\bar{X})=7=\mu$. $E(\bar{X}^2)=54.6$, $\mathrm{Var}=54.6-49=5.6$. Formula: $\frac{8}{2}(1-\frac{1}{4})=4\times0.75=3$? Wait: $\frac{\sigma^2}{n}(1-\frac{n-1}{N-1}) = \frac{8}{2}(1-\frac{1}{4}) = 4 \times 0.75 = 3$. But computed variance is 5.6? Let's recalculate.
$\bar{x}$ values: 4,5,6,7,6,7,8,8,9,10. Sum = 70, mean = 7. Squares: 16,25,36,49,36,49,64,64,81,100. Sum = 520. Mean square = 52. Var = 52 - 49 = 3. Yes, 3. My previous $E(\bar{X}^2)$ was wrong. Correct: 520/10=52. Var=3. Matches.
\textbf{5.} With replacement: $N^n=25$ ordered samples. $\mathrm{Var}(\bar{X}) = \sigma^2/n = 8/2 = 4$. Distribution has more spread.
\end{corrige}

\begin{corrige}[Integration Exercise 2]
Valid numbers 001-500: 123, 456, 001, 499, 234, 012, 345, 099, 100, 200, 300, 400, 002, 003, 004, 005, 006, 007, 008, 009, 010, 011, 013, 014, 015, 016 (012 duplicate, 501>500, 678>500, 500 is valid? 500 is within 1-500, yes. But 500 appears after 400. Let's list carefully: 123, 456, 001, 499, 234, (501 reject), 012, 345, (678 reject), 099, 100, 200, 300, 400, 500, 002, 003, 004, 005, 006, 007, 008, 009, 010, 011, (012 duplicate), 013, 014, 015, 016. That's 30 numbers. $\hat{p}=18/30=0.6$. $N=500, n=30$. $\widehat{SE} = \sqrt{\frac{0.6\times0.4}{30}(1-\frac{29}{499})} = \sqrt{0.008 \times 0.9419} = \sqrt{0.007535} = 0.0868$.
\end{corrige}

\begin{corrige}[Integration Exercise 3]
\textbf{1.} $\bar{X} \sim N(25, 5^2/64) = N(25, 0.390625)$, $\sigma_{\bar{X}}=0.625$.
\textbf{2.} $Z = (26-25)/0.625 = 1.6$. $P(Z>1.6) = 1-0.9452 = 0.0548$.
\textbf{3.} $Z_1 = (24.5-25)/0.625 = -0.8$, $Z_2 = 0.8$. $P(-0.8<Z<0.8) = 2\times0.7881-1 = 0.5762$.
\textbf{4.} $P(|\bar{X}-\mu|<0.5) = P(|Z| < 0.5\sqrt{n}/5) = P(|Z| < 0.1\sqrt{n}) \ge 0.95$. Need $0.1\sqrt{n} \ge 1.96 \Rightarrow \sqrt{n} \ge 19.6 \Rightarrow n \ge 384.16 \to n=385$.
\end{corrige}

\begin{corrige}[Integration Exercise 4]
\textbf{1.} $E(\hat{\mu}_1)=\mu$, $E(\hat{\mu}_2)=\frac{1+2+3+4}{10}\mu=\mu$, $E(\hat{\mu}_3)=\mu$. All unbiased.
\textbf{2.} $\mathrm{Var}(\hat{\mu}_1)=\sigma^2/4=0.25\sigma^2$. $\mathrm{Var}(\hat{\mu}_2)=\frac{1+4+9+16}{100}\sigma^2 = 30/100\sigma^2=0.3\sigma^2$. $\mathrm{Var}(\hat{\mu}_3)=\frac{1+1}{4}\sigma^2=0.5\sigma^2$.
\textbf{3.} $\hat{\mu}_1$ most efficient (smallest variance).
\textbf{4.} $\mathrm{Var}(\hat{\mu}_4)=\sigma^2/3 \approx 0.333\sigma^2$. Relative efficiency of $\hat{\mu}_4$ w.r.t $\hat{\mu}_1$ = $(0.25)/(0.333) = 0.75$. $\hat{\mu}_1$ uses all 4 observations, so more efficient.
\end{corrige}

\begin{corrige}[Integration Exercise 5]
\textbf{1.} $\bar{x}_p = \frac{8\times15.2 + 12\times16.5}{20} = \frac{121.6+198}{20} = 15.98$.
\textbf{2.} $s_p^2 = \frac{7\times3.6 + 11\times4.4}{18} = \frac{25.2+48.4}{18} = 73.6/18 = 4.0889$.
\textbf{3.} CI for common $\mu$: $\nu=18$, $t_{0.025,18}=2.101$. $SE = s_p/\sqrt{20} = \sqrt{4.0889/20} = \sqrt{0.2044} = 0.4521$. $E=2.101\times0.4521=0.950$. CI: $15.98 \pm 0.95 \Rightarrow (15.03, 16.93)$.
\textbf{4.} CI for $\mu_1-\mu_2$ (pooled): point estimate $15.2-16.5=-1.3$. $SE = s_p\sqrt{1/8+1/12} = 2.022\sqrt{0.20833} = 2.022\times0.4564 = 0.923$. $\nu=18$, $t=2.101$. $E=1.939$. CI: $-1.3 \pm 1.939 \Rightarrow (-3.239, 0.639)$.
\end{corrige}

\begin{corrige}[Integration Exercise 6]
\textbf{1.} $n=25$, $\nu=24$, $t_{0.005,24}=2.797$. $SE=8/5=1.6$. $E=2.797\times1.6=4.475$. CI: $(45.525, 54.475)$.
\textbf{2.} $n=100$, large sample. Can use $z=2.576$ or $t_{0.005,99}\approx2.626$. $SE=8/10=0.8$. Using $z$: $E=2.576\times0.8=2.061$. CI: $(47.939, 52.061)$. Using $t$: $E=2.626\times0.8=2.101$. CI: $(47.899, 52.101)$. Both acceptable; $t$ is more conservative.
\textbf{3.} Pooled: $s_p^2 = \frac{19\times100 + 24\times144}{43} = \frac{1900+3456}{43} = 5356/43 = 124.558$. $s_p=11.16$. $\nu=43$, $t_{0.025,43}\approx2.017$. $SE = 11.16\sqrt{1/20+1/25} = 11.16\sqrt{0.09} = 11.16\times0.3 = 3.348$. $E=2.017\times3.348=6.75$. CI: $5 \pm 6.75 \Rightarrow (-1.75, 11.75)$.
\textbf{4.} Welch: $SE = \sqrt{100/20 + 144/25} = \sqrt{5 + 5.76} = \sqrt{10.76} = 3.280$. $\nu \approx \frac{(5+5.76)^2}{5^2/19 + 5.76^2/24} = \frac{115.78}{1.316+1.382} = 115.78/2.698 = 42.9 \to 43$. $t_{0.025,43}\approx2.017$. $E=2.017\times3.280=6.616$. CI: $5 \pm 6.616 \Rightarrow (-1.616, 11.616)$. Similar to pooled because sample sizes and variances are not too different.
\end{corrige}

\newpage

\coursec{Exercises}

\courssub{I check my knowledge}

\begin{exercice}[Sampling with replacement from a finite population]
A population consists of the four values $1,\ 3,\ 5,\ 7$.
\begin{enumerate}
\item List all possible samples of size $2$ drawn \emph{with replacement}.
\item Construct the sampling distribution of the sample mean $\bar{X}$.
\item Verify that $\operatorname{E}(\bar{X}) = \mu$ and $\operatorname{Var}(\bar{X}) = \dfrac{\sigma^2}{2}$.
\end{enumerate}
\end{exercice}

\begin{exercice}[Using random number tables]
An extract from a random number table is given below:
\[
84\ 12\ 57\ 03\ 91\ 46\ 28\ 73\ 65\ 19\ 04\ 88\ 32\ 50\ 77
\]
A researcher wishes to select a random sample of $6$ candidates from a numbered list of $85$ applicants (numbered $01$ to $85$). Using the numbers in the order given, select the sample and explain clearly why any numbers are rejected.
\end{exercice}

\begin{exercice}[True/False with justification]
State whether each of the following statements is true or false. Justify your answer in each case.
\begin{enumerate}
\item The Central Limit Theorem states that the sample mean $\bar{X}$ is normally distributed for any sample size $n$.
\item If $\hat{\theta}$ is an unbiased estimator of $\theta$, then $\operatorname{E}(\hat{\theta}) = \theta$ for all possible samples.
\item The pooled estimate of the common variance from two independent samples is always the arithmetic mean of the two sample variances.
\item A $95\%$ confidence interval for a population mean $\mu$ means that there is a $0.95$ probability that $\mu$ lies in the calculated interval.
\end{enumerate}
\end{exercice}

\begin{exercice}[Direct calculations on point estimation]
A random sample of $8$ observations yields $\sum x = 164$ and $\sum x^2 = 3\,412$.
\begin{enumerate}
\item Calculate an unbiased estimate of the population mean $\mu$.
\item Calculate an unbiased estimate of the population variance $\sigma^2$.
\end{enumerate}
\end{exercice}

\courssub{I practise}

\begin{exercice}[Probability using the Central Limit Theorem]
The masses of loaves produced by a bakery are distributed with mean $500\ \mathrm{g}$ and standard deviation $12\ \mathrm{g}$. A random sample of $36$ loaves is selected.
\begin{enumerate}
\item State the theorem that allows you to approximate the distribution of the sample mean.
\item Find the probability that the mean mass of the sample exceeds $504\ \mathrm{g}$.
\end{enumerate}
\end{exercice}

\begin{exercice}[Sampling from a normal distribution]
The heights of students in a large school are normally distributed with mean $170\ \mathrm{cm}$ and variance $25\ \mathrm{cm}^2$. A random sample of $16$ students is taken.
Find the probability that the sample mean height lies between $168\ \mathrm{cm}$ and $173\ \mathrm{cm}$.
\end{exercice}

\begin{exercice}[Stratified sampling]
A school has $400$ Lower Sixth students and $600$ Upper Sixth students. The head teacher wants to select a stratified sample of $50$ students using proportional allocation.
\begin{enumerate}
\item Determine how many students should be selected from each stratum.
\item Describe briefly how the selection within each stratum should be carried out.
\item State one advantage of this stratified sampling method over simple random sampling of $50$ students from the whole school.
\end{enumerate}
\end{exercice}

\begin{exercice}[Comparing estimators]
Let $X_1, X_2$ be a random sample of size $2$ from a population with mean $\mu$ and variance $\sigma^2$. Consider the two estimators of $\mu$:
\[
\bar{X} = \frac{X_1 + X_2}{2}, \qquad T = \frac{X_1 + 2X_2}{3}.
\]
\begin{enumerate}
\item Show that both $\bar{X}$ and $T$ are unbiased estimators of $\mu$.
\item Calculate $\operatorname{Var}(\bar{X})$ and $\operatorname{Var}(T)$.
\item Which estimator is more efficient? Justify your answer.
\end{enumerate}
\end{exercice}

\courssub{I prepare for the GCE}

\begin{exercice}[Pooled estimation and confidence interval for difference of means] \hfill (12 marks)
Two independent random samples are drawn from normal populations with a common variance $\sigma^2$.
Sample 1: $n_1 = 10,\ \bar{x}_1 = 25.2,\ s_1^2 = 4.1$
Sample 2: $n_2 = 15,\ \bar{x}_2 = 24.6,\ s_2^2 = 3.6$
\begin{enumerate}
\item Assuming the two populations have a common mean $\mu$, obtain the pooled unbiased estimate of $\mu$. \hfill (2 marks)
\item Obtain the pooled unbiased estimate of the common population variance $\sigma^2$. \hfill (3 marks)
\item Construct a $95\%$ confidence interval for the difference of the population means $\mu_1 - \mu_2$. \hfill (4 marks)
\item Interpret your interval in the context of the problem. \hfill (3 marks)
\end{enumerate}
\end{exercice}

\begin{exercice}[Confidence intervals and sample size determination] \hfill (14 marks)
A sample of $100$ bags of rice has a mean mass of $49.8\ \mathrm{kg}$. The population standard deviation is known to be $1.2\ \mathrm{kg}$.
\begin{enumerate}
\item Construct a $95\%$ confidence interval for the true mean mass of all bags. \hfill (4 marks)
\item Determine the minimum sample size required so that a $95\%$ confidence interval for the mean has a total width of at most $0.4\ \mathrm{kg}$. \hfill (5 marks)
\item If the population standard deviation were unknown and the sample standard deviation $s = 1.2\ \mathrm{kg}$ were used instead, explain how the confidence interval in (i) would change and why. \hfill (5 marks)
\end{enumerate}
\end{exercice}

\begin{exercice}[GCE-style: Hypothesis test via confidence interval] \hfill (16 marks)
A factory fills sugar packets labelled $500\ \mathrm{g}$. A random sample of $64$ packets is weighed. The results are summarised by $\sum x = 32\,064\ \mathrm{g}$ and $\sum (x - \bar{x})^2 = 1\,008\ \mathrm{g}^2$.
\begin{enumerate}
\item Calculate the sample mean and the unbiased estimate of the population variance. \hfill (3 marks)
\item Construct a $90\%$ confidence interval for the true mean content of the packets. \hfill (5 marks)
\item The factory manager claims that the mean content is $502\ \mathrm{g}$. Use your confidence interval to test this claim at the $10\%$ significance level. State your conclusion clearly. \hfill (4 marks)
\item State one assumption you have made in constructing the confidence interval. \hfill (2 marks)
\item If the sample size had been $16$ instead of $64$, with the same sample mean and unbiased variance estimate, explain how the width of the $90\%$ confidence interval would change. \hfill (2 marks)
\end{enumerate}
\end{exercice}

\newpage

\coursec{Solutions to the exercises}

\coursec{Activities of Integration, Assessment and Remediation}

\courssub{Worked Examination Questions}

\begin{corrige}[Exercise 1 — Sampling with replacement from a finite population]
The population consists of the four values \(1, 3, 5, 7\). We draw samples of size \(n=2\) \emph{with replacement}.

\begin{enumerate}
\item \textbf{All possible samples (ordered pairs):}
Since each draw has 4 possibilities and draws are independent, there are \(4^2 = 16\) equally likely ordered samples.
\[
\begin{array}{cccc}
(1,1), & (1,3), & (1,5), & (1,7), \\
(3,1), & (3,3), & (3,5), & (3,7), \\
(5,1), & (5,3), & (5,5), & (5,7), \\
(7,1), & (7,3), & (7,5), & (7,7).
\end{array}
\]

\item \textbf{Sampling distribution of the sample mean \(\bar{X}\):}
For each sample we compute \(\bar{x} = \dfrac{x_1+x_2}{2}\). Grouping samples yielding the same mean:

\[
\begin{array}{c|c}
\text{Sample(s)} & \bar{x} \\ \hline
(1,1) & 1 \\
(1,3),\ (3,1) & 2 \\
(1,5),\ (3,3),\ (5,1) & 3 \\
(1,7),\ (3,5),\ (5,3),\ (7,1) & 4 \\
(3,7),\ (5,5),\ (7,3) & 5 \\
(5,7),\ (7,5) & 6 \\
(7,7) & 7
\end{array}
\]

Each ordered pair has probability \(\frac{1}{16}\). The probability distribution of \(\bar{X}\) is therefore:

\[
\begin{array}{c|ccccccc}
\bar{x} & 1 & 2 & 3 & 4 & 5 & 6 & 7 \\ \hline
P(\bar{X}=\bar{x}) & \frac{1}{16} & \frac{2}{16} & \frac{3}{16} & \frac{4}{16} & \frac{3}{16} & \frac{2}{16} & \frac{1}{16}
\end{array}
\]

\item \textbf{Verification of \(\operatorname{E}(\bar{X})=\mu\) and \(\operatorname{Var}(\bar{X})=\frac{\sigma^2}{n}\):}
Population parameters:
\[
\mu = \frac{1+3+5+7}{4} = 4,\qquad
\sigma^2 = \frac{(1-4)^2+(3-4)^2+(5-4)^2+(7-4)^2}{4} = \frac{9+1+1+9}{4} = 5.
\]

Expectation of \(\bar{X}\):
\[
\operatorname{E}(\bar{X}) = \sum \bar{x}\,P(\bar{X}=\bar{x})
= \frac{1\cdot1 + 2\cdot2 + 3\cdot3 + 4\cdot4 + 5\cdot3 + 6\cdot2 + 7\cdot1}{16}
= \frac{64}{16} = 4 = \mu.
\]

Variance of \(\bar{X}\):
\[
\operatorname{E}(\bar{X}^2) = \frac{1^2\cdot1 + 2^2\cdot2 + 3^2\cdot3 + 4^2\cdot4 + 5^2\cdot3 + 6^2\cdot2 + 7^2\cdot1}{16}
= \frac{296}{16} = 18.5.
\]
\[
\operatorname{Var}(\bar{X}) = \operatorname{E}(\bar{X}^2) - [\operatorname{E}(\bar{X})]^2 = 18.5 - 16 = 2.5.
\]
Theoretical value for sampling \emph{with replacement}: \(\dfrac{\sigma^2}{n} = \dfrac{5}{2} = 2.5\). The result is verified.
\end{enumerate}
\end{corrige}

\begin{corrige}[Exercise 2 — Using random number tables for simple random sampling without replacement]
We have the following random number sequence (grouped in two‑digit blocks):
\[
84\ 12\ 57\ 03\ 91\ 46\ 28\ 73\ 65\ 19\ 04\ 88\ 32\ 50\ 77
\]
Population: \(85\) applicants numbered \(01\) to \(85\). Required: a simple random sample of size \(6\) \emph{without replacement}.

\textbf{Procedure:} Read the two‑digit numbers sequentially. Accept a number if it lies between \(01\) and \(85\) inclusive \emph{and} has not been selected before; otherwise reject it. Stop when \(6\) distinct numbers are obtained.

\begin{center}
\begin{tabular}{|c|c|l|}
\hline
\textbf{Random Number} & \textbf{Decision} & \textbf{Reason} \\ \hline
84 & Accept & In range, not previously chosen \\
12 & Accept & In range, not previously chosen \\
57 & Accept & In range, not previously chosen \\
03 & Accept & In range, not previously chosen \\
91 & Reject & \(91 > 85\) (outside population range) \\
46 & Accept & In range, not previously chosen \\
28 & Accept & In range, not previously chosen \\ \hline
\end{tabular}
\end{center}

We have now obtained the required \(6\) distinct numbers. The sampling stops here; the remaining numbers (\(73, 65, 19, 04, 88, 32, 50, 77\)) are not used.

\textbf{Selected sample (sorted):} Applicants numbered \(03,\ 12,\ 28,\ 46,\ 57,\ 84\).
\end{corrige}

\begin{corrige}[Exercise 3 — True/False with justification]
\begin{enumerate}
\item \textbf{Statement:} ``The Central Limit Theorem guarantees that the sample mean is normally distributed for any sample size.'' \\
\textbf{False.} The Central Limit Theorem states that for a \emph{large} sample size \(n\) (usually \(n \ge 30\)), the sampling distribution of the sample mean \(\bar{X}\) is \emph{approximately} normal, regardless of the shape of the population distribution. For small \(n\), normality of \(\bar{X}\) is guaranteed only if the population itself is normal.

\item \textbf{Statement:} ``An unbiased estimator always gives the exact value of the parameter for a given sample.'' \\
\textbf{False.} An estimator \(\hat{\theta}\) is unbiased if its \emph{expected value} over all possible samples equals the parameter: \(\operatorname{E}(\hat{\theta}) = \theta\). This does \emph{not} mean that the estimate from a particular sample equals \(\theta\); it means that on average, over repeated sampling, the estimator hits the target.

\item \textbf{Statement:} ``The pooled estimate of a common variance is simply the arithmetic mean of the two sample variances.'' \\
\textbf{False.} The pooled estimate of a common variance from two independent samples is a \emph{weighted} average of the two sample variances, with weights equal to the degrees of freedom \((n_1-1)\) and \((n_2-1)\):
\[
s_p^2 = \frac{(n_1-1)s_1^2 + (n_2-1)s_2^2}{n_1+n_2-2}.
\]
It reduces to the arithmetic mean only when \(n_1 = n_2\).

\item \textbf{Statement:} ``A \(95\%\) confidence interval for \(\mu\) means there is a \(95\%\) probability that \(\mu\) lies in the calculated interval.'' \\
\textbf{False.} A \(95\%\) confidence interval for \(\mu\) means that if we were to repeat the sampling process many times and construct a \(95\%\) interval each time, approximately \(95\%\) of those intervals would contain the true \(\mu\). For a \emph{specific} calculated interval, the true \(\mu\) is either inside it (probability \(1\)) or outside it (probability \(0\)); \(\mu\) is a fixed unknown constant, not a random variable.
\end{enumerate}
\end{corrige}

\begin{corrige}[Exercise 4 — Direct calculations on point estimation]
Given a random sample of size \(n = 8\) with \(\sum x = 164\) and \(\sum x^2 = 3412\).

\begin{enumerate}
\item \textbf{Unbiased estimate of the population mean \(\mu\):}
The sample mean \(\bar{x}\) is an unbiased estimator of \(\mu\).
\[
\bar{x} = \frac{\sum x}{n} = \frac{164}{8} = 20.5.
\]

\item \textbf{Unbiased estimate of the population variance \(\sigma^2\):}
The sample variance \(s^2 = \dfrac{\sum x^2 - \frac{(\sum x)^2}{n}}{n-1}\) is an unbiased estimator of \(\sigma^2\).
\[
s^2 = \frac{3412 - \frac{164^2}{8}}{7}
= \frac{3412 - \frac{26896}{8}}{7}
= \frac{3412 - 3362}{7}
= \frac{50}{7} \approx 7.143.
\]
\end{enumerate}
\end{corrige}

\begin{corrige}[Exercise 5 — Probability calculation using the Central Limit Theorem]
Population: mean \(\mu = 500\,\mathrm{g}\), standard deviation \(\sigma = 12\,\mathrm{g}\). Sample size \(n = 36\).

\begin{enumerate}
\item Because \(n = 36 \ge 30\), the \textbf{Central Limit Theorem} allows us to approximate the distribution of the sample mean \(\bar{X}\) by a normal distribution:
\[
\bar{X} \sim N\left(\mu,\ \frac{\sigma^2}{n}\right) \quad\text{approximately}.
\]

\item Standard error of the mean:
\[
\sigma_{\bar{X}} = \frac{\sigma}{\sqrt{n}} = \frac{12}{\sqrt{36}} = 2\,\mathrm{g}.
\]
Required probability:
\[
P(\bar{X} > 504) = P\left(Z > \frac{504-500}{2}\right) = P(Z > 2).
\]
From standard normal tables, \(\Phi(2) = 0.9772\), hence
\[
P(Z > 2) = 1 - 0.9772 = 0.0228.
\]
\end{enumerate}
\end{corrige}

\begin{corrige}[Exercise 6 — Sampling from a normal population (exact distribution)]
Population: heights \(\sim N(\mu=170\,\mathrm{cm},\ \sigma^2=25\,\mathrm{cm}^2)\), so \(\sigma = 5\,\mathrm{cm}\). Sample size \(n = 16\).

Since the population is \emph{exactly} normal, the sample mean \(\bar{X}\) is \emph{exactly} normally distributed for any \(n\):
\[
\bar{X} \sim N\left(170,\ \frac{25}{16}\right),\qquad
\sigma_{\bar{X}} = \frac{5}{\sqrt{16}} = 1.25\,\mathrm{cm}.
\]

\[
P(168 < \bar{X} < 173) = P\left(\frac{168-170}{1.25} < Z < \frac{173-170}{1.25}\right)
= P(-1.6 < Z < 2.4).
\]
Using standard normal tables:
\[
\Phi(2.4) = 0.9918,\qquad \Phi(-1.6) = 0.0548.
\]
\[
P = 0.9918 - 0.0548 = 0.9370.
\]
\end{corrige}

\begin{corrige}[Exercise 7 — Stratified sampling with proportional allocation]
A school has \(400\) Lower Sixth and \(600\) Upper Sixth students (total \(N=1000\)). A sample of \(n=50\) is to be selected using stratified sampling with proportional allocation.

\begin{enumerate}
\item \textbf{Sample sizes per stratum:}
\[
n_{\text{Lower}} = \frac{400}{1000} \times 50 = 20,\qquad
n_{\text{Upper}} = \frac{600}{1000} \times 50 = 30.
\]

\item \textbf{Selection procedure within each stratum:}
Obtain a complete list of all Lower Sixth students (numbered \(1\) to \(400\)) and a separate list of all Upper Sixth students (numbered \(1\) to \(600\)). Use simple random sampling (e.g., random number tables or a computer random number generator) \emph{independently} on each list to select \(20\) Lower Sixth and \(30\) Upper Sixth students.

\item \textbf{Advantage over simple random sampling:}
Stratified sampling guarantees that both strata are represented in the sample in the correct proportions. If the variable of interest (e.g., height, opinion on a policy) differs systematically between Lower and Upper Sixth students, stratified sampling typically yields a more precise estimate (smaller sampling variance) than a simple random sample of the same total size, because it eliminates the between‑strata variability from the sampling error.
\end{enumerate}
\end{corrige}

\begin{corrige}[Exercise 8 — Comparing two unbiased estimators]
Let \(X_1, X_2\) be a random sample of size \(2\) from a population with mean \(\mu\) and variance \(\sigma^2\). Consider the estimators:
\[
\bar{X} = \frac{X_1+X_2}{2},\qquad T = \frac{X_1+2X_2}{3}.
\]

\begin{enumerate}
\item \textbf{Unbiasedness:}
\[
\operatorname{E}(\bar{X}) = \frac{\operatorname{E}(X_1)+\operatorname{E}(X_2)}{2} = \frac{\mu+\mu}{2} = \mu.
\]
\[
\operatorname{E}(T) = \frac{\operatorname{E}(X_1)+2\operatorname{E}(X_2)}{3} = \frac{\mu+2\mu}{3} = \mu.
\]
Both \(\bar{X}\) and \(T\) are unbiased estimators of \(\mu\).

\item \textbf{Variances (using independence):}
\[
\operatorname{Var}(\bar{X}) = \frac{1}{4}\bigl(\operatorname{Var}(X_1)+\operatorname{Var}(X_2)\bigr) = \frac{1}{4}(\sigma^2+\sigma^2) = \frac{\sigma^2}{2}.
\]
\[
\operatorname{Var}(T) = \frac{1}{9}\bigl(\operatorname{Var}(X_1)+4\operatorname{Var}(X_2)\bigr) = \frac{1}{9}(\sigma^2+4\sigma^2) = \frac{5\sigma^2}{9}.
\]

\item \textbf{Efficiency comparison:}
\[
\frac{\sigma^2}{2} = \frac{4.5\sigma^2}{9} < \frac{5\sigma^2}{9} \quad\Longrightarrow\quad \operatorname{Var}(\bar{X}) < \operatorname{Var}(T).
\]
The estimator with the smaller variance is more efficient. Therefore \(\bar{X}\) is more efficient than \(T\).
\end{enumerate}
\end{corrige}

\begin{corrige}[Exercise 9 — Pooled estimation and confidence interval for the difference of two means]
Two independent random samples are drawn from normal populations with a \emph{common} variance \(\sigma^2\).
Sample 1: \(n_1=10,\ \bar{x}_1=25.2,\ s_1^2=4.1\)
Sample 2: \(n_2=15,\ \bar{x}_2=24.6,\ s_2^2=3.6\)

\begin{enumerate}
\item \textbf{Pooled unbiased estimate of the common mean \(\mu\) (2 marks):}
\[
\hat{\mu} = \frac{n_1\bar{x}_1 + n_2\bar{x}_2}{n_1+n_2}
= \frac{10\times25.2 + 15\times24.6}{25}
= \frac{252 + 369}{25}
= \frac{621}{25} = 24.84.
\]

\item \textbf{Pooled unbiased estimate of the common variance \(\sigma^2\) (3 marks):}
\[
s_p^2 = \frac{(n_1-1)s_1^2 + (n_2-1)s_2^2}{n_1+n_2-2}
= \frac{9\times4.1 + 14\times3.6}{23}
= \frac{36.9 + 50.4}{23}
= \frac{87.3}{23} \approx 3.7957.
\]

\item \textbf{95\% confidence interval for \(\mu_1-\mu_2\) (4 marks):}
Degrees of freedom \(\nu = n_1+n_2-2 = 23\). From \(t\)-tables, \(t_{0.025,23} = 2.069\).
Pooled standard deviation: \(s_p = \sqrt{3.7957} \approx 1.9482\).
Standard error of the difference:
\[
\text{SE} = s_p\sqrt{\frac{1}{n_1}+\frac{1}{n_2}}
= 1.9482\sqrt{\frac{1}{10}+\frac{1}{15}}
= 1.9482\sqrt{0.16667} \approx 1.9482 \times 0.40825 \approx 0.7953.
\]
Margin of error:
\[
E = t_{0.025,23} \times \text{SE} \approx 2.069 \times 0.7953 \approx 1.645.
\]
Difference of sample means: \(\bar{x}_1-\bar{x}_2 = 25.2 - 24.6 = 0.6\).
\[
\text{95\% CI} = 0.6 \pm 1.645 \quad\Rightarrow\quad (-1.045,\ 2.245).
\]

\item \textbf{Interpretation (3 marks):}
We are \(95\%\) confident that the true difference \(\mu_1-\mu_2\) lies between \(-1.045\) and \(2.245\). Since the interval includes \(0\), there is no statistically significant difference between the two population means at the \(5\%\) significance level.
\end{enumerate}
\end{corrige}

\begin{corrige}[Exercise 10 — Confidence intervals and sample size determination]
A random sample of \(n=100\) bags of cement gives \(\bar{x}=49.8\,\mathrm{kg}\). The population standard deviation is known to be \(\sigma=1.2\,\mathrm{kg}\).

\begin{enumerate}
\item \textbf{95\% confidence interval for \(\mu\) (4 marks):}
For known \(\sigma\) we use the standard normal distribution. \(z_{0.025} = 1.96\).
Standard error \(= \dfrac{\sigma}{\sqrt{n}} = \dfrac{1.2}{10} = 0.12\,\mathrm{kg}\).
Margin of error \(= 1.96 \times 0.12 = 0.2352\,\mathrm{kg}\).
\[
\text{CI} = 49.8 \pm 0.2352 = (49.5648,\ 50.0352)\,\mathrm{kg}.
\]

\item \textbf{Minimum sample size for a 95\% CI width \(\le 0.4\,\mathrm{kg}\) (5 marks):}
Width \(= 2 \times z_{0.025} \times \dfrac{\sigma}{\sqrt{n}} \le 0.4\).
\[
2 \times 1.96 \times \frac{1.2}{\sqrt{n}} \le 0.4
\;\Rightarrow\; \frac{4.704}{\sqrt{n}} \le 0.4
\;\Rightarrow\; \sqrt{n} \ge \frac{4.704}{0.4} = 11.76
\;\Rightarrow\; n \ge 138.30.
\]
The minimum integer sample size is \(n = 139\).

\item \textbf{If \(\sigma\) is unknown and we use the sample standard deviation \(s=1.2\,\mathrm{kg}\) (5 marks):}
We must use the \(t\)-distribution with \(\nu = n-1 = 99\) degrees of freedom.
\(t_{0.025,99} \approx 1.984\) (slightly larger than \(1.96\)).
The confidence interval becomes
\[
\bar{x} \pm t_{0.025,99} \frac{s}{\sqrt{n}} = 49.8 \pm 1.984 \times 0.12 = 49.8 \pm 0.2381 = (49.5619,\ 50.0381).
\]
The interval is slightly wider because the \(t\)-distribution has heavier tails than the normal distribution, reflecting the extra uncertainty from estimating \(\sigma\) by \(s\).
\end{enumerate}
\end{corrige}

\begin{corrige}[Exercise 11 — GCE-style: Hypothesis test via confidence interval]
A random sample of \(n=64\) packets of biscuits yields \(\sum x = 32064\,\mathrm{g}\) and \(\sum(x-\bar{x})^2 = 1008\,\mathrm{g}^2\).

\begin{enumerate}
\item \textbf{Sample mean and unbiased variance estimate (3 marks):}
\[
\bar{x} = \frac{32064}{64} = 501\,\mathrm{g}.
\]
\[
s^2 = \frac{\sum(x-\bar{x})^2}{n-1} = \frac{1008}{63} = 16\,\mathrm{g}^2,\quad s = 4\,\mathrm{g}.
\]

\item \textbf{90\% confidence interval for \(\mu\) (5 marks):}
Since \(n=64\) is large, the sampling distribution of \(\bar{X}\) is approximately normal by the Central Limit Theorem. We may use the \(t\)-distribution with \(\nu=63\) (or the normal approximation).
\(t_{0.05,63} \approx 1.669\) (using \(z_{0.05}=1.645\) is also acceptable at this sample size).
Standard error \(= \dfrac{s}{\sqrt{n}} = \dfrac{4}{8} = 0.5\,\mathrm{g}\).
Using \(t\): margin \(= 1.669 \times 0.5 = 0.8345\,\mathrm{g}\).
\[
\text{90\% CI} = 501 \pm 0.8345 = (500.1655,\ 501.8345)\,\mathrm{g}.
\]
(If \(z\) is used: margin \(= 1.645 \times 0.5 = 0.8225\), giving \((500.1775,\ 501.8225)\).)

\item \textbf{Test the claim \(\mu = 502\,\mathrm{g}\) at the \(10\%\) significance level (4 marks):}
The \(90\%\) confidence interval is \((500.17,\ 501.83)\) (using either \(t\) or \(z\)). The claimed value \(502\,\mathrm{g}\) lies \emph{outside} this interval. Therefore we \textbf{reject} the manager's claim at the \(10\%\) significance level. There is sufficient evidence that the true mean content is not \(502\,\mathrm{g}\).

\item \textbf{Assumption made (2 marks):}
The sample is a random sample from the population. Because \(n=64\) is large, the Central Limit Theorem ensures the sampling distribution of \(\bar{X}\) is approximately normal even if the population distribution is not normal. (If \(n\) were small, we would need to assume the population is normally distributed.)

\item \textbf{Effect of reducing the sample size to \(16\) (2 marks):}
With \(n=16\), standard error \(= \dfrac{s}{\sqrt{16}} = \dfrac{4}{4} = 1\,\mathrm{g}\).
The critical value would be \(t_{0.05,15} \approx 1.753\).
Margin of error \(= 1.753 \times 1 = 1.753\,\mathrm{g}\), which is more than double the previous margin (\(0.8345\)).
The width of the confidence interval is inversely proportional to \(\sqrt{n}\); halving \(\sqrt{n}\) (from \(8\) to \(4\)) roughly doubles the width.
\end{enumerate}
\end{corrige}

\newpage

\coursec{Summary sheet}

\begin{popfigure}
\begin{tikzpicture}[grow cyclic, mindmap, concept color=popblue,
level 1/.append style={level distance=3.6cm, sibling angle=51.4},
level 2/.append style={level distance=2.2cm},
every node/.style={concept, scale=0.75},
root concept/.append style={font=\small\bfseries}]
\node[root concept]{Sampling \& Estimation}
child[concept color=poporange]{ node{Finite population}
  child{ node{With / without replacement} }
  child{ node{$\binom{N}{n}$ or $N^n$ samples} }
}
child[concept color=popgreen]{ node{Sampling distributions}
  child{ node{$\bar X$ statistic} }
  child{ node{Random numbers} }
}
child[concept color=poppurple]{ node{Central limit theorem}
  child{ node{$\bar X\sim N\!\left(\mu,\frac{\sigma^2}{n}\right)$} }
  child{ node{Normal if $n\ge 30$} }
}
child[concept color=popgold]{ node{Sampling methods}
  child{ node{Random, stratified} }
  child{ node{Systematic, quota} }
}
child[concept color=poppink]{ node{Point estimation}
  child{ node{Unbiased: $\hat\mu=\bar x$, $\hat\sigma^2=\frac{n}{n-1}s^2$} }
  child{ node{Most efficient} }
  child{ node{Pooled estimates} }
}
child[concept color=popgreen!70!black]{ node{Confidence intervals}
  child{ node{$\sigma^2$ known: $z$} }
  child{ node{$\sigma^2$ unknown: $t$} }
  child{ node{Difference of means} }
};
\end{tikzpicture}
\legende{Mind map of the chapter: from sampling to confidence intervals.}
\end{popfigure}

\begin{bilan}
\textbf{Key definitions}
\begin{itemize}
\item A \cle{population} is the whole set of individuals or items under study; a \cle{sample} is a subset actually observed; $N$ = population size, $n$ = sample size.
\item \cle{Sampling with replacement}: each item is returned before the next draw; draws are independent; number of ordered samples: $N^n$. \cle{Without replacement}: number of unordered samples: $\binom{N}{n}$.
\item A \cle{statistic} (e.g.\ $\bar X$, $S^2$) is a random variable computed from the sample; its probability distribution is its \cle{sampling distribution}.
\item An \cle{estimator} $\hat\theta$ is \cle{unbiased} if $E(\hat\theta)=\theta$; among unbiased estimators, the \cle{most efficient} is the one with smallest variance.
\end{itemize}

\textbf{Formulas and theorems to know}
\begin{itemize}
\item Sample mean: $\bar X=\dfrac{1}{n}\sum X_i$; \quad $E(\bar X)=\mu$, \quad $\mathrm{Var}(\bar X)=\dfrac{\sigma^2}{n}$ (infinite population or sampling with replacement).
\item Finite population, without replacement: $\mathrm{Var}(\bar X)=\dfrac{\sigma^2}{n}\cdot\dfrac{N-n}{N-1}$ (finite population correction).
\item Unbiased estimate of variance: $\hat\sigma^2=\dfrac{1}{n-1}\sum (x_i-\bar x)^2=\dfrac{n}{n-1}s^2$.
\item \cle{Central limit theorem}: for large $n$ ($n\ge 30$), $\bar X\approx N\!\left(\mu,\dfrac{\sigma^2}{n}\right)$ whatever the population distribution; if the population is already $N(\mu,\sigma^2)$, the result is exact for all $n$.
\item Pooled estimates from two samples $(n_1,\bar x_1,s_1^2)$, $(n_2,\bar x_2,s_2^2)$:
\[ \hat\mu=\frac{n_1\bar x_1+n_2\bar x_2}{n_1+n_2},\qquad \hat\sigma^2=\frac{n_1s_1^2+n_2s_2^2}{n_1+n_2-2}. \]
\item Confidence interval for $\mu$, $\sigma^2$ known (or $n\ge 30$):
\[ \bar x\pm z_{\alpha/2}\,\frac{\sigma}{\sqrt{n}}. \]
Common values: $90\%\!:1.645$; $95\%\!:1.96$; $99\%\!:2.576$.
\item Confidence interval for $\mu$, $\sigma^2$ unknown, small sample from a normal population:
\[ \bar x\pm t_{\alpha/2}(n-1)\,\frac{\hat\sigma}{\sqrt{n}}. \]
\item Difference of two means (independent samples):
\[ (\bar x_1-\bar x_2)\pm z_{\alpha/2}\sqrt{\frac{\sigma_1^2}{n_1}+\frac{\sigma_2^2}{n_2}} \]
(use $t$ with pooled $\hat\sigma^2$ and $n_1+n_2-2$ degrees of freedom when variances are unknown but assumed equal).
\end{itemize}

\textbf{Methods}
\begin{itemize}
\item To build a sampling distribution: list all possible samples, compute the statistic for each, then tabulate values with their probabilities.
\item To use random number tables: number the population, read digits in groups, ignore repeats and out-of-range numbers.
\item For a confidence interval: identify whether $\sigma$ is known, choose $z$ or $t$, compute the standard error, apply the formula, and interpret: ``we are $95\%$ confident that $\mu$ lies in ...''.
\end{itemize}

\textbf{Common mistakes}
\begin{itemize}
\item Confusing $s^2$ (divisor $n$) with the unbiased estimate $\hat\sigma^2$ (divisor $n-1$).
\item Forgetting $\sqrt{n}$ in the standard error, or using $\sigma^2/n$ instead of $\sigma/\sqrt{n}$.
\item Using $z$ when $\sigma$ is unknown and $n<30$: use the $t$-distribution.
\item Forgetting the finite population correction $\frac{N-n}{N-1}$ when sampling without replacement from a small population.
\item Mixing up the confidence level with the probability that $\mu$ is fixed: $\mu$ is constant, the interval is random.
\end{itemize}

\textbf{GCE tips}
\begin{itemize}
\item Quote the distribution before computing: $\bar X\sim N(\mu,\sigma^2/n)$.
\item Show the standardisation $Z=\dfrac{\bar X-\mu}{\sigma/\sqrt{n}}$ explicitly to earn method marks.
\item In pooled estimation, check the denominator $n_1+n_2-2$ and give answers to at least 2 decimal places.
\item Always end an interval question with a clear sentence of interpretation in context.
\end{itemize}
\end{bilan}

\findecours

\end{document}
